\documentclass[11pt]{article}
\usepackage[T1]{fontenc}
\usepackage[utf8]{inputenc}
\usepackage{lmodern}
\usepackage{microtype}
\usepackage[a4paper,margin=30mm]{geometry}
\usepackage{amsmath,amssymb,amsthm,mathtools}
\usepackage{booktabs,array,tabularx}
\usepackage{enumitem}
\usepackage{xcolor}
\usepackage{hyperref}
\usepackage[nameinlink,noabbrev]{cleveref}

\hypersetup{
  colorlinks=true,
  linkcolor=blue!55!black,
  citecolor=blue!55!black,
  urlcolor=blue!55!black,
  pdftitle={A Finite Constructive Fundamental Theorem of Algebra and Rational-Function Argument Principle over the Gaussian Rationals},
  pdfauthor={Anonymous},
  pdfsubject={Winding certificates and rational-box root-approximation tracks}
}
\setlist{nosep,leftmargin=*}
\setlength{\parskip}{0.35em}
\setlength{\parindent}{1.25em}
\setlength{\emergencystretch}{3em}

\newtheorem{theorem}{Theorem}[section]
\newtheorem{lemma}[theorem]{Lemma}
\newtheorem{proposition}[theorem]{Proposition}
\newtheorem{corollary}[theorem]{Corollary}
\newtheorem{definition}[theorem]{Definition}
\newtheorem{remark}[theorem]{Remark}
\newtheorem{example}[theorem]{Example}
\newtheorem{construction}[theorem]{Construction}

\newcommand{\Qi}{\mathbb{Q}(i)}
\newcommand{\Zi}{\mathbb{Z}[i]}
\newcommand{\Q}{\mathbb{Q}}
\newcommand{\N}{\mathbb{N}}
\newcommand{\Z}{\mathbb{Z}}
\newcommand{\Norm}[1]{\operatorname{N}\!\left(#1\right)}
\newcommand{\Ree}{\operatorname{Re}}
\newcommand{\Imm}{\operatorname{Im}}
\newcommand{\Wind}{\operatorname{Wind}}
\newcommand{\Ind}{\operatorname{Ind}}
\newcommand{\Var}{\operatorname{Var}}
\newcommand{\Disc}{\operatorname{Disc}}
\newcommand{\diamq}{\operatorname{diag}^{2}}
\newcommand{\Clear}{\operatorname{Clear}}
\newcommand{\Count}{\operatorname{Count}}
\newcommand{\eps}{\varepsilon}
\newcommand{\half}{\tfrac12}
\newcommand{\Defref}[1]{Definition~\ref{#1}}
\newcommand{\Thmref}[1]{Theorem~\ref{#1}}
\newcommand{\Corref}[1]{Corollary~\ref{#1}}
\newcommand{\leanref}[1]{\par\smallskip\noindent
  \emph{Formal verification:}\par\nobreak
  {\raggedright\nolinkurl{#1}.\par}}

\title{A Finite Constructive Fundamental Theorem of Algebra and Rational-Function Argument Principle over the Gaussian Rationals\\
\large Signed winding certificates and rational-box approximation tracks}
\author{Rich Lawrence}
\date{August 8, 2026}

\begin{document}
\maketitle

\begin{abstract}
We prove a finite, computational analogue of the Fundamental Theorem of
Algebra for polynomials over the Gaussian rationals.  For a squarefree
polynomial $f\in\Qi[Z]$ of degree $d\ge2$, a finite algorithm returns exactly
$d$ pairwise separated refinement machines.  Given a positive rational
accuracy request, each machine returns a strictly refining rational box with
clear boundary, algebraic winding number one, and squared diagonal below the
requested bound.  Canonical Gaussian-rational anchors in these boxes satisfy
explicit coefficient-dependent residual estimates, while the boxes also
control Hermitian projective quadrance.  In degree one the zero is computed
exactly.
The construction does not assume that an ``irrational root'' is an
object at all, completed or otherwise: the finite interface answering each
finite request is the entire construction.  The construction computes winding
through signed Euclidean remainder chains and uses fraction-free Sylvester
eliminants, rational Taylor bounds, bounded searches, and recursion on natural
numbers.

We also derive a finite rational-function argument principle.  After
executable cancellation and squarefree-multiplicity decomposition, the
algorithm returns a jointly separated family of zero and pole sites, exact
for linear factors and given by shrinking box tracks otherwise.  A site of
multiplicity $m$ has winding $+m$ for a numerator factor and $-m$ for a
denominator factor; the weighted totals are the reduced degrees, and a
coefficient-computed outer square has winding
$\deg F_{\rm red}-\deg G_{\rm red}$.
\end{abstract}

\tableofcontents

\section{Introduction}

The Fundamental Theorem of Algebra is ordinarily regarded as settled: every
nonconstant polynomial over the complex numbers is said to have a complex
zero.  That statement, however, is formulated inside a completed complex
field, and every proof route known to us uses assumptions that this paper
deliberately declines.  Analytic proofs begin with completed real and complex
fields and use continuity there, whether through limits, the intermediate
value theorem, compactness or minimization, integration, or related analytic
machinery.  Algebraic and topological proofs may avoid some of those devices,
but still appeal to infinitary structures or existence principles such as
algebraic closures, infinite extensions, global collections of roots, or
choice.  We do not dispute those proofs relative to their intended
foundations.  We ask instead what can be proved when no completed real or
complex field, actual infinite object or process, analytic principle, or
choice-based existence argument is admitted.

After fixing that discipline, this introductory section introduces, in
dependency order, the finite definitions needed to state our constructive
FTA and rational-function argument principle.  The remainder of the paper
then serves as a full proof of those two theorems: each later section
constructs and verifies part of the finite algorithms and the certificates
they return.

\subsection{Statement, terminology, and constructive discipline}

Under our proposed discipline the classical conclusion cannot simply be restated:
to assert an arbitrary complex zero would already assume the sort of object
we have declined to introduce.  Our aim is the closest finite,
calculation-based counterpart.  From finite Gaussian-rational coefficient
data we compute coefficient-checked winding certificates and finitely many
independent programs.  Each program answers any displayed finite accuracy
request with a smaller rational box and finite proofs of refinement,
clearance, winding, separation, and quantitative error bounds.  No
irrational root is assumed to exist behind these answers, as a completed
point or as an actual infinite sequence of approximations.  The finite
request-and-answer interface is itself the result.

\subsubsection{Related work}

Choice-free constructive versions of the FTA already reach much further than
a merely numerical approximation theorem.  Ruitenburg \cite{Ruitenburg}
constructs roots over modulated Cauchy complex numbers without a choice
principle.  Richman \cite{Richman} develops a completion suitable for a
choiceless setting and identifies a monic polynomial with a completed limit
multiset: a coherent method of approximately factoring the polynomial, with
no advance selection of one member as a distinguished root.  This is
conceptually close to the present output.  Indeed, each refinement procedure
eventually returned here could be used as a computable name for a completed
value: it answers every displayed positive rational accuracy request with
finer rational data, from which one could subsequently form a Cauchy sequence
and then a completed complex value.  We stop before doing so: no completion is
formed, equivalent names are not quotiented, and convergence to or equality
of completed points is not used.  Instead the finite program and each checked
response are the theorem's objects.  For squarefree displayed input the
output also keeps the degree-many refinement procedures separated and
independently refinable, rather than completing one unordered multiset.
Conversely, those
constructive developments admit coefficients in their completed complex
field, whereas our coefficient input is the narrower, fully displayed field
$\Qi$.  The construction in Richman's paper is presented in prose, whereas
here the finite interfaces, the programs producing their certificates, and
their stated guarantees are formalized together in Lean.

There is also a substantial formal lineage.  Geuvers, Wiedijk, and
Zwanenburg \cite{GeuversWiedijkZwanenburg} formalize Kneser's constructive
FTA in Coq.  Their real-number structure is axiomatically a Cauchy-complete
Archimedean ordered field, and their complex numbers are pairs of such reals;
the formal theorem therefore retains the completed carrier and its root
conclusion.  Li and Paulson \cite{LiPaulson} formally verify in Isabelle/HOL
that winding numbers can be evaluated through Cauchy indices, and combine
that calculation with the classical argument principle to count complex
roots in rectangles and half-planes.  Their work gives an effective,
machine-checked root-counting procedure, but its correctness theorem still
uses complex-analytic winding and existing complex roots.

The closest algebraic route begins with Eisermann \cite{Eisermann}, who
constructs a real-algebraic winding number from Sturm chains, derives the FTA
over real-closed fields, and explains the resulting root-finding algorithm.
This makes the classical theorem effective, but still assumes an ambient
real-closed field and the polynomial intermediate value theorem.  Perrucci
and Roy \cite{PerrucciRoyQuantitative} quantify that dependence, showing that
degree-$d$ FTA by this route needs the intermediate value theorem only through
degree $d^2$.  Perrucci and Roy \cite{PerrucciRoyWinding} later correct the
unrestricted auxiliary product formula, give a fully additive algebraic
winding, and extend the count to rational functions, including zeros and
poles on rectangle boundaries.  Their result remains a
theorem over $\mathbf C=\mathbf R[i]$, with $\mathbf R$ real closed, and counts
zeros and poles already present in that algebraic closure.  Our finite
Cauchy-index calculations begin from the same algebraic insight but construct
signed sites without adopting that carrier or its collection of points.

Finally, recursive rectangular isolation is itself classical algorithmic
territory.  Wilf \cite{Wilf} gives a global bisection algorithm that repeatedly
uses argument-principle root counts to refine complex isolating rectangles.
The subdivision shape is therefore not new here.  What changes is its proof
boundary and output: safe rational cuts, winding and clearance certificates,
total independently callable refiners, and their machine-checked assembly
without an analytic root-counting oracle.  Across these related lines, the
distinctive claim of the present work is the combined package, not the
novelty of constructive approximation, algebraic winding, or box subdivision
taken separately.

\subsubsection{Finite geometric language}

This restriction governs the paper's geometric language as well.  Membership
and path descriptions abbreviate calculations on displayed rational data.
We use the following conventions.

\begin{itemize}
\item A rational interval or box is a finite record of endpoints.  To say
that a displayed point is inside, outside, or on its boundary means that the
corresponding decidable rational equalities and inequalities hold.

\item A statement such as \emph{every $z\in B$} means that, given any
displayed $z$, the box-membership inequalities imply the stated conclusion.
It does not construct or range over a completed collection of all points of
$B$.

\item An oriented rational path is a finite list of affine
parametrizations.  Its boundary image under a polynomial is the corresponding
finite list of composed coordinate polynomials; it is not a completed curve
or a set of all its point-values.

\item A region is a rational box or cell record together with its predicates,
and a family is always an explicitly finite list or finite-indexed record.
Thus phrases such as \emph{on the path}, \emph{throughout the region}, and
\emph{the image of the boundary} assert conditional identities or
inequalities for displayed rational parameters.  They introduce no carrier
containing all possible parameter values or points.
\end{itemize}

\subsubsection{Polynomial equality}
\label{sec:polynomial-equality}

A displayed polynomial is a normalized finite coefficient list, written from
highest degree to constant term.  Leading zero entries are removed.  We
define two displayed polynomials to be equal exactly when their normalized
coefficient lists agree entry by entry.  For polynomials whose coefficients
are themselves polynomials, and for the bivariate coordinate polynomials
used below, the same definition is applied recursively to the finite
coefficient data.

Consequently every polynomial identity in this paper---including B\'ezout,
factorization, derivative, substitution, and polynomial-matrix
identities---means equality of normalized coefficient lists.  Evaluation by
finite Horner calculation preserves such an identity, but agreement of
values is never used as the definition of polynomial equality.  By contrast,
an expression such as $p(q)=0$ is explicitly an equality of two displayed
Gaussian-rational values after evaluation.

\subsection{Machine verification}

The finite machinery developed in this paper has also been implemented and
verified in Lean.  The accompanying formalization contains the executable
procedures together with machine-checked proofs of the paper-facing
guarantees; the \emph{Formal verification} notes throughout identify the
corresponding Lean declarations.  This is not only a transcription of the
finished prose proof, but a second presentation in which the constructions
themselves can be run and their full dependency boundary inspected.

This is precisely the kind of mathematics that ought to benefit from machine
verification.  Every object in the construction is finite data; every search
has a displayed finite bound; every recursive call decreases an explicit
natural measure; and every certificate is checked by finite algebra, integer
arithmetic, or rational comparison.  There is no appeal to a picture, a
limiting object, or an unnamed witness whose properties must be taken on
trust.  A proof assistant can therefore check the argument at the same level
at which the theorem is stated.  Lean also permits the programs and their
correctness proofs to share one environment, so the formalization checks not
only that the certificates imply the advertised conclusions, but that the
finite procedures producing them are total and executable.  Machine
verification is consequently part of the constructive discipline of this
paper, not decorative evidence added after the mathematics was finished.

The consequences of that choice for Lean, Mathlib, and the proof's audited
axiom boundary are discussed in \cref{sec:formalization-audit}.

\subsection{Definitions and theorem statements}

The motivational and foundational preliminaries are now complete.  We begin
the chain of finite definitions needed to state the two main theorems.  The
definitions are arranged in dependency order and culminate first in the
finite constructive FTA, and then in the finite rational-function argument
principle.  This subsection deliberately prioritizes reaching those
statements quickly: it retains the data needed to read them, while deferring
most intuition, motivation, and derivation to the later section where each
construction is first used.  The short links following each group point to
that deferred development.

Polynomial names are role-based.  The letter $f$ denotes the polynomial
currently under the main winding, isolation, or FTA argument, whereas $p$
denotes an arbitrary polynomial supplied to a reusable bound, auxiliary
algorithm, or preprocessing step.  Letters such as $g,h$ denote temporary
comparison polynomials, GCDs, or factors, and $U,V$ denote displayed Bézout
cofactors.  In the rational-function statements, $F,G$ are the original
numerator and denominator, $F_0,G_0$ their reduced cofactors, $P,Q$ a generic
polynomial pair, and indexed $p_j,q_k$ individual squarefree factors.

\subsubsection{Polynomial input certificates}

\begin{definition}[squarefree certificate]
\label{def:squarefree-certificate}
For $f\in\Qi[Z]$, a \emph{squarefree certificate} is a checked nonzero
constant output of the finite Euclidean GCD computation on $f,f'$.  Rescaling
the computed Bézout cofactors gives displayed $U,V\in\Qi[Z]$ satisfying
\[
                       Uf+Vf'=1.
\]
\end{definition}

\leanref{Purelib.Polynomial.Squarefree.SquarefreeCertificate;
Purelib.Polynomial.Squarefree.CoprimeCertificate.identity}

Its quantitative role is developed in \cref{sec:finite-bounds}.

\subsubsection{Finite metric and box data}

\begin{definition}[Gaussian-rational quadrance]
\label{def:quadrance}
For a displayed Gaussian rational $z=x+iy$, its \emph{quadrance} is
\[
             \Norm z=x^2+y^2\in\Q_{\ge0}.
\]
\end{definition}

\leanref{Purelib.Data.Gaussian.normSq}

\begin{definition}[Hermitian projective quadrance]
\label{def:hermitian-projective-quadrance}
For displayed $z,w\in\Qi$, their \emph{Hermitian projective quadrance} is
\begin{equation}
\label{eq:hermitian-projective-affine}
 \chi^2(z,w)=
 \frac{\Norm{z-w}}{(1+\Norm z)(1+\Norm w)}.
\end{equation}
This is the Gaussian-Hermitian affine analogue of Wildberger's projective
quadrance \cite{Wildberger}; here it is used only as a rational error measure.
\end{definition}

\leanref{Purelib.Polynomial.Geometry.Chordal.affineChordalSq}

\begin{definition}[rational box]
\label{def:rational-box}
A \emph{rational box} is a tuple
\[
 B=(x_-,x_+,y_-,y_+)\in\Q^4,
 \qquad x_-<x_+,\quad y_-<y_+.
\]
A point $z=x+iy$ belongs to $B$, written $z\in B$, when
\[
        x_-\le x\le x_+,\qquad y_-\le y\le y_+.
\]
The \emph{squared diagonal} of $B$ is
\[
       \diamq(B)=(x_+-x_-)^2+(y_+-y_-)^2.
\]
One box is contained in another, written $D\subseteq B$, when both of its
coordinate intervals are contained; this symbol allows equality.  If
$D=(u_-,u_+,v_-,v_+)$, write $D\subset B$ and call $D$ a \emph{strict box
refinement} of $B$ when
\[
 x_-<u_-,\qquad u_+<x_+,\qquad
 y_-<v_-,\qquad v_+<y_+.
\]
Thus $\subset$ always denotes strict coordinatewise containment between
rational boxes in this paper, whereas $\subseteq$ denotes containment that
may include equality.
\end{definition}

\leanref{Purelib.Polynomial.Geometry.Boxes.Box; Purelib.Polynomial.Geometry.Boxes.Box.Proper;
Purelib.Polynomial.Geometry.Boxes.Box.Contains; Purelib.Polynomial.Geometry.Boxes.Box.Refines;
Purelib.Polynomial.Geometry.Boxes.Box.StrictInterior; Purelib.Polynomial.Geometry.Boxes.Box.diagonalSq}

\begin{definition}[box anchor and opposite corner]
\label{def:box-corners}
For the rational box $B$ of \cref{def:rational-box}, put
\[
 q_B=x_-+iy_-,
 \qquad
 v_B=x_++iy_+.
\]
Call $q_B$ the canonical rational \emph{anchor} of $B$.
\end{definition}

\leanref{Purelib.Polynomial.Geometry.Boxes.Box.anchor; Purelib.Polynomial.Geometry.Chordal.oppositeCorner}

\begin{definition}[coefficient isolation square]
\label{def:coefficient-isolation-square}
Let
\[
       f(Z)=a_dZ^d+\sum_{k=0}^{d-1}a_kZ^k,
       \qquad d\ge1,\quad a_d\ne0.
\]
Compute the two positive rational quantities
\begin{equation}
\label{eq:isolation-radius}
 C_f=d\sum_{k=0}^{d-1}2^k\Norm{a_k},
 \qquad
 R_f=\max\left\{1,\frac{C_f}{\Norm{a_d}}+1\right\}.
\end{equation}
The \emph{coefficient isolation square} of $f$ is the rational box
\[
                     K_f=(-R_f,R_f,-R_f,R_f).
\]
\end{definition}

\leanref{Purelib.RootIsolation.Winding.Global.denseLargeSquareRadius;
Purelib.RootIsolation.Winding.Global.largeSquareBox}

\begin{definition}[ambient box]
\label{def:ambient-box}
An \emph{ambient box} is a fixed rational box $A$
together with the stored containment certificates $B\subseteq A$ for the
working boxes, supplying one common domain for coefficient bounds.  In the
main polynomial theorem $A=K_f$.
\end{definition}

\begin{definition}[common coefficient isolation square]
\label{def:common-coefficient-isolation-square}
For two nonzero polynomials $P,Q$, let $R_{P,Q}$ be the larger of the radii
given by \eqref{eq:isolation-radius} for their nonconstant members.  A
constant imposes no radius condition; if both polynomials are constant, put
$R_{P,Q}=1$.  Their
\emph{common coefficient isolation square} is
\[
 K_{P,Q}=(-R_{P,Q},R_{P,Q},-R_{P,Q},R_{P,Q}).
\]
\end{definition}

For box records, $A$ denotes a fixed ambient box, $B$ a generic current or
parent box, $C$ an intermediate child selected during subdivision, and $D$
a returned refined box.  The symbols $K_f,K_R,K_{P,Q}$ are reserved for
coefficient-computed outer isolation squares rather than arbitrary boxes.
Displayed Gaussian-rational samples are denoted by $z,w,q$; $q_B$ and $v_B$
are respectively the anchor and opposite corner of $B$.

\begin{definition}[quadrance-Lipschitz certificate]
\label{def:quadrance-lipschitz}
Let $p\in\Qi[Z]$ and let $A$ be a rational box.  A
\emph{quadrance-Lipschitz certificate} on $A$ is a displayed rational
$L>0$ together with a finite proof that, for any displayed $z,w\in\Qi$,
the two box-membership predicates imply
\begin{equation}
\label{eq:quadrance-lipschitz}
             \Norm{p(z)-p(w)}\le L\Norm{z-w}.
\end{equation}
We write $L_1(p,A)$ for the particular coefficient-computed certificate
constructed in \cref{sec:finite-bounds} from the finite divided-difference
identity \eqref{eq:divdiff}.
\end{definition}

\leanref{Purelib.Polynomial.Bounds.denseLipschitzBound;
Purelib.Polynomial.Bounds.normSq_eval_sub_eval_le_denseLipschitzBound}

The metric estimates are proved in \cref{sec:finite-bounds}, and the degree
winding of the coefficient isolation square in \cref{sec:global-winding}.

\subsubsection{Certified states and finite tracks}

\begin{definition}[boundary-clearance certificate]
\label{def:boundary-clearance}
For $f\in\Qi[Z]$ and a rational box $B$, a
\emph{boundary-clearance certificate} is a displayed rational $m_B>0$
together with four finite proofs that
\[
                 m_B\le\Norm{f(z)}
\]
on the four rational edge parametrizations of $B$.
\end{definition}

\leanref{Purelib.Polynomial.Geometry.BoxClearance.BoundaryClearance}

\begin{definition}[box winding]
\label{def:rational-function-winding}
Let $P,Q\in\Qi[Z]$, and let a rational box $B$ carry boundary-clearance
certificates for both polynomials.  On each oriented affine edge, represent
the direction of the displayed quotient $P(z)/Q(z)$ by the
Gaussian-rational coordinate pair
\[
                         P(z)\overline{Q(z)}.
\]
The four finite signed Cauchy--Sturm edge calculations,
realized explicitly in \cref{sec:winding-realization}, give the integer
$\Wind_4(P/Q;B)$.  Its conventional normalization is
\[
                         \Wind_{P/Q}(B)=\Wind_4(P/Q;B)/4.
\]
For a polynomial $f$, the case $P=f,Q=1$ is denoted
$\Wind_4(f;B)$ and $\Wind_f(B)$.  The finite constructions below use
$\Wind_4$; the normalized notation is reserved for the headline statements.
No winding is assigned without the required boundary-clearance certificates.
\end{definition}

\leanref{Purelib.RationalFunctions.Winding.rationalQuadrupledWinding;
Purelib.RationalFunctions.Winding.rationalBoxWinding;
Purelib.RootIsolation.Winding.quadrupledBoxWinding; Purelib.RootIsolation.Winding.boxWinding}

\begin{definition}[winding-one certificate]
\label{def:winding-one-certificate}
A \emph{winding-one certificate} for $f$ and $B$ consists of its four finite
edge-chain calculations, their endpoint signs, and the integer check
\[
                           \Wind_4(f;B)=4.
\]
A finite root state also carries the clearance certificate of
\cref{def:boundary-clearance}.
\end{definition}

\leanref{Purelib.RootIsolation.Isolation.ShrinkingTrack.WindingOneState;
Purelib.RootIsolation.Winding.quadrupledBoxWinding}

\begin{definition}[finite root state]
\label{def:finite-root-state}
Let $f\in\Qi[Z]$.  A \emph{finite root state} consists of
\[
       (B,\mathcal C,\mathcal W),
\]
where $B$ is a rational box, $\mathcal C$ is a boundary-clearance
certificate, and $\mathcal W$ is a winding-one certificate.
Its displayed sample is the anchor $q_B$ of
\cref{def:box-corners}; the defining rational inequalities give
$q_B\in B$.
\end{definition}

\leanref{Purelib.RootIsolation.Isolation.ShrinkingTrack.WindingOneState}

\begin{definition}[refinement of root states]
\label{def:root-state-refinement}
A finite root state $\sigma'$ \emph{refines} a state $\sigma$ when its box
$B'$ satisfies $B'\subseteq B$ in the sense of
\Defref{def:rational-box}.  The state refinement is \emph{strict} when
$B'\subset B$.  In either case its refinement certificate is the finite list
of coordinate inequalities witnessing the indicated box relation.
\end{definition}

From this point $\sigma,\sigma'$ denote complete certified state records;
$B,B'$ denote their box components when no other state data is needed.

\begin{definition}[root-approximation track, RAT]
\label{def:rat}
A \emph{root-approximation track}, abbreviated \emph{RAT}, for $f$ consists
of a finite initial root state $\sigma_0$ and a deterministic refinement rule.
Given the current root state $\sigma$ and a displayed $\eps\in\Q_{>0}$, the rule
terminates and returns one new root state
\[
                    (\sigma,\eps)\longmapsto
                    \sigma'=\operatorname{refine}(\sigma,\eps).
\]
The returned state $\sigma'$ is a strict refinement of $\sigma$, and its box $B'$
satisfies
\[
                  \diamq(B')\le\eps^2.
\]
\end{definition}

\leanref{Purelib.RootIsolation.Isolation.ShrinkingTrack.WindingOneTrack;
Purelib.RootIsolation.Isolation.ShrinkingTrack.ofInitial}

\begin{definition}[finite observations and separation]
\label{def:finite-observation-separation}
An \emph{exact terminal} for a polynomial $p$ is a displayed $q\in\Qi$
with exact evaluation $p(q)=0$.  A \emph{finite observation}
is either an exact terminal or a RAT.  At any finite stage it has a rational
\emph{support record}
\[
                   S=(x_-^S,x_+^S,y_-^S,y_+^S).
\]
For an exact terminal $q=a+ib$, put $S=(a,a,b,b)$.  For a RAT, make a
finite call and use its returned box as $S$.

Two finite observations are \emph{separated} when we can produce support
records $S,T$ and a displayed $g\in\Q_{>0}$ satisfying one of
\[
\begin{array}{ll}
 x_+^S+g\le x_-^T, & x_+^T+g\le x_-^S,\\[2pt]
 y_+^S+g\le y_-^T, & y_+^T+g\le y_-^S.
\end{array}
\]
A finite family is \emph{jointly separated} when such a certificate is
stored for every pair of distinct observations.  Refining a RAT preserves
every stored witness.
\end{definition}

\leanref{Purelib.Polynomial.Geometry.Boxes.Box.Separated;
Purelib.Polynomial.Geometry.Boxes.Box.SeparationCertificate;
Purelib.Polynomial.Geometry.Boxes.Box.separationCertificate;
Purelib.Polynomial.Geometry.Boxes.Box.Separated.of_refines;
Purelib.RationalFunctions.Sites.PositiveSideInput.ObservationsSeparated;
Purelib.RationalFunctions.Sites.PositiveSideInput.separateObservations;
Purelib.RationalFunctions.SignedMultiplicity.TrackPairSeparation}

Certified states, separation, and total tracks are constructed in
\cref{sec:local-extraction,sec:track-separation,sec:exact-isolation}.

\subsubsection{The finite theorem}

\begin{theorem}[finite constructive FTA]
\label{thm:main}
Let $f\in\Qi[Z]$ have degree $d\ge1$, nonzero leading coefficient, and a
squarefree certificate.  A finite algorithm returns one of the following.
\begin{enumerate}[label=(\roman*)]
\item If $d=1$, it returns a directly computed exact terminal $q\in\Qi$ for
$f$.
\item If $d\ge2$, it returns a list of exactly $d$ RATs.
\end{enumerate}
In the second case the returned family has four certified properties.
\begin{enumerate}[label=(\alph*)]
\item \emph{Separation.}  Its initial boxes are pairwise separated in the
sense of \Defref{def:finite-observation-separation}, viewing each RAT as a
finite observation, with displayed positive rational gaps.
\item \emph{Refinement.}  Every call terminates and returns the strict
refinement required by \Defref{def:rat}, including the anchor $q_B$ of
\Defref{def:box-corners}.
\item \emph{Quantitative control.}  Let $K_f$ be the coefficient isolation
square of \Defref{def:coefficient-isolation-square}.  Every returned box
satisfies $B\subseteq K_f$.  The quadrance-Lipschitz certificate of
\Defref{def:quadrance-lipschitz} gives
\[
             \Norm{f(q_B)}\le L_1(f,K_f)\diamq(B).
\]
For every displayed $z\in B$, one also has
\[
 \begin{aligned}
 \chi^2(q_B,z)&\le\diamq(B),\\
 \Norm{f(z)}&\le4L_1(f,K_f)\diamq(B),\\
 \Norm{f(q_B)}&\le
 L_1(f,K_f)(1+\Norm{q_B})(1+\Norm{v_B})\chi^2(q_B,v_B).
 \end{aligned}
\]
In particular a request with $\diamq(B)\le\eps^2$ bounds both quantities by
$\chi^2(q_B,z)\le\eps^2$ and
$\Norm{f(q_B)}\le L_1(f,K_f)\eps^2$.  Consecutive anchors $q_B,q_{B'}$ satisfy
$\Norm{q_{B'}-q_B}\le\diamq(B)$ for the earlier box.
\item \emph{Finite maximality.}  No finite family of pairwise separated RATs
supported in the same coefficient isolation square $K_f$ has length greater
than $d$.
\end{enumerate}
The algorithm uses only finite arithmetic, bounded searches over displayed
finite lists, and recursion on natural numbers.
\end{theorem}

\leanref{Purelib.RootIsolation.Isolation.BoxFTA.provePositive; Purelib.Papers.RFAP.Paper.isolateByBoxes}

The proof is completed in \cref{sec:exact-isolation}.

\subsubsection{Rational-function data and statement}

The rational-function theorem uses the same polynomial tracks.  The extra
records below say exactly what cancellation, multiplicity, sign, and
cross-side separation mean.

\begin{definition}[reduced quotient presentation]
\label{def:reduced-quotient}
Let $F,G\in\Qi[Z]$ be displayed nonzero polynomials.  A
\emph{reduced quotient presentation} consists of five displayed polynomials
\[
                        H,F_0,G_0,U,V\in\Qi[Z]
\]
and the identities
\[
 \begin{gathered}
 F=HF_0,\qquad G=HG_0,\\
 UF_0+VG_0=1,
 \end{gathered}
\]
together with the finite-division degree checks
\[
 \deg H+\deg F_0=\deg F,
 \qquad
 \deg H+\deg G_0=\deg G.
\]
\end{definition}

\leanref{Purelib.RationalFunctions.Reduction.ReducedPair;
Purelib.RationalFunctions.Reduction.reduce}

\begin{definition}[finite squarefree-multiplicity presentation]
\label{def:squarefree-multiplicity-presentation}
For a displayed nonzero polynomial $P$, a
\emph{finite squarefree-multiplicity presentation} is an identity
\[
                 P=u\prod_{j=1}^{r}p_j^{m_j},
\]
where $u\in\Qi^\times$, every $p_j$ has positive degree and a squarefree
certificate, every $m_j$ is a displayed positive integer, and every pair of
distinct factors $p_j,p_k$ carries displayed $A_{jk},B_{jk}\in\Qi[Z]$ with
the Bézout identity
\[
                 A_{jk}p_j+B_{jk}p_k=1.
\]
The record also contains the checked degree identity
\[
                  \sum_{j=1}^{r}m_j\deg p_j=\deg P.
\]
A nonzero constant has $r=0$ and the empty product.
\end{definition}

\leanref{Purelib.Polynomial.Factorization.Decomposition.positiveMultiplicityEntries;
Purelib.Polynomial.Factorization.Decomposition.positiveMultiplicityDegree;
Purelib.Polynomial.Multiplicity.PositiveEntry}

\begin{definition}[signed rational-function site]
\label{def:signed-site}
Fix one factor $p_j$ in a presentation of a reduced numerator or
denominator.  A \emph{signed site} records its positive multiplicity $m_j$,
the sign $+$ for a numerator factor or $-$ for a denominator factor, and a
finite observation in the sense of
\Defref{def:finite-observation-separation}, with the following additional
data.
\begin{enumerate}[label=(\roman*)]
\item If $\deg p_j=1$, the observation is an exact terminal $q$, together
with a certificate that the opposite reduced cofactor does not vanish at
$q$.
\item If $\deg p_j\ge2$, the observation is one of the $\deg p_j$ RATs for
$p_j$, together with
finite boundary certificates showing that the full reduced quotient has
winding $+4m_j$ on its returned boxes in the numerator case, or $-4m_j$ in
the denominator case.
\end{enumerate}
\end{definition}

\leanref{Purelib.RationalFunctions.Sites.PositiveSideInput.FactorObservation;
Purelib.RationalFunctions.SignedMultiplicity.TrackSiteInput}

The rational-function theorem uses the box winding of
\Defref{def:rational-function-winding}.

\begin{theorem}[finite constructive rational-function action principle]
\label{thm:rfap-general}
For displayed $F,G\in\Qi[Z]$ with $G\ne0$, a total finite algorithm returns
one of the following.
\begin{enumerate}[label=(\alph*)]
\item If $F=0$, it returns the exact zero-numerator certificate.  It assigns
no boundary winding to the identically zero quotient.
\item If $F\ne0$, it returns a reduced quotient presentation in the sense of
\Defref{def:reduced-quotient}; finite squarefree-multiplicity presentations
of $F_0$ and $G_0$ in the sense of
\Defref{def:squarefree-multiplicity-presentation}; and all their signed
sites in the sense of \Defref{def:signed-site}.  The sites are jointly
separated as finite observations in the sense of
\Defref{def:finite-observation-separation}.

The degree-weighted positive numerator multiplicities total $\deg F_0$, and
the degree-weighted absolute denominator multiplicities total $\deg G_0$, as in
\Defref{def:squarefree-multiplicity-presentation}.  The common coefficient
isolation square $K_{F_0,G_0}$ of
\Defref{def:common-coefficient-isolation-square} is boundary-clear for both
cofactors and satisfies
\[
             \Wind_{F_0/G_0}(K_{F_0,G_0})
                    =\deg F_0-\deg G_0
\]
in the sense of \Defref{def:rational-function-winding}.
\end{enumerate}
For every higher squarefree factor, any competing winding-one RAT supported
in that factor's coefficient isolation square has a deterministically
computed owner among the returned tracks.  Consequently every
pairwise-separated competing family for that factor has cardinality at most
its degree.
\end{theorem}

\leanref{Purelib.RationalFunctions.Sites.prove;
Purelib.Papers.RFAP.Paper.rationalFunctionArgumentPrinciple}

The multiplicity presentations are constructed in
\cref{sec:general-polynomials}; the remaining signed-site construction and
the proof of \Thmref{thm:rfap-general} are completed in \cref{sec:rfap}.

\subsection{Proof-route glossary}
\label{sec:proof-route-glossary}

These entries form a reader's map, not a second set of formal definitions:
their dependency order makes the direction of the proof visible.

\paragraph{Displayed data.}
An object is \emph{displayed} when it is given by finite rational or integer
data on which the stated calculation can actually be run.  A displayed
polynomial is its normalized finite coefficient list; a displayed point has
Gaussian-rational coordinates.

\paragraph{Certificate.}
A \emph{certificate} is displayed data together with finite checks implying
the claimed property.  Later algorithms consume the certificate itself and
return further certificates.

\paragraph{Polynomial identity.}
A \emph{polynomial identity} is a certificate that the normalized
coefficient lists agree, as fixed in \cref{sec:polynomial-equality}.
Substitution may turn such an identity into an
identity of displayed values, but value agreement is not polynomial
equality.

\paragraph{Squarefree certificate.}
The B\'ezout identity $Uf+Vf'=1$ is the \emph{squarefree certificate}.  Finite
coefficient bounds for $U$ and $V$ turn it into quantitative squarefreeness:
if the displayed value of $f$ is small in a box, the displayed value of
$f'$ there cannot also be small.  The resulting conditional positive lower
bound for $\Norm{f'}$ is the \emph{derivative margin} used by the Newton step.

\paragraph{Quadrance, projective quadrance, and coefficient bound.}
For a displayed Gaussian rational $z=x+iy$, we write
$\Norm z=x^2+y^2\in\Q_{\ge0}$ for its norm, which will also refer to as its \emph{quadrance}. A
\emph{Hermitian projective quadrance} is the normalized rational error
\[
             \chi^2(z,w)=
             \frac{\Norm{z-w}}{(1+\Norm z)(1+\Norm w)};
\]
it is the bounded metric quantity controlled by the rational boxes.  A
\emph{coefficient bound} is a rational number computed by a finite fold over
a polynomial's coefficients and proved to bound a stated quadrance on a
rational interval or box.  A \emph{quadrance-Lipschitz certificate} packages
such bounds for $f(z)-f(w)$ and its first-order Taylor remainder.  A
\emph{Horner majorant} is a coefficient bound computed by running a rational
Horner recursion on the displayed coefficient list; it is not a maximum over
an interval or box.

\paragraph{Rational box, anchor, and ambient box.}
A \emph{rational box} is four rational endpoints; its \emph{anchor} is its
lower-left Gaussian-rational corner, and its \emph{squared diagonal} is the
sum of the squares of its two side lengths.  An \emph{ambient box} is a fixed
outer box in which coefficient bounds have been certified and inside which
working boxes are refined.  A \emph{support record} is the same four-endpoint
shape used uniformly for separation: a RAT supplies one of its boxes, while
an exact terminal supplies the degenerate box at its displayed value.

\paragraph{Boundary clearance.}
A box is \emph{boundary-clear} for a polynomial when a certificate gives one
positive rational lower bound for its quadrance on all four displayed edge
parametrizations.  Clearance is what permits the winding calculation below
and is preserved when a cleared edge is reused by subdivision.  Combining
clearance with coefficient bounds gives a \emph{protected boundary radius}:
within that rational coordinate distance, a new boundary still cannot make
$f$ vanish.  A \emph{protected parent} is a proper boundary-clear box together
with its refinement into a fixed ambient box; it retains exactly the finite
local data needed for residual-driven Newton extraction, whether or not its
winding is recorded.

\paragraph{Signed remainder chain and Cauchy index.}
A \emph{signed remainder chain} is the finite Euclidean remainder program
for two rational coordinate polynomials.  Its endpoint sign variations give
their \emph{Cauchy index}, a signed count of crossing events on a displayed
interval.

\paragraph{Box winding.}
The \emph{box winding} $\Wind_4$ is the integer obtained by adding the four
Cauchy indices on a boundary-clear box.  The subscript records the retained
factor of four; a \emph{winding-one} box has $\Wind_4=4$.

\paragraph{Eliminant, safe cut, and subdivision.}
An \emph{eliminant} is a displayed nonzero rational polynomial whose nonzero
value at a coordinate produces a boundary-clearance certificate for the
corresponding vertical or horizontal line.  A bounded scan finds such a
\emph{safe cut}.  Safe cuts divide a boundary-clear box into a finite grid of
boundary-clear boxes, and winding additivity distributes the parent's box
winding among those children.

\paragraph{Leading--tail dominance, global winding budget, and local
extraction.}
Outside the \emph{coefficient isolation square}, the polynomial's displayed
\emph{leading monomial} dominates its \emph{lower tail}; their explicit
quadrance gap is the \emph{leading--tail dominance margin}.  On the boundary
of that square this gives the \emph{global winding budget}
$\Wind_4=4\deg f$.  \emph{Local extraction} repeatedly subdivides a child
with nonzero winding and uses finite Newton and Taylor certificates to obtain
a \emph{winding-one root state}: a rational box carrying both
boundary-clearance and winding-one certificates.

\paragraph{Newton centre, Newton band, and finite Rouch\'e certificate.}
When the derivative margin proves $f'(a)\ne0$, the \emph{Newton centre} is
the displayed Gaussian rational $c=a-f(a)/f'(a)$.  At radius $r>0$, a
\emph{Newton band} is a rational rectangle whose left, right, bottom, and top
sides are chosen respectively at coordinate distances between $r$ and $2r$
from $c$.  Safe-cut searches choose its four sides.  A Taylor
\emph{dominance certificate} proves on the complete boundary that the affine
Newton term is larger than the nonlinear remainder; the finite Rouch\'e rule
then transfers winding four from the affine term to the selected band.

\paragraph{Root-approximation track (RAT) and finite observation.}
A \emph{RAT} packages a winding-one root state with a total refinement
program.  A request for a displayed accuracy returns a strictly smaller
boundary-clear, winding-one box and its anchor.  A \emph{finite observation}
is either such a RAT or an \emph{exact terminal}, which is a displayed
Gaussian-rational point whose polynomial value is exactly zero.

\paragraph{Residual.}
The polynomial \emph{residual} at a displayed $q$ is the rational quadrance
$\Norm{f(q)}$.  Small winding-one boxes give small residuals; with a protected
parent, a sufficiently small residual conversely constructs a nearby RAT.

\paragraph{Separation, ownership, and finite maximality.}
\emph{Separation} is a positive rational coordinate gap between the support
boxes of two observations.  Given output boxes $B_1,\ldots,B_d$ and a new
refined box $D$, an \emph{ownership certificate} gives the unique index
$j$, with $1\le j\le d$, for which $D\subset B_j$; the $j$th returned RAT is
called the \emph{owner}.  For a separated family these assignments are
distinct, giving \emph{finite
maximality}: there can be no more competing RATs than output boxes.

\paragraph{Squarefree preprocessing, multiplicity presentation, and signed
site.}
\emph{Squarefree preprocessing} is the terminating finite gcd recursion that
splits a displayed polynomial into squarefree occurrences and then folds
overlaps into pairwise-coprime factors.  Its output is a \emph{multiplicity
presentation}: a polynomial identity built from those factors and positive
integer powers.  A \emph{signed site} is an exact terminal or RAT for one such
factor, labelled by its multiplicity and by $+$ for a numerator or $-$ for a
denominator.

\paragraph{Reduced quotient and rational-function action principle.}
A \emph{reduced quotient} cancels the displayed common factor of numerator
and denominator and certifies the remaining pair as coprime.  Applying the
polynomial construction to its multiplicity presentations produces jointly
separated signed sites; adding their signed winding contributions is the
finite rational-function action principle.

\paragraph{Protected quotient residual.}
At a numerator site for a reduced quotient $P/Q$, a finite B\'ezout
certificate, a source RAT, and its current state form a \emph{local quotient
context}.  Its finite refinement gives a \emph{protected state} on which the
opposite polynomial $Q$ has displayed positive lower and finite upper
quadrance bounds.  The \emph{quotient residual} is $\Norm{P(q)/Q(q)}$.  At a
denominator site the roles are exchanged: the shrinking quantity is the
reciprocal residual $\Norm{Q(q)/P(q)}$, and the protected opposite polynomial
is $P$.

\section{Finite polynomial bounds over \texorpdfstring{$\Qi$}{Q(i)}}
\label{sec:finite-bounds}

This section computes coefficient bounds for displayed polynomials on
rational boxes.  These bounds turn squarefree certificates into the
boundary-clearance and local-extraction certificates used later.

The estimates in this section have three simple jobs.  The coefficient
$L_1$ controls how much a polynomial value can change across a small box;
$L_2$ controls the error made by replacing the polynomial there with its
affine Taylor expression; and the pair $(\tau,\eta)$ turns squarefreeness
into the implication ``small value, non-small derivative.''  Later
constructions repeatedly choose a rational scale at which these three
controls are simultaneously valid.  The formulas are lengthy because every
constant is computed from coefficients, not because they introduce three
different notions of approximation.

\begin{lemma}[finite quadrance algebra]
\label{lem:finite-quadrance-algebra}
For displayed $z,w,u,v\in\Qi$,
\[
 \Norm{zw}=\Norm z\Norm w,
 \qquad
 \Norm{u+v}\le 2\bigl(\Norm u+\Norm v\bigr).
\]
For every displayed finite list $u_1,\dots,u_m\in\Qi$,
\begin{equation}
\label{eq:finite-sum}
 \Norm{u_1+\cdots+u_m}\le m\sum_{j=1}^{m}\Norm{u_j}.
\end{equation}
\end{lemma}

\begin{proof}
Expand the two rational coordinates for the product identity.  The
two-term inequality follows after subtracting the left side from the right:
the difference is a sum of two rational squares.  For the finite-list bound,
apply
\[
       (r_1+\cdots+r_m)^2\le m\sum_{j=1}^{m}r_j^2
\]
separately to the real and imaginary coordinate lists and add the two
inequalities.  The difference in the displayed scalar inequality is a finite
sum of squares.
\end{proof}

\leanref{Purelib.Data.Gaussian.normSq_mul;
Purelib.Data.Gaussian.normSq_add_le_two;
Purelib.LinearAlgebra.Norm.normSq_gaussianSum_le}

\begin{lemma}[rational-box norm bound]
\label{lem:box-norm-bound}
For a rational box $B$, put
\[
  R_B=2\max\{x_-^2,x_+^2,y_-^2,y_+^2\}.
\]
Every displayed $z\in B$ satisfies $\Norm z\le R_B$.
\end{lemma}

\begin{proof}
The rational interval inequalities for the two coordinates give
\[
 \Ree(z)^2\le\max\{x_-^2,x_+^2\},\qquad
 \Imm(z)^2\le\max\{y_-^2,y_+^2\}.
\]
Both right-hand sides are at most the maximum of all four endpoint squares.
Adding the two bounds proves the claim.
\end{proof}

\leanref{Purelib.Polynomial.Geometry.Boxes.Box.normSqBound;
Purelib.Polynomial.Geometry.Boxes.Box.normSq_le_normSqBound}

Quadrance is squared Euclidean magnitude kept rational, so no square root is
required.  The Hermitian projective quadrance used here is the
Gaussian-Hermitian analogue, on the affine chart, of Wildberger's algebraic
projective separation of homogeneous directions: $z,w$ represent $[z:1]$
and $[w:1]$.  Its displayed rational value is the error measure used below.

\begin{proposition}[box calibration of Hermitian projective quadrance]
\label{prop:box-projective-calibration}
For a rational box $B$ and every displayed $z,w\in B$,
\begin{equation}
\label{eq:projective-box-upper}
       \chi^2(z,w)\le\Norm{z-w}\le\diamq(B).
\end{equation}
At the distinguished opposite corners the exact comparison is
\begin{equation}
\label{eq:box-projective-exact}
 \diamq(B)=
 (1+\Norm{q_B})(1+\Norm{v_B})\chi^2(q_B,v_B).
\end{equation}
\end{proposition}

\begin{proof}
The denominator in \cref{eq:hermitian-projective-affine} is at least one.
Membership in $B$ bounds the real and imaginary coordinate differences by the two side
lengths, proving \cref{eq:projective-box-upper}.  For $q_B,v_B$, the numerator
$\Norm{q_B-v_B}$ is exactly $\diamq(B)$; clearing the explicitly positive
denominator gives \cref{eq:box-projective-exact}.
\end{proof}

\leanref{Purelib.Polynomial.Geometry.Chordal.affineChordalSq_le_box_diagonal;
Purelib.Polynomial.Geometry.Chordal.diagonalSq_eq_cornerWeight_mul_affineChordalSq}

\begin{lemma}[finite Horner evaluation bound]
\label{lem:eval-bound}
Let $p\in\Qi[Z]$ and let $B$ be a rational box.  Write the displayed
high-to-low coefficient list of $p$ as
$[a_n,\ldots,a_0]$.  Starting from $M_0=0$, perform the finite recurrence
\[
        M_{k+1}=2\bigl(M_kR_B+\Norm{a_{n-k}}\bigr).
\]
Then every $z$ satisfying the predicate of $B$ obeys
$\Norm{p(z)}\le M_{n+1}$.  Simultaneously folding
$S_{k+1}=2(S_kR_B+M_k)$ gives the finite divided-difference coefficient used
below.  Both outputs are rational computations on the finite coefficient
list.
\end{lemma}

\begin{proof}
Induct over the Horner fold and use
$\Norm{uv}=\Norm u\Norm v$ and
$\Norm{u+v}\le2(\Norm u+\Norm v)$ at each step.  Applying the same induction
to two Horner evaluations gives the displayed $S_k$ recurrence.
\end{proof}

\leanref{Purelib.Polynomial.Bounds.denseHornerNormBounds;
Purelib.Polynomial.Bounds.normSq_eval_le_denseEvalNormBound}

The same finite fold has the following one-variable rational specialization.
We name it separately because it is the form used below for B\'ezout
cofactors on parameter intervals.

\begin{definition}[rational-interval Horner majorant]
\label{def:interval-horner-majorant}
Let $a\le b$ be displayed rationals, write $I=[a,b]$, and put
\[
                         R_I=\max\{a^2,b^2\}.
\]
For $W\in\Q[T]$ with displayed high-to-low coefficient list
$[w_n,\ldots,w_0]$, start from $H_0=0$ and perform the finite recurrence
\[
                  H_{k+1}=2\bigl(H_kR_I+w_{n-k}^2\bigr).
\]
The resulting rational
\[
                         \mathcal H_I(W):=H_{n+1}
\]
is the \emph{Horner majorant} of $W$ on $I$.  It is computed from the finite
coefficient list and the two endpoints; it is not a maximum taken over the
parameters satisfying the interval predicate.
\end{definition}

\begin{lemma}[rational-interval Horner bound]
\label{lem:interval-horner-majorant}
For the data of \cref{def:interval-horner-majorant} and every displayed
$t\in\Q$ satisfying $a\le t\le b$,
\[
                         W(t)^2\le\mathcal H_I(W).
\]
The majorant is a nonnegative rational computed by one finite fold over the
displayed coefficient list.
\end{lemma}

\begin{proof}
The endpoint inequalities give $t^2\le R_I$.  Let $E_0=0$ and
$E_{k+1}=E_kt+w_{n-k}$ be the corresponding Horner evaluation.  If
$E_k^2\le H_k$, then
\[
 \begin{aligned}
 E_{k+1}^2
   &\le 2\bigl(E_k^2t^2+w_{n-k}^2\bigr)\\
   &\le 2\bigl(H_kR_I+w_{n-k}^2\bigr)=H_{k+1}.
 \end{aligned}
\]
Finite induction gives $W(t)^2=E_{n+1}^2\le H_{n+1}$.  The same induction,
starting from $H_0=0$, proves $H_k\ge0$ at every step.
\end{proof}

\leanref{Purelib.Polynomial.Geometry.PathClearance.evalSqBoundStep;
Purelib.Polynomial.Geometry.PathClearance.evalSqBound;
Purelib.Polynomial.Geometry.PathClearance.eval_sq_le_evalSqBound}

\begin{lemma}[explicit geometric depth]
\label{lem:reciprocal-search}
For displayed rationals $w\ge0$ and $\eps>0$, define
\[
 \operatorname{geometricDepth}(w,\eps)
       :=\left\lceil\frac{2w}{\eps}\right\rceil\in\N.
\]
Then
\[
 \left(\frac23\right)^{\operatorname{geometricDepth}(w,\eps)}w
       \le\eps.
\]
Finite minima of displayed positive rationals remain positive.
\end{lemma}

\begin{proof}
Put $N=\operatorname{geometricDepth}(w,\eps)$.
The constructive rational ceiling gives $2w/\eps\le N$.  The finite
Bernoulli induction
$(3/2)^N\ge1+N/2$ then implies
$(2/3)^Nw\le\eps$.  Positivity of a finite minimum is proved by induction on
the displayed list.
\end{proof}

\leanref{Purelib.RootIsolation.Isolation.SafeCandidates.geometricDepth;
Purelib.RootIsolation.Isolation.SafeCandidates.geometricDepth_sufficient}

\begin{lemma}[finite divided-difference bound]
\label{lem:divided-difference-bound}
Let $p\in\Qi[Z]$ and let $B$ be a rational box.  Write
$p(Z)=\sum_{j=0}^{n}a_jZ^j$.  The finite identity
\begin{equation}
\label{eq:divdiff}
 p(z)-p(w)=(z-w)D_p(z,w),\qquad
 D_p(z,w)=\sum_{j=1}^{n}a_j\sum_{r=0}^{j-1}z^{j-1-r}w^r
\end{equation}
computes a rational $L_1(p,B)>0$ such that, whenever $z,w\in B$,
\begin{equation}
\label{eq:lip1}
 \Norm{p(z)-p(w)}\le L_1(p,B)\Norm{z-w}
\end{equation}
and hence supplies the certificate of
\Defref{def:quadrance-lipschitz}.
\end{lemma}

\begin{proof}
Factor $z^j-w^j$ as
$(z-w)\sum_{r=0}^{j-1}z^{j-1-r}w^r$ and sum the finitely many monomial
identities to obtain \eqref{eq:divdiff}.  The resulting $D_p$ is a displayed
finite bivariate polynomial.  Apply the Horner bounds of
\cref{lem:eval-bound} in the two variables using the box norm bound of
\cref{lem:box-norm-bound}, and combine the result with norm
multiplicativity from \cref{lem:finite-quadrance-algebra}.
\end{proof}

\leanref{Purelib.Polynomial.Bounds.denseLipschitzBound;
Purelib.Polynomial.Bounds.normSq_eval_sub_eval_le_denseLipschitzBound}

\begin{lemma}[finite Taylor remainder bound]
\label{lem:finite-taylor-bound}
Let $p\in\Qi[Z]$ and let $B$ be a rational box.  There is a displayed finite
polynomial $T_p(z,w)$ satisfying
\begin{equation}
\label{eq:taylor}
 p(z)-p(w)-p'(w)(z-w)=(z-w)^2T_p(z,w),
\end{equation}
and the same finite coefficient calculation computes $L_2(p,B)>0$ such
that, whenever $z,w\in B$,
\begin{equation}
\label{eq:lip2}
 \Norm{p(z)-p(w)-p'(w)(z-w)}
       \le L_2(p,B)\Norm{z-w}^2.
\end{equation}
\end{lemma}

\begin{proof}
Expand every monomial $z^j=(w+(z-w))^j$ by the finite binomial identity and
remove its constant and linear terms.  Every remaining term contains
$(z-w)^2$, giving the displayed $T_p$.  Apply the same finite Horner bound as
in \cref{lem:divided-difference-bound} and use
$\Norm{(z-w)^2}=\Norm{z-w}^2$.
\end{proof}

\leanref{Purelib.Polynomial.Bounds.denseTaylorRemainderBound;
Purelib.Polynomial.Bounds.normSq_dense_taylor_remainder_le_denseTaylorRemainderBound}

The first estimate is the squared-distance form of the usual Lipschitz
bound, retained so that every quantity remains rational.  Both results are
finite polynomial identities followed by rational inequalities.

The identity $Uf+Vf'=1$ is the useful content of the squarefree certificate.
Direct substitution says that $f$ and $f'$ cannot both vanish at one
displayed argument; finite bounds on $U,V$ strengthen this to the numerical
rule needed below: a small value of $f$ forces a non-small derivative.  This
is the local derivative control used by the extraction step.

\begin{lemma}[finite squarefree-cofactor bounds]
\label{lem:squarefree-cofactor-bounds}
Let the squarefree certificate supply displayed $U,V\in\Qi[Z]$ satisfying
the identity
\[
                         Uf+Vf'=1.
\]
For any rational box $B$, apply the Horner recurrence of
\cref{lem:eval-bound} separately to the high-to-low coefficient lists of
$U,V$, obtaining $H_U,H_V$, and put
\[
                       M_U=\max\{1,H_U\},
                       \qquad
                       M_V=\max\{1,H_V\}.
\]
Then $M_U,M_V\ge1$, and every displayed $z\in B$ satisfies
\[
                  \Norm{U(z)}\le M_U,
                  \qquad
                  \Norm{V(z)}\le M_V.
\]
\end{lemma}

\begin{proof}
The two applications of \cref{lem:eval-bound} give
$\Norm{U(z)}\le H_U$ and $\Norm{V(z)}\le H_V$.  Each output is at most its
maximum with $1$, which also proves the two positivity bounds.
\end{proof}

\leanref{Purelib.Polynomial.Bounds.denseHornerNormBounds;
Purelib.Polynomial.Bounds.normSq_eval_le_denseEvalNormBound}

From these bounds define
\begin{equation}
\label{eq:bez-margins}
       \tau=\frac{1}{16M_U},\qquad \eta=\frac{1}{16M_V}.
\end{equation}

\begin{lemma}[derivative margin from the Bézout certificate]
\label{lem:derivative-margin}
For every rational $z$ satisfying the predicate of $B$,
\[
  \Norm{f(z)}\le\tau \quad\Longrightarrow\quad \Norm{f'(z)}\ge\eta.
\]
\end{lemma}

\begin{proof}
From $1=U(z)f(z)+V(z)f'(z)$ and \eqref{eq:finite-sum},
\[
  1\le 2M_U\Norm{f(z)}+2M_V\Norm{f'(z)}.
\]
If the displayed conclusion failed, then both products
$M_U\Norm{f(z)}$ and $M_V\Norm{f'(z)}$ would be at most $1/16$.
The right side would therefore be at most $1/4$, contradicting $1\le1/4$.
\end{proof}

\leanref{Purelib.RootIsolation.Isolation.LocalExtraction.derivativeMargin;
Purelib.RootIsolation.Isolation.LocalExtraction.derivativeMargin_le_of_small_residual;
Purelib.RootIsolation.Isolation.LocalExtraction.derivative_ne_zero_of_small_residual}

\section{Sturm chains as finite programs}

This section constructs the signed remainder chains whose endpoint
variations give Cauchy indices.  These indices are the inputs to box winding.

We use the signed Euclidean remainder sequence as the \emph{definition} of
the index over $\Q$.  The proof below represents every crossing event by a
rational isolating interval and a refinement routine.  Throughout this paper,
every Cauchy--Sturm index is
normalized at twice the conventional scale: an interior crossing contributes
$\pm2$ and an endpoint crossing contributes $\pm1$.  This normalization is
built directly into the endpoint variation defined below, avoiding an
intermediate halving followed by a doubling in the index.  We write simply
$\Var_a$ and $\Ind_a^b$, since no competing undoubled variation or index is
used.

\begin{definition}[rational interval partition]
\label{def:rational-interval-partition}
A rational interval $I=[a,b]$ is the endpoint pair $(a,b)$ together with the
predicate $a\le t\le b$ on a displayed rational $t$.  An \emph{interval
partition} rooted at $I$ is a finite binary tree of such endpoint records.
An internal node $J=[u,v]$ stores a displayed rational cut $c$ with $u<c<v$
and has the two children $[u,c]$ and $[c,v]$.  A \emph{leaf} is a terminal
node, and the finitely many leaves of the tree are extracted as a finite
list.
\end{definition}

\leanref{Purelib.RootIsolation.Sturm.EventPartition.Cell;
Purelib.RootIsolation.Sturm.EventPartition.SafeSplit;
Purelib.RootIsolation.Sturm.EventPartition.Subdivision;
Purelib.RootIsolation.Sturm.EventPartition.Subdivision.leaves}

\begin{definition}[central-third cut]
\label{def:central-third-cut}
For a proper rational interval $J=[u,v]$, its \emph{open central third} is
specified by the two rational bounds
\[
 u+\frac{v-u}{3}<c< u+2\frac{v-u}{3}.
\]
A \emph{central-third cut} is a displayed rational $c$ satisfying these
inequalities.
\end{definition}

\begin{lemma}[central-third contraction]
\label{lem:central-third-contraction}
If $c$ is a central-third cut of $J=[u,v]$, then its two child widths obey
\[
 c-u<\frac23(v-u),\qquad v-c<\frac23(v-u).
\]
\end{lemma}

\begin{proof}
Subtract the left central-third inequality from $c$, and subtract the right
one from $v$.
\end{proof}

\begin{lemma}[finite safe-cut search]
\label{lem:safe-central-third-cut}
Let $H\in\Q[T]$ be a displayed nonzero polynomial, let $D\in\N$ be a
displayed bound with $\deg H\le D$, and let $J$ be a proper rational
interval.  A bounded search returns an $H$-safe central-third cut $c$:
$c$ lies in the open central third of $J$ and $H(c)\ne0$.
\end{lemma}

\begin{proof}
Place $D+1$ distinct, equally spaced rational candidates strictly between
the two central-third bounds and evaluate $H$ on that finite list.  If all
evaluations vanished, repeated exact division by the corresponding linear
factors $T-c$ would force $H$ to have degree at least $D+1$.  Hence one
candidate has nonzero evaluation.  Return the first such candidate in the
displayed list.
\end{proof}

\leanref{Purelib.RootIsolation.Isolation.SafeCandidates.centralLower;
Purelib.RootIsolation.Isolation.SafeCandidates.centralUpper;
Purelib.RootIsolation.Isolation.SafeCandidates.centralBounds;
Purelib.RootIsolation.Isolation.SafeCandidates.centralSearchValue;
Purelib.RootIsolation.Isolation.SafeCandidates.centralSearchValue_spec;
Purelib.RootIsolation.Isolation.SafeCandidates.centralSearchValue_contraction}

The central-third condition and a finite guard polynomial solve two
independent problems at once.  Choosing near the middle guarantees visible
geometric progress, while avoiding the finitely many guarded zeros makes the
new endpoint safe for the sign calculation encoded by $H$.  The scan of
$D+1$ candidates is the finite substitute for saying ``choose a generic
cut.''  For a signed chain, the product of its members will be exactly such
a guard.

\subsection{Signs, variations, and signed remainders}

\begin{definition}[doubled sign variation]
\label{def:sign-variation}
For $r\in\Q$, $\operatorname{sgn}(r)\in\{-1,0,1\}$ is decidable.  For a
finite list $s=(s_0,\dots,s_n)$ define
\begin{equation}
\label{eq:variation}
 \Var(s)=\sum_{j=1}^{n}|\operatorname{sgn}(s_{j-1})-
                                  \operatorname{sgn}(s_j)|.
\end{equation}
This doubled variation is an integer: a direct switch between positive and
negative contributes two, while a comparison between zero and a nonzero
value contributes one.  For polynomials $S_0,\dots,S_n\in\Q[T]$ and
$a\in\Q$, write
\[
 \Var_a(S_0,\dots,S_n)=\Var(S_0(a),\dots,S_n(a)),\qquad
 \Var_a^b=\Var_a-\Var_b.
\]
Note this is twice the normal definition, but removing the $1/2$ here removes the
need to double the contributions for the Cauchy--Sturm index calculation.
\end{definition}

\leanref{Purelib.RootIsolation.Sturm.signJump; Purelib.RootIsolation.Sturm.variation2;
Purelib.RootIsolation.Sturm.endpointVariation2}

To interpret this number, freeze the parameter at $t$ and read the signs of
the chain from left to right.  $\Var_t$ records sign switches on the
doubled scale fixed at the start of this section: a zero at one end of a
comparison is a half-switch on the conventional scale and therefore
contributes one here.  As $t$ moves, the signed remainder identities force
changes at an interior chain member to
cancel between its two neighbours.  A net change survives only at the first
member $S_0$, namely the denominator of $R/S$ after GCD reduction.  Thus
$\Var_a-\Var_b$ is a finite algebraic ledger of the oriented
pole-crossing events of $R/S$ between the two endpoints.  It extracts from
endpoint signs the information ordinarily read from the graph of the
quotient.

\begin{definition}[reduced signed remainder chain]
\label{def:reduced-signed-remainder-chain}
Given $R,S\in\Q[T]$ with $S\ne0$, run the finite extended Euclidean program
to remove their displayed common GCD $g$, and put $S_0=S/g$.  If $R=0$,
the returned chain is $[S_0]$.  Otherwise put $S_1=R/g$ and repeatedly
perform
\begin{equation}
\label{eq:srem}
       S_{j+1}=Q_jS_j-S_{j-1},\qquad \deg S_{j+1}<\deg S_j,
\end{equation}
stopping immediately before the zero remainder.  The resulting finite list
of nonzero polynomials is the \emph{reduced signed remainder chain} of
$(R,S)$.
\end{definition}

\begin{lemma}[termination and coprimality of the signed chain]
\label{lem:signed-chain-termination}
The computation in \cref{def:reduced-signed-remainder-chain} terminates after
finitely many steps.  Every consecutive pair in the returned chain is
coprime, and hence no two consecutive members can vanish at the same
displayed rational parameter.
\end{lemma}

\begin{proof}
At every nonterminal Euclidean step, the natural degree count of the next
remainder is strictly smaller than that of the current divisor.  Recursion
on that displayed natural number therefore terminates.  The initial reduced
pair comes with the B\'ezout identity returned by extended Euclid, and the
identity in \eqref{eq:srem} transports coprimality to each consecutive pair.
Indeed, if consecutive members $S_j,S_{j+1}$ have displayed cofactors $U,V$
with
\[
                       US_j+VS_{j+1}=1,
\]
then substitution of any displayed $t\in\Q$ preserves this polynomial
identity and gives
\[
             U(t)S_j(t)+V(t)S_{j+1}(t)=1.
\]
If both chain members vanished at $t$, the left-hand side would be zero,
contradicting $0\ne1$ in $\Q$.
\end{proof}

\leanref{Purelib.RootIsolation.Sturm.signedRemainderChain;
Purelib.RootIsolation.Sturm.signedRemainderChain_length_le;
Purelib.RootIsolation.Sturm.signedRemainderChain_all_nonzero;
Purelib.RootIsolation.Sturm.noAdjacentZero_signedRemainderChain;
Purelib.RootIsolation.Sturm.RootCount.reducedChain}

\begin{corollary}[safe cut for a signed remainder chain]
\label{cor:sturm-safe-cut}
For a reduced signed remainder chain $S_0,\ldots,S_m$, put
\[
                         G:=\prod_{j=0}^{m}S_j.
\]
There is a finite central-third search returning a cut $c$ at which every
$S_j(c)$ is nonzero.  Both child intervals contract by the factor in
\cref{lem:central-third-contraction}.
\end{corollary}

\begin{proof}
By definition, all members of the returned chain are nonzero polynomials,
so their finite product $G$ is nonzero.  Apply
\cref{lem:safe-central-third-cut} to $G$ and its computed degree.  Since a
nonzero product has no zero factor, $G(c)\ne0$ proves $S_j(c)\ne0$ for every
member of the finite chain.  The contraction conclusion is
\cref{lem:central-third-contraction}.
\end{proof}

\leanref{Purelib.RootIsolation.Sturm.RootCount.chainGuard;
Purelib.RootIsolation.Sturm.RootCount.safeCut;
Purelib.RootIsolation.Sturm.RootCount.safeCut_avoids_chain;
Purelib.RootIsolation.Sturm.RootCount.safeCut_contraction}

\begin{definition}[finite Cauchy--Sturm index]
\label{def:index}
For rational endpoints $a,b$, set the index to zero when $R=0$ or $S=0$.
Otherwise, for the chain $S_0,\ldots,S_m$ of
\cref{def:reduced-signed-remainder-chain}, define
\[
       \Ind_a^b\!\left(\frac{R}{S}\right)
          =\Var_a^b(S_0,S_1,\dots,S_m).
\]
This is exactly the integer stored by the finite program on rational
coefficients and endpoints.  The doubled
normalization is already built into \eqref{eq:variation}, so no additional
factor is inserted here.  Reversing the endpoints changes the sign by
definition.
\end{definition}

\leanref{Purelib.RootIsolation.Sturm.cauchyIndex2}

The word \emph{index} therefore refers to a signed event count, not to a
location.  A nonzero index certifies that the interval contains an
uncancelled crossing event.  Its sign records orientation, and its magnitude
is the net count retained by the refinement machinery.

The local event scanner used below is entirely finite.  It
refines a rational interval to a coefficient-computed natural depth, records
stable signs for the complete reduced chain on every leaf, and checks the
resulting endpoint variation there.

\begin{theorem}[finite Sturm calculation]
\label{thm:sturm}
For the displayed $R,S$ of \cref{def:index}, the index and the safe
partition program of \cref{cor:sturm-safe-cut} satisfy the following
identities.
\begin{enumerate}[label=(\alph*)]
\item Subdivision and reversal:
\[
 \Ind_a^b(R/S)+\Ind_b^c(R/S)=\Ind_a^c(R/S),\qquad
 \Ind_b^a(R/S)=-\Ind_a^b(R/S).
\]
\item If the displayed denominator is nonzero at both endpoints, the
reduced index is even.
\item If $S\ne0$, then for an input rational interval $I=[a,b]$ a
coefficient-computed finite event partition returns a finite leaf list.  For each leaf
$J=[u_J,v_J]$, its classification stores formal signs for
the GCD-reduced numerator and denominator, proves that at least one of those
signs is nonzero, proves their stability against the displayed endpoint
evaluations at $u_J$ and $v_J$, and identifies
$\Ind_{u_J}^{v_J}(R/S)$ with the resulting finite endpoint event formula.
The partition also proves the exact telescoping identity
\[
 \Ind_a^b(R/S)
 =\sum_{J\text{ a leaf}}\Ind_{u_J}^{v_J}(R/S).
\]
\end{enumerate}
\end{theorem}

\begin{proof}
Part (a) follows by addition and subtraction of the three displayed endpoint
variations.  For (b), the finite parity lemma for a coprime signed remainder
chain is applied at the two clear endpoints.  For (c), use the safe
central-third cut of \cref{cor:sturm-safe-cut} at every internal node.
A coefficient-computed depth makes the finite sign classifier succeed on
every terminal interval.  Its stored stability equalities give the endpoint
event formula for that leaf.  Subdivision distributes the parent index to
its two children, and induction over the same finite depth telescopes the
leaf list.
\end{proof}

\leanref{Purelib.RootIsolation.Sturm.cauchyIndex2_add; Purelib.RootIsolation.Sturm.cauchyIndex2_reverse;
Purelib.RootIsolation.Sturm.RootCount.reducedCauchyIndex2_even;
Purelib.RootIsolation.Sturm.EventPartition.classifyEvents;
Purelib.RootIsolation.Sturm.EventPartition.leafEvent_index2;
Purelib.RootIsolation.Sturm.EventPartition.eventSubdivision_cellIndex2_sum}

\subsection{Univariate approximation tracks and positive lower bounds}

\begin{definition}[one-dimensional nonzero-event track]
\label{def:one-dimensional-track}
A nonzero-event track for a displayed nonzero polynomial $p$ is a rational
interval $[u,v]$ carrying the explicit certificate
\[
                         \Ind_u^v(p'/p)\ne0,
\]
together with the deterministic refinement program from
\cref{thm:sturm}.  A positive rational width request is answered by one
finite nested interval of at most that width, again carrying an explicit
nonzero-index certificate.
\end{definition}

\leanref{Purelib.RootIsolation.Sturm.OneDimensional.ZeroEventTrack}

The retained certificate is exactly the nonzero signed event that survives
each refinement request; it is not silently strengthened to count one.

Termination uses only the natural geometric-depth computation of
\cref{lem:reciprocal-search}.  Each safe
central-third cut of \cref{cor:sturm-safe-cut} contracts the retained
interval by the fixed rational factor $2/3$;
the depth needed to cross a positive requested width is computed before the
finite recursion starts.

\begin{theorem}[finite one-dimensional event machinery]
\label{thm:1d-isolation}
For rational polynomials $R,S$ with $S\ne0$ and a rational input interval
$I=[a,b]$, the finite event-tree algorithm of \cref{thm:sturm} classifies
all leaves $J=[u_J,v_J]$ and proves
\[
 \Ind_a^b(R/S)
   =\sum_{J\text{ a leaf}}\Ind_{u_J}^{v_J}(R/S).
\]
For every explicitly selected nonzero-event track of
\cref{def:one-dimensional-track}, a total refinement function answers each
positive rational width request by a nested interval of at most that width
while retaining a nonzero Sturm index.  Successive answers are nested by
composition of their displayed rational refinement certificates.
\end{theorem}

\begin{proof}
Let $V_t$ denote the doubled endpoint variation of the GCD-reduced signed
chain at $t$.  At an internal interval $[u,v]$, the safe cut $c$ of
\cref{cor:sturm-safe-cut} gives the exact integer calculation
\[
\begin{aligned}
 \Ind_u^v(R/S)
   &=V_u-V_v\\
   &=(V_u-V_c)+(V_c-V_v)\\
   &=\Ind_u^c(R/S)+\Ind_c^v(R/S).
\end{aligned}
\]
Induction over the finite binary tree now replaces each parent by the sum of
its two children.  Write
$\mathcal I([r,s])=\Ind_r^s(R/S)$, and let $L$ and $M$ be the finite leaf
lists of the left and right subtrees.  Associativity of integer addition and
the finite-list identity
\[
 \sum_{J\in L\mathbin{++}M}\mathcal I(J)
   =\sum_{J\in L}\mathcal I(J)+\sum_{J\in M}\mathcal I(J)
\]
give exactly the displayed leaf sum.  This is the subdivision part of
\cref{thm:sturm} written out algebraically.

Now let a track for $p$ have current interval $[u,v]$ and nonzero count
\[
                         C=\Ind_u^v(p'/p)\ne0.
\]
At its safe cut $c$, compute the two displayed integers
\[
 C_L=\Ind_u^c(p'/p),\qquad C_R=\Ind_c^v(p'/p),
 \qquad C=C_L+C_R.
\]
Integer equality is decidable.  If $C_L\ne0$, retain $[u,c]$.  Otherwise
$C_L=0$, and $C_R=0$ would force $C=0$; hence retain $[c,v]$ with the
certificate $C_R\ne0$.  Thus one finite step preserves the nonzero count
and refines the current interval.

If $w_n$ is the retained width after $n$ steps, the central-third
inequalities give
\[
                 w_{n+1}<\frac23w_n,
 \qquad          w_N\le\left(\frac23\right)^Nw_0.
\]
For a request $\eps>0$, the natural number
$N=\operatorname{geometricDepth}(w_0,\eps)$ of
\cref{lem:reciprocal-search} makes the last quantity at most $\eps$.
Executing exactly those $N$ steps therefore
returns the requested nonzero-event interval.  For successive requests,
the coordinate inequalities for the second returned interval refine those
for the first, and those for the first refine the original; transitivity of
the two endpoint inequalities proves the claimed nesting.
\end{proof}

\leanref{Purelib.RootIsolation.Sturm.OneDimensional.classify;
Purelib.RootIsolation.Sturm.OneDimensional.ClassifiedPartition.indexSum;
Purelib.RootIsolation.Sturm.NonzeroEvent.step;
Purelib.RootIsolation.Sturm.NonzeroEvent.refineN_refines;
Purelib.RootIsolation.Sturm.NonzeroEvent.refineN_width_le;
Purelib.RootIsolation.Isolation.SafeCandidates.geometricDepth_sufficient;
Purelib.RootIsolation.Sturm.OneDimensional.ZeroEventTrack.request;
Purelib.RootIsolation.Sturm.OneDimensional.ZeroEventTrack.request_refines;
Purelib.RootIsolation.Sturm.OneDimensional.ZeroEventTrack.request_width_le;
Purelib.RootIsolation.Sturm.OneDimensional.ZeroEventTrack.request_countNonzero;
Purelib.RootIsolation.Sturm.OneDimensional.ZeroEventTrack.successive_refines}

The next construction converts symbolic coprimality into a numerical safety
margin.  The identity $UP+VQ=1$ says that $P$ and $Q$ cannot vanish
simultaneously.  Bounding the displayed cofactors $U,V$ on the interval says
more: the pair $(P(t),Q(t))$ cannot even become simultaneously too small.
This is the recurring mechanism by which a finite algebraic certificate
becomes boundary clearance.

\begin{lemma}[finite Bézout clearance]
\label{lem:positive-lower}
Let $P,Q\in\Q[T]$, let $I=[a,b]$, and supply displayed $U,V$ with the
polynomial identity $UP+VQ=1$.  Using
\cref{def:interval-horner-majorant}, put
\[
 M=\max\{1,\mathcal H_I(U)+\mathcal H_I(V)\},
 \qquad m=\frac1M.
\]
Then $m\in\Q_{>0}$ and
\[
             P(t)^2+Q(t)^2\ge m
\]
for every rational parameter satisfying $a\le t\le b$.  The same certificate
restricts unchanged to every displayed rational subinterval.
\end{lemma}

\begin{proof}
By \cref{lem:interval-horner-majorant},
\[
                  U(t)^2+V(t)^2\le M.
\]
The exact two-square identity gives
\[
 1=(U(t)P(t)+V(t)Q(t))^2
 \le (U(t)^2+V(t)^2)(P(t)^2+Q(t)^2)
 \le M(P(t)^2+Q(t)^2).
\]
Since $M\ge1$, division gives $m=1/M>0$ and the required lower bound.  A
displayed rational subinterval has no new members, so the same inequality
and the same $m$ restrict unchanged.
\end{proof}

\leanref{Purelib.Polynomial.Geometry.PathClearance.bezoutSqMajorant;
Purelib.Polynomial.Geometry.PathClearance.certificateOfCoprime;
Purelib.Polynomial.Geometry.PathClearance.Certificate.restrict}

\section{Algebraic winding on rational paths}
\label{sec:winding-realization}

This section combines the four Cauchy indices on a rational-box boundary
into box winding and proves the finite rules used by subdivision and the
rational-function action principle.

\begin{definition}[oriented affine edges and box boundary]
\label{def:oriented-box-boundary}
An oriented segment from $u$ to $v$ is the affine expression
$e(t)=u+t(v-u)$ under the parameter predicate $0\le t\le1$.  The boundary
of a rational box $B$ is its positively oriented four-element list
$\operatorname{edges}(B)$ of such expressions.  The notation $\partial B$
used in prose abbreviates this finite list and its four parameter predicates.
\end{definition}

Clearance is the numerical guardrail around this boundary.  Its positive
margin says not only that the boundary value avoids zero, but by how much;
that margin later permits perturbation, subdivision, and comparison without
changing the winding calculation.

The intended winding picture is familiar: traverse the four edges and watch
the direction of the polynomial or quotient value.  Its net signed number of
turns about zero is recovered here without angles or a continuous image
curve, using only finitely many signed real-axis crossings.

\begin{definition}[division-free quotient edge]
\label{def:division-free-quotient-edge}
For displayed $P,Q\in\Qi[Z]$ and an oriented edge $e$ of
\cref{def:oriented-box-boundary}, write
\[
 \begin{aligned}
 P(e(t))&=A_{P,e}(t)+iB_{P,e}(t),\\
 Q(e(t))&=A_{Q,e}(t)+iB_{Q,e}(t),
 \end{aligned}
 \qquad A_{P,e},B_{P,e},A_{Q,e},B_{Q,e}\in\Q[t].
\]
The division-free direction polynomial on the edge is
\[
 \begin{aligned}
 H_{P/Q,e}(t)
   &:=P(e(t))\overline{Q(e(t))}\\
   &=A_{P,e}A_{Q,e}+B_{P,e}B_{Q,e}
     +i\bigl(B_{P,e}A_{Q,e}-A_{P,e}B_{Q,e}\bigr)\\
   &=:C_{P/Q,e}(t)+iD_{P/Q,e}(t).
 \end{aligned}
\]
Here conjugation is performed on the two rational coordinate polynomials
of the affine edge; it is not treated as a holomorphic polynomial operation
in $Z$.
\end{definition}

\begin{lemma}[division-free quotient direction]
\label{lem:quotient-direction}
At a displayed edge parameter $t$, put $z=e(t)$.  If $Q(z)\ne0$, then
\begin{equation}
\label{eq:quotient-direction}
             \frac{P(z)}{Q(z)}
             =\frac{H_{P/Q,e}(t)}{\Norm{Q(z)}},
             \qquad \Norm{Q(z)}>0.
\end{equation}
If also $P(z)\ne0$, then $H_{P/Q,e}(t)\ne0$, so
$C_{P/Q,e}(t)$ and $D_{P/Q,e}(t)$ do not vanish simultaneously.
\end{lemma}

\begin{proof}
The Gaussian-rational identities are
\[
 Q(z)\overline{Q(z)}=\Norm{Q(z)},\qquad
 Q(z)^{-1}=\frac{\overline{Q(z)}}{\Norm{Q(z)}}.
\]
The norm is a sum of two rational squares and is strictly positive when
$Q(z)\ne0$, proving
\eqref{eq:quotient-direction}.  If $P(z)$ and $Q(z)$ are both nonzero, their
product $P(z)\overline{Q(z)}$ is nonzero in the Gaussian-rational field;
therefore its two rational coordinates cannot both be zero.
\end{proof}

\leanref{Purelib.RationalFunctions.Winding.quotientPathReal;
Purelib.RationalFunctions.Winding.quotientPathImag;
Purelib.RationalFunctions.Winding.quotientEdge}

\begin{definition}[quotient winding]
\label{def:quotient-winding}
Let $B$ carry boundary-clearance certificates for $P$ and $Q$.  The
contribution of an oriented edge and the box winding are the finite
Cauchy--Sturm calculations
\[
 \Wind_4(P/Q;e)
   :=\Ind_0^1\!\left(\frac{C_{P/Q,e}}{D_{P/Q,e}}\right),
 \qquad
 \Wind_4(P/Q;B)
   :=\sum_{e\text{ in }\operatorname{edges}(B)}\Wind_4(P/Q;e),
\]
where the index is that of \cref{def:index}.  This formula uses only the
displayed rational coordinate polynomials; it never evaluates a division
by $Q$.  The clearance certificates required in
\cref{def:rational-function-winding}, together with
\cref{lem:quotient-direction}, prove that $H_{P/Q,e}$ never has both
coordinates zero under an edge predicate.
\end{definition}

\leanref{Purelib.RationalFunctions.Winding.quotientBoundary;
Purelib.RationalFunctions.Winding.rationalQuadrupledWinding}

Concretely, along one edge the finite Sturm program reads the endpoint signs
of $(C_{P/Q,e},D_{P/Q,e})$ on rational event intervals.  A change recorded
in the sign ledger for $D_{P/Q,e}$ is the division-free surrogate for a
real-axis crossing; the stored sign data for $C_{P/Q,e}$ and the orientation
of that change determine its signed contribution.  The program preserves
this integer event count together with its rational interval and sign
certificates.  When the four oriented edges are joined, endpoint half-events at
a shared vertex cancel, and the remaining signed crossing total is exactly
the net number of turns.  This is the bridge between the sign variations of
the preceding section and the box winding used throughout the rest of the
paper.

\begin{lemma}[polynomial specialization]
\label{lem:polynomial-winding-specialization}
For $P=f$ and $Q=1$, the quotient-edge construction satisfies
$H_{f/1,e}=f(e)=A_{f,e}+iB_{f,e}$ and hence
\[
 \Wind_4(f;e)=\Ind_0^1\!\left(\frac{A_{f,e}}{B_{f,e}}\right),
 \qquad
 \Wind_4(f;B)=\sum_{e\text{ in }\operatorname{edges}(B)}\Wind_4(f;e).
\]
This is the polynomial specialization of \cref{def:quotient-winding}.
\end{lemma}

\begin{proof}
The coordinate polynomials of the constant denominator $1$ are $1$ and
$0$.  Hence \cref{def:division-free-quotient-edge} gives
\[
 \begin{aligned}
 H_{f/1,e}
   &=f(e)\overline{1}=f(e),\\
 C_{f/1,e}
   &=A_{f,e}\cdot1+B_{f,e}\cdot0=A_{f,e},\\
 D_{f/1,e}
   &=B_{f,e}\cdot1-A_{f,e}\cdot0=B_{f,e}.
 \end{aligned}
\]
Unfolding the edge definition therefore gives
\[
 \Wind_4(f;e)
   =\Wind_4(f/1;e)
   =\Ind_0^1\!\left(\frac{C_{f/1,e}}{D_{f/1,e}}\right)
   =\Ind_0^1\!\left(\frac{A_{f,e}}{B_{f,e}}\right).
\]
Summing this equality over the four entries of
$\operatorname{edges}(B)$ proves the box identity.
\end{proof}

\leanref{Purelib.RationalFunctions.Winding.rationalQuadrupledWinding;
Purelib.RootIsolation.Winding.quadrupledBoxWinding}

Every $\Wind_4$ value is obtained by four finite Euclidean chains, and every
winding-one state stores the exact equality $\Wind_4(f;B)=4$.

\begin{lemma}[edge reversal]
\label{lem:winding-edge-reversal}
For an oriented affine edge $e$, put $e^{\mathrm{rev}}(t):=e(1-t)$.
Whenever the edge winding of a displayed polynomial or division-free
quotient expression $F$ is defined,
\[
             \Wind_4(F;e^{\mathrm{rev}})=-\Wind_4(F;e).
\]
\end{lemma}

\begin{proof}
The coordinate polynomials on $e^{\mathrm{rev}}$ are those on $e$ after the
affine change $t\mapsto1-t$.  Thus traversing $e^{\mathrm{rev}}$ from
$t=0$ to $t=1$ is the same finite Cauchy--Sturm calculation as traversing
$e$ from $t=1$ to $t=0$.  The reversal identity in \cref{thm:sturm} negates
that index.
\end{proof}

\leanref{Purelib.RootIsolation.Sturm.cauchyIndex2_reverse}

\begin{lemma}[box subdivision]
\label{lem:winding-add}
If a vertical or horizontal rational cut divides a proper box $B$ into
$B_1$ and $B_2$, then
\[
       \Wind_4(f;B)=\Wind_4(f;B_1)+\Wind_4(f;B_2).
\]
The corresponding statement holds for a finite rectangular grid.
\end{lemma}

\begin{proof}
In the sum of the two child boundary lists, the internal cut occurs once in
each orientation.  Its two contributions cancel by
\cref{lem:winding-edge-reversal}.  Each outer edge of $B$ is either unchanged
or divided into two consecutive rationally parametrized segments; the
subdivision identity in \cref{thm:sturm} replaces the contributions of those
segments by the contribution of the original edge.  The remaining sum is
therefore exactly $\Wind_4(f;B)$.  Iterating the same finite argument over
the cuts proves the rectangular-grid case.
\end{proof}

\leanref{Purelib.RootIsolation.Winding.quadrupledBoxWinding_vertical_split;
Purelib.RootIsolation.Winding.quadrupledBoxWinding_horizontal_split}

\begin{proposition}[coordinate winding]
\label{prop:coordinate-normalization}
For every displayed $R>0$, let $B=[-R,R]^2$.  For the coordinate polynomial
$Z$, one has $\Wind_4(Z;B)=4$.
\end{proposition}

\begin{proof}
Substitute the four affine edge parametrizations of the symmetric square.
Each index is the variation of a chain of degree at most one; the finite sign
table gives edge contributions $0,2,0,2$, whose sum is four.
\end{proof}

\leanref{Purelib.RootIsolation.Winding.Global.coordinate_symmetric_normalization}

\begin{proposition}[winding parity]
\label{prop:parity}
If a displayed polynomial $f$ carries a boundary-clearance certificate on a
proper rational box $B$, then $\Wind_4(f;B)$ is even.
\end{proposition}

\begin{proof}
The four endpoint parity corrections telescope at the shared vertices of
the oriented boundary, so the box sum is even.
\end{proof}

\leanref{Purelib.RootIsolation.Winding.quadrupledBoxWinding_even_of_boundaryCoprime}

\begin{proposition}[product formula]
\label{prop:product}
If $f$ and $g$ carry quantitative boundary-clearance certificates on a
proper rational box $B$, then
\[
       \Wind_4(fg;B)=\Wind_4(f;B)+\Wind_4(g;B).
\]
In particular, for every nonzero constant $c$,
$\Wind_4(c;B)=0$.
\end{proposition}

\begin{proof}
On one edge write $f=A+iB$ and $g=C+iD$.  Then
\[
 \Ree(fg)=AC-BD,\qquad \Imm(fg)=AD+BC.
\]
The finite local product calculation has one correction determined by the
endpoint signs of $(A,B,C,D)$.  Its generic and coordinate-degenerate
branches are finite sign tables.  Summing around the four sides cancels each
outgoing correction with the incoming correction at the same vertex; the
clearance hypotheses remove simultaneous-zero cases.  What remains is the
displayed identity.  This follows the algebraic product strategy of
Eisermann \cite{Eisermann}, but does not invoke its unrestricted auxiliary
formula.  Perrucci and Roy \cite{PerrucciRoyWinding} explain why that formula
needs additional cases.  Here a finite guard partition resolves simultaneous
coordinate events, the coordinate-degenerate branches are checked
separately, and only then are the endpoint corrections telescoped.
\end{proof}

\leanref{Purelib.RootIsolation.Winding.Product.quadrupledBoxWinding_mul_of_boundaryClearances}

\begin{theorem}[finite quotient winding formula]
\label{thm:rfap-quotient}
Let $P,Q\in\Qi[Z]$, and let $C$ be a proper rational box whose boundary
carries positive clearance certificates for both polynomials.  The generic
division-free construction of \cref{def:quotient-winding} satisfies
\[
       \Wind_4(P/Q;C)=\Wind_4(P;C)-\Wind_4(Q;C).
\]
\end{theorem}

\begin{proof}
On one oriented affine edge write
\[
 P(e)=A+iB,
 \qquad
 Q(e)=C+iD.
\]
The direction polynomial from the generic construction is
\[
 P(e)\overline{Q(e)}=(AC+BD)+i(BC-AD).
\]
Apply the finite product calculation underlying
\cref{prop:product} to the coordinate pairs $(A,B)$ and $(C,-D)$.
It expresses the edge index of their product as the sum of their two edge
indices plus a correction determined entirely by the endpoint signs.
Conjugation changes the sign of the second edge index.  After the four
oriented edges are added, the outgoing correction at each vertex cancels
the incoming correction at that same vertex.  Boundary clearance excludes
the simultaneous-zero cases in every finite sign table.  The remaining
identity is therefore the stated subtraction of windings.  Every step is a
signed Euclidean-remainder calculation over $\Q$.
\end{proof}

\leanref{Purelib.RationalFunctions.Winding.rationalQuadrupledWinding_eq_sub_of_boundaryClearances}

\subsection{Certified homotopies and a finite Rouché lemma}

\begin{definition}[dominance certificate]
\label{def:dominance-certificate}
A \emph{dominance certificate} for displayed polynomial expressions $F,G$
on a rationally parametrized boundary consists of a positive rational margin
$m$ and a finite polynomial-inequality proof of
\[
        \Norm{F(t)}-\Norm{G(t)}\ge m>0
\]
under each rational edge predicate.
\end{definition}

Every later use constructs the margin in
\cref{def:dominance-certificate} directly from finite coefficient bounds;
no unbounded search is hidden in the certificate.  Intuitively, it records
enough room to add $G$ to $F$ without allowing the boundary value to reach
zero.

\begin{lemma}[half-plane winding]
\label{lem:half-plane}
If the four coordinate-polynomial expressions obtained by composing a
displayed polynomial $h$ with the edges of a proper rational box $B$ have
quantitative clearance and satisfy $\Ree h(e(t))>0$ under their rational
parameter predicates, then $\Wind_4(h;B)=0$.
\end{lemma}

\begin{proof}
The imaginary part may cross zero, but the real part is positive at every such crossing.  In the endpoint variation formula the two half-crossings on adjacent edges cancel, and no crossing of the negative real ray occurs.  Equivalently, apply the product formula after multiplying by a fixed rational constant; the four composed edge expressions then carry certificates for the displayed half-plane predicate.
\end{proof}

\leanref{Purelib.Polynomial.Geometry.HalfPlane.quadrupledBoxWinding_eq_zero_of_positive_real}

\begin{lemma}[finite dominance homotopy]
\label{lem:dominance-homotopy}
Let $F,G$ be displayed polynomial expressions on a proper rational-box
boundary, with clearance for $F$.  A dominance certificate in the sense of
\cref{def:dominance-certificate} proves
\[
       \Wind_4(F+G;B)=\Wind_4(F;B).
\]
\end{lemma}

\begin{proof}
On each affine edge form the displayed coordinate polynomial for
$\overline F(F+G)$.  The exact two-square identity and the supplied positive
margin make its real coordinate positive throughout the edge.  Finite Horner
bounds turn the same margin into a clearance certificate for $F+G$.
The half-plane formula and the finite product formula then identify the edge
index of $F+G$ with that of $F$, up to one explicit endpoint correction.
The four corrections telescope around the box, leaving the displayed winding
identity.
\end{proof}

\leanref{Purelib.RootIsolation.Winding.Homotopy.quadrupledBoxWinding_add_of_boundaryNormDominance}

\begin{corollary}[finite Rouché rule]
\label{cor:rouche}
Whenever a rational norm bound proves $\Norm{g(z)}<\Norm{f(z)}$ on the clear
boundary of a proper rational box $B$,
\[
        \Wind_4(f+g;B)=\Wind_4(f;B).
\]
\end{corollary}

\begin{proof}
The supplied strict rational inequality is a dominance certificate.
Lemma~\ref{lem:dominance-homotopy} gives the displayed identity.
\end{proof}

\leanref{Purelib.RootIsolation.Winding.Homotopy.quadrupledBoxWinding_add_of_boundaryNormDominance}

\section{Clearing axis-parallel lines}

This section constructs eliminants and uses them to compute safe vertical and
horizontal cuts.  Their boundary-clearance certificates make rational-box
subdivision available to local extraction.

Write
\[
 f(x+iy)=A_f(x,y)+iB_f(x,y),\qquad A_f,B_f\in\Q[x,y].
\]

A vertical coordinate $x=c$ is usable when an explicit B\'ezout calculation
rules out simultaneous vanishing of the two displayed coordinate
polynomials along that line; a horizontal coordinate is treated analogously.
The two eliminants below are univariate guard polynomials for those
calculations.  A nonzero guard value at $c$ lets us construct the required
B\'ezout and quantitative-clearance certificates using finite arithmetic.
The guards thereby choose safe grid lines.

\subsection{Safe axis-parallel cuts}

\begin{definition}[fraction-free scaled-inverse recursion]
\label{def:fraction-free-scaled-inverse}
Let $M$ be a square matrix over $\Qi[T]$.  Its scaled-inverse data $(s,N)$
are defined by recursion on its dimension.  In dimension zero put $s=1$ and
take the empty matrix for $N$.  If $M$ has positive dimension $n+1$, write
\[
 M=\begin{pmatrix}a&u\\ v&M_{\mathrm{tail}}\end{pmatrix},
 \qquad S=aM_{\mathrm{tail}}-vu,
\]
where $a$ is a scalar, $u$ is a $1\times n$ row, $v$ is an $n\times1$
column, and $M_{\mathrm{tail}}$ and $S$ are $n\times n$ matrices.  If
recursion on the smaller matrix $S$ returns $(s',N')$, return
\[
 s=as',\qquad
 N=\begin{pmatrix}
      s'+uN'v&-auN'\\
      -aN'v&a^2N'
    \end{pmatrix}.
\]
\end{definition}

\begin{lemma}[scaled-inverse identity]
\label{lem:scaled-inverse-identity}
If the recursion of \cref{def:fraction-free-scaled-inverse} returns $(s,N)$
from $M$, then the following identities hold over $\Qi[T]$:
\[
                         NM=sI=MN.
\]
\end{lemma}

\begin{proof}
Induct on the dimension.  The zero-dimensional case is immediate.  In
positive dimension, for $S=aM_{\mathrm{tail}}-vu$, the induction hypothesis
gives the polynomial-matrix identities
\[
                 N'S=s'I=SN'.
\]
Multiplying the two block matrices in
\cref{def:fraction-free-scaled-inverse}, the off-diagonal blocks of $NM$
reduce to
\[
 s'u-uN'(aM_{\mathrm{tail}}-vu)=s'u-us'=0,
 \qquad -a^2N'v+a^2N'v=0,
\]
while its diagonal blocks reduce to $as'$ and $aN'S=as'I$.
Thus $NM=as'I=sI$.  Multiplication in the other order uses $SN'=s'I$ and
gives $MN=as'I$ in the same way.
\end{proof}

\leanref{Purelib.Polynomial.SeparatedVariables.scaledInverseData_coefficientIdentity}

\begin{definition}[separated-variable Schur--Sylvester eliminants]
\label{def:axis-eliminants}
Put
\[
       P_T(U)=f(U),\qquad Q_T(U)=\overline f(2T-U),
\]
where $\overline f$ means coefficientwise conjugation, and let $S_x(T)$ be
their displayed $2d\times2d$ Sylvester coefficient matrix over $\Qi[T]$.
Write $(s_x,N_x)$ for the scaled-inverse data of $S_x$ from
\cref{def:fraction-free-scaled-inverse}.  Define $E_x\in\Q[T]$ by taking
coefficientwise the imaginary part of $s_x$ when the real part of its
leading coefficient is zero, and otherwise taking its coefficientwise real
part.

Apply the same construction to the quarter-turned polynomial
$f^{\circlearrowleft}(Z)=f(iZ)$; write $(s_y,N_y)$ for its scaled-inverse
data and $E_y\in\Q[T]$ for the rational coordinate selected by the same
rule.  Finally put
\[
             D_f=\max\{\deg E_x,\deg E_y\}.
\]
\end{definition}

\begin{lemma}[nonzero Schur--Sylvester eliminants]
\label{lem:eliminant-nonzero}
If $f$ is nonconstant, then $E_x$ and $E_y$ of
\cref{def:axis-eliminants} are nonzero.
\end{lemma}

\begin{proof}
Write $a_d$ for the nonzero leading coefficient of $f$.  In each of the
first $d$ columns of $S_x(T)$, coming from $Q_T$, use $d$ as the displayed
parameter-degree bound; in each of the last $d$ columns, coming from
$P_T=f$, use degree zero.  Reading the finite coefficient lists shows that
the matrix of these column-leading coefficients is upper triangular.  Its
first $d$ diagonal entries are
\[
                  2^d\overline{a_d},
\]
the top parameter coefficient created by repeatedly substituting $2T-U$
in $\overline f$, and its last $d$ diagonal entries are $a_d$.  Both are
nonzero in $\Qi$.  Entries strictly below the diagonal have degree strictly
below their column's displayed bound.

Now consider one Schur step of \cref{def:fraction-free-scaled-inverse}.
Suppose the current matrix has column-degree bounds $\delta_j$ and nonzero
diagonal leading coefficients $\lambda_j$.  The head entry has leading term
$\lambda_0T^{\delta_0}$.  In column $j$ of
\[
                 S=aM_{\mathrm{tail}}-vu,
\]
the first product has degree at most $\delta_0+\delta_{j+1}$ and diagonal
leading coefficient $\lambda_0\lambda_{j+1}$.  In the correcting product
$vu$, the entry from the left column lies strictly below the original
diagonal, so its degree is too small to cancel that leading term.  Thus $S$
again has an upper-triangular leading-coefficient matrix, now with diagonal
entries $\lambda_0\lambda_{j+1}$, all nonzero.

Induction therefore shows that the head entry at every recursive stage is a
nonzero polynomial.  Since the recursion starts with scalar $1$ and replaces
$s'$ by $as'$, its output $s_x$ is a finite product of those nonzero pivots
and is itself nonzero.  Let $\alpha+i\beta$ be the leading coefficient of
the normalized $s_x$.  If $\alpha\ne0$, the coefficientwise real part has
leading coefficient $\alpha$; if $\alpha=0$, then $\beta\ne0$ and the
coefficientwise imaginary part has leading coefficient $\beta$.  The rule in
\cref{def:axis-eliminants} therefore makes $E_x$ nonzero.  Finally,
$f^{\circlearrowleft}(Z)=f(iZ)$ has the same degree and leading coefficient
$a_di^d\ne0$, so the identical argument gives $E_y\ne0$.
\end{proof}

\leanref{Purelib.Polynomial.SeparatedVariables.separatedSylvester_leading;
Purelib.LinearAlgebra.PolynomialMatrix.Leading.scaledInverseData_scalar_nonzero;
Purelib.Polynomial.SeparatedVariables.separatedScalar_nonzero;
RFAP.ConstructiveSylvesterBezout.scalarEliminant_nonzero;
RFAP.ConstructiveSylvesterBezout.horizontalEliminant_nonzero;
Purelib.RootIsolation.Isolation.AxisSafeCuts.separatedEliminationDegree}

\begin{lemma}[specialized axis B\'ezout certificates]
\label{lem:axis-bezout}
For every displayed $c\in\Q$ with $E_x(c)\ne0$, finite calculation returns
$U_c,V_c\in\Q[Y]$ with the polynomial identity
\[
             1=U_c(Y)A_f(c,Y)+V_c(Y)B_f(c,Y).
\]
Likewise, $E_y(c)\ne0$ returns displayed polynomials satisfying the
corresponding polynomial identity for $A_f(X,c)$ and $B_f(X,c)$.
\end{lemma}

\begin{proof}
Coefficientwise real and imaginary projection commute with Horner
evaluation at the rational parameter $c$.  Consequently $E_x(c)$ is either
the real or the imaginary coordinate of $s_x(c)$.  If $s_x(c)$ were zero,
that selected coordinate would be zero as well; hence
\[
                    E_x(c)\ne0\quad\Longrightarrow\quad s_x(c)\ne0.
\]

Use the right-hand identity of \cref{lem:scaled-inverse-identity} and take
the first column after specializing at $c$:
\[
             S_x(c)N_x(c)e_0=s_x(c)e_0.
\]
Split the displayed column $N_x(c)e_0$ into its first and last $d$
entries, say $(\ell_0,\ldots,\ell_{d-1})$ and
$(r_0,\ldots,r_{d-1})$, and form
\[
       L_c(U)=\sum_{j<d}\ell_jU^j,
       \qquad R_c(U)=\sum_{j<d}r_jU^j.
\]
Multiplying the Sylvester rows by $1,U,\ldots,U^{2d-1}$ and summing is a
finite rearrangement of coefficients.  The first $d$ columns reconstruct
$Q_c(U)$ and the last $d$ columns reconstruct $P_c(U)$, so the displayed
matrix identity gives the polynomial identity
\[
              L_c(U)Q_c(U)+R_c(U)P_c(U)=s_x(c)
\]
in $\Qi[U]$.  Division by the now certified nonzero scalar
$s_x(c)$ gives polynomials $\widehat L_c,\widehat R_c$ satisfying
\begin{equation}
\label{eq:normalized-separated-bezout}
              \widehat L_c(U)Q_c(U)+
              \widehat R_c(U)P_c(U)=1.
\end{equation}

Substitute the polynomial $c+iY$ into
\eqref{eq:normalized-separated-bezout}.  Since
$\overline{2c-(c+iY)}=c+iY$, the definitions of the separated pair give
\[
 P_c(c+iY)=f(c+iY),
 \qquad
 Q_c(c+iY)=\overline{f(c+iY)}.
\]
Write $f(c+iY)=A_f(c,Y)+iB_f(c,Y)$.  Taking the real part of the substituted
polynomial identity yields
\[
\begin{aligned}
1={}&\bigl(\Ree\widehat L_c(c+iY)+\Ree\widehat R_c(c+iY)\bigr)A_f(c,Y)\\
 &+\bigl(\Imm\widehat L_c(c+iY)-\Imm\widehat R_c(c+iY)\bigr)B_f(c,Y).
\end{aligned}
\]
The two parenthesized expressions are displayed polynomials
$U_c,V_c\in\Q[Y]$, proving the vertical claim.

For the horizontal claim apply the same construction to
$f^{\circlearrowleft}(w)=f(iw)$.  A horizontal point $t+ic$ is obtained from
the displayed parameter $w=c-it$ because $iw=t+ic$.  The vertical identity
for the quarter-turned polynomial therefore becomes the asserted identity
for $A_f(X,c)$ and $B_f(X,c)$ after the same real-coordinate calculation.
\end{proof}

\leanref{RFAP.ConstructiveSylvesterBezout.scalar_ne_zero_of_eliminant_eval_ne_zero;
RFAP.ConstructiveSylvesterBezout.verticalCoprimeCertificate;
RFAP.ConstructiveSylvesterBezout.horizontalCoprimeCertificate}

\begin{lemma}[bounded safe-cut search]
\label{lem:safe-cut}
Let $[a,b]$ be a nondegenerate rational $x$-interval and let $[r,s]$ be a
bounded rational $y$-interval.  Apply \cref{lem:safe-central-third-cut} to
$H=E_x$ with the degree bound $D_f$ of \cref{def:axis-eliminants}.
The resulting finite scan returns $c$ in the open central third of $[a,b]$
with $E_x(c)\ne0$ and computes a positive rational $m_c$ such that
\[
       A_f(c,t)^2+B_f(c,t)^2\ge m_c
       \qquad(r\le t\le s).
\]
Thus the full vertical line segment $c+it$ carries a quantitative boundary
clearance certificate.  Interchanging $x$ and $y$ gives the analogous safe
horizontal cut.
\end{lemma}

\begin{proof}
By \cref{lem:eliminant-nonzero}, $E_x$ is nonzero and
$\deg E_x\le D_f$.  The bounded search in
\cref{lem:safe-central-third-cut} therefore returns the stated $c$ after at most
$D_f+1$ exact evaluations.

Apply \cref{lem:axis-bezout} to $E_x(c)\ne0$ to obtain the displayed
polynomials $U_c,V_c$ and their Bézout identity.
Apply the finite Bézout-clearance calculation of
\cref{lem:positive-lower} to this identity on $[r,s]$.
It returns $m_c>0$ with
$A_f(c,t)^2+B_f(c,t)^2\ge m_c$ for every displayed rational $t$ satisfying
$r\le t\le s$.  This is exactly $m_c\le\Norm{f(c+it)}$.  The calculation
for $E_y$ is identical after interchanging the coordinates.
\end{proof}

\leanref{Purelib.RootIsolation.Isolation.AxisSafeCuts.separatedAxisEliminationCertificate;
Purelib.RootIsolation.Isolation.AxisSafeCuts.verticalSafeCoordinate_spec;
Purelib.RootIsolation.Isolation.AxisSafeCuts.horizontalSafeCoordinate_spec;
Purelib.RootIsolation.Isolation.AxisSafeCuts.verticalSafeCoordinate_contraction;
Purelib.RootIsolation.Isolation.AxisSafeCuts.horizontalSafeCoordinate_contraction;
Purelib.RootIsolation.Isolation.AxisSafeCuts.buildClearFourGrid}

The inequalities in \cref{lem:central-third-contraction} make both child widths
strictly less than $2/3$ of the parent width.  This is the contraction used
by every refinement routine.

\section{The degree winding on a large square}
\label{sec:global-winding}

This section computes the global winding budget on the coefficient isolation
square from displayed coefficient bounds.

Let
\[
       f(Z)=a_dZ^d+\sum_{k=0}^{d-1}a_kZ^k,
       \qquad a_d\ne0.
\]
Recall from \Defref{def:coefficient-isolation-square} that
\[
              C_f=d\sum_{k=0}^{d-1}2^k\Norm{a_k}
\]
is the displayed weighted norm bound for the nonleading coefficients.  Take
the coefficient radius $R=R_f$ from the same definition.  It satisfies
\begin{equation}
\label{eq:large-R}
             1\le R,\qquad C_f<\Norm{a_d}R\le\Norm{a_d}R^2.
\end{equation}
These are direct rational comparisons from the displayed formula
\eqref{eq:isolation-radius}; no search over a geometric object is involved.
The square is called \emph{global} only because every output track uses it as
a common containing box.  It remains finite endpoint data, and the displayed
inequalities are the leading-term margin used below.

\begin{definition}[leading monomial, lower tail, and dominance margin]
\label{def:leading-tail-margin}
For the displayed polynomial $f$ and coefficient radius $R$, put
\[
 L_f(Z)=a_dZ^d,
 \qquad
 T_f(Z)=\sum_{k=0}^{d-1}a_kZ^k,
 \qquad
 f=L_f+T_f.
\]
The associated \emph{leading--tail dominance margin} is the displayed
rational
\[
 M_f=\bigl(\Norm{a_d}R^2-C_f\bigr)(R^2)^{d-1}.
\]
\end{definition}

\begin{lemma}[leading--tail annular margin]
\label{lem:leading-tail-annular-margin}
The margin $M_f$ is positive.  Every displayed $q\in\Qi$ satisfying
\[
                   R^2\le\Norm q\le2R^2
\]
also satisfies
\[
       M_f\le \Norm{L_f(q)}-\Norm{T_f(q)}.
\]
\end{lemma}

\begin{proof}
Positivity follows from $C_f<\Norm{a_d}R^2$ in \eqref{eq:large-R} and
$R>0$.  For the annular estimate, \eqref{eq:finite-sum}, the upper bound on
$\Norm q$, and $R\ge1$ give
\[
\begin{aligned}
 \Norm{T_f(q)}
 &\le d\sum_{k<d}\Norm{a_k}(2R^2)^k\\
 &=d\sum_{k<d}2^k\Norm{a_k}R^{2k}\\
 &\le C_fR^{2d-2}.
\end{aligned}
\]
The lower bound on $\Norm q$ gives
\[
       \Norm{L_f(q)}=\Norm{a_d}\Norm q^d
       \ge\Norm{a_d}R^{2d}.
\]
Subtracting the two rational bounds yields
\[
 \Norm{L_f(q)}-\Norm{T_f(q)}
 \ge\bigl(\Norm{a_d}R^2-C_f\bigr)R^{2d-2}=M_f.
\]
\end{proof}

\leanref{Purelib.RootIsolation.Winding.Global.leadingMonomial;
Purelib.RootIsolation.Winding.Global.lowerTermsPolynomial;
Purelib.RootIsolation.Winding.Global.leadingMonomial_add_lowerTerms_eval;
Purelib.RootIsolation.Winding.Global.largeSquareDominanceMargin;
Purelib.RootIsolation.Winding.Global.largeSquareDominanceMargin_pos;
Purelib.RootIsolation.Winding.Global.largeSquare_radial_dominance_margin}

Let $K_R=(-R,R,-R,R)$.  Its boundary satisfies
$R^2\le\Norm z\le2R^2$, so
Lemma~\ref{lem:leading-tail-annular-margin} gives
\[
       0<M_f\le\Norm{L_f(z)}-\Norm{T_f(z)}.
\]

The meaning of this estimate is simple: on the outer square the leading
monomial is larger than the sum of all lower terms, so the finite Rouché rule
allows the tail to be discarded without changing winding.  Multiplication by
$Z$ contributes one turn, and the product rule therefore makes $Z^d$
contribute $d$.  This is the sole source of the global count used later.

\begin{theorem}[global degree winding]
\label{thm:global-wind}
The boundary of $K_R$ is clear and
\[
             \Wind_4(f;K_R)=4d,\qquad \Wind_f(K_R)=d.
\]
\end{theorem}

\begin{proof}
The dominance inequality of Lemma~\ref{lem:leading-tail-annular-margin} and
\cref{cor:rouche} give
\[
                 \Wind_4(f;K_R)=\Wind_4(L_f;K_R).
\]
A nonzero constant has winding zero, $Z$ has winding four by
\cref{prop:coordinate-normalization}, and repeated use of
\cref{prop:product} gives $\Wind_4(L_f;K_R)=4d$.  The conventional value
$\Wind_f(K_R)=d$ follows from
\cref{def:rational-function-winding}.
\end{proof}

\leanref{Purelib.RootIsolation.Winding.Global.quadrupledBoxWinding_largeSquare_eq_degree}

\section{Local extraction from a nonzero-winding box}
\label{sec:local-extraction}

This section performs local extraction: from a boundary-clear box with
nonzero box winding, it extracts one winding-one root state and retains a
boundary-clear subdivision carrying all the remaining winding.  Applied to
a winding-one state, the same construction is the refinement step later
packaged as a RAT.

More precisely, put $w=\Wind_4(f;B)\ne0$.  For any displayed request
$\eps>0$, the algorithm returns four safe cut lines and their full clear
$3\times3$ subdivision of $B$.  Its central box $D$ is strictly inside $B$
and satisfies
\[
                 \Wind_4(f;D)=4,
                 \qquad \diamq(D)\le\eps^2.
\]
If $\mathcal R$ is the list of the eight noncentral boxes,
\cref{lem:winding-add} gives
\[
             w=4+\sum_{E\in\mathcal R}\Wind_4(f;E).
\]
Thus only the selected box is normalized to winding four.  If $w>4$, total
winding $w-4$ remains in the complementary boxes; if $w=4$, their total is
zero.  Before nonnegativity is proved, $w$ may also be negative; the same
equation then gives the complementary total $w-4$.  The fact that each
extraction removes exactly four from the retained remainder is the step
iterated in the counting argument of \cref{sec:exact-isolation}.

The construction has five stages.

\begin{enumerate}[label=(\roman*)]
\item Safely subdivide $B$ until a nonzero-winding child $C$ is small and lies
away from the parent's boundary.
\item Use the small diameter and nonzero winding to prove that the rational
anchor $a_C$ has a small residual $\Norm{f(a_C)}$.
\item Use the squarefree Bézout certificate to prove that $f'(a_C)$ is not
small, and hence compute a rational Newton center.
\item Use the explicit Taylor bound to show that the affine part about that
center dominates the nonlinear remainder on a small surrounding band.
\item Choose four eliminant-safe lines in that band.  Their central rectangle
is the returned box $D$; the full grid records both $D$ and its eight
complementary boxes.
\end{enumerate}

The constants introduced below merely allocate enough rational margin for
these five steps to hold simultaneously.

\subsection{Finite local bounds and the protected boundary}

\begin{definition}[local finite Horner bounds]
\label{def:local-horner-bounds}
Fix a current rational box $B$ and write
\[
 R_B=2\max\{x_-^2,x_+^2,y_-^2,y_+^2\}.
\]
For the displayed high-to-low coefficient list
$[a_n,\ldots,a_0]$ of $f$, start with
$(M_0,S_0,T_0)=(0,0,0)$ and, for $0\le k\le n$, compute
\[
\begin{aligned}
 M_{k+1}&=2\bigl(M_kR_B+\Norm{a_{n-k}}\bigr),\\
 S_{k+1}&=2\bigl(S_kR_B+M_k\bigr),\\
 T_{k+1}&=2\bigl(S_k+R_BT_k\bigr).
\end{aligned}
\]
Define the rationals
\[
             L_1(B)=\max\{1,S_{n+1}\},\qquad
             L_2(B)=\max\{1,T_{n+1}\}.
\]
\end{definition}

\begin{lemma}[local Horner estimates]
\label{lem:local-horner-estimates}
The quantities of \cref{def:local-horner-bounds} are positive, since
\[
                         L_1(B),L_2(B)\ge1.
\]
For every pair of displayed Gaussian rationals $z,w$ satisfying the
membership predicate of $B$, they also satisfy
\begin{align}
 \Norm{f(z)-f(w)}
   &\le L_1(B)\Norm{z-w},
   \label{eq:local-l1}\\
 \Norm{f(z)-f(w)-f'(w)(z-w)}
   &\le L_2(B)\Norm{z-w}^2.
   \label{eq:local-l2}
\end{align}
\end{lemma}

\begin{proof}
Run Horner evaluation simultaneously on the displayed coefficient prefixes
at $z$ and $w$.  At stage $k$, $M_k$ bounds the two partial values, $S_k$
bounds their divided difference, and $T_k$ bounds the remainder after
removing its constant and linear terms.  Multiplication by the next argument
uses $\Norm z,\Norm w\le R_B$; adjoining the next coefficient uses the finite
two-term inequality $\Norm{x+y}\le2(\Norm x+\Norm y)$.  These are exactly the
three updates defining $M_{k+1},S_{k+1},T_{k+1}$.  Induction over the finite
coefficient list therefore gives the two displayed estimates with
$S_{n+1}$ and $T_{n+1}$.  Replacing them by their maxima with $1$ preserves
the inequalities and gives $L_1(B),L_2(B)>0$.
\end{proof}

Thus $L_1(B)$ is the squared-distance Lipschitz coefficient and $L_2(B)$
is the quadratic Taylor-remainder coefficient.  They are obtained by one
finite fold over the coefficient list; neither is a maximum or supremum over
the points satisfying the box predicate.  In the notation of
\eqref{eq:lip1} and \eqref{eq:lip2}, they are $L_1(f,B)$ and $L_2(f,B)$.

\leanref{Purelib.Polynomial.Bounds.HornerNormBounds.step;
Purelib.Polynomial.Bounds.denseHornerNormBounds;
Purelib.Polynomial.Bounds.denseLipschitzBound;
Purelib.Polynomial.Bounds.HornerTaylorNormBounds.step;
Purelib.Polynomial.Bounds.denseHornerTaylorNormBounds;
Purelib.Polynomial.Bounds.denseTaylorRemainderBound;
Purelib.Polynomial.Bounds.normSq_eval_sub_eval_le_denseLipschitzBound;
Purelib.Polynomial.Bounds.normSq_dense_taylor_remainder_le_denseTaylorRemainderBound}

We continue to write $\tau(B),\eta(B)$ for the quantities in
\eqref{eq:bez-margins} computed on $B$, and $\mu(B)>0$ for the bound in its
four-edge clearance certificate.

\begin{definition}[protected boundary radius]
\label{def:protected-boundary-radius}
Define
\[
 \rho(B)=\min\left\{\frac1{512},
       \frac{\mu(B)}{2L_1(B)},
       \frac{\operatorname{width}_x(B)}{16},
       \frac{\operatorname{width}_y(B)}{16}\right\}
\]
from the displayed clearance, Horner bound, and endpoint data of $B$.
\end{definition}

\begin{lemma}[positivity of the protected radius]
\label{lem:protected-radius-positive}
If $B$ is proper and carries the stated clearance certificate, then
$\rho(B)>0$.
\end{lemma}

\begin{proof}
Properness of $B$ makes both coordinate widths positive.  The clearance
certificate gives $\mu(B)>0$, while
\cref{lem:local-horner-estimates} gives $L_1(B)>0$.  Hence all four entries
in the finite minimum of \cref{def:protected-boundary-radius} are positive.
\end{proof}

\leanref{Purelib.Polynomial.Geometry.ProtectedBoundary.protectedBoundaryRadius;
Purelib.Polynomial.Geometry.ProtectedBoundary.protectedBoundaryRadius_pos}

\begin{lemma}[protected boundary strip]
\label{lem:protected-strip}
Let $B$ be proper and carry the stated clearance certificate.  If a rational
subbox $D\subseteq B$ has $\diamq(D)\le\rho(B)^2$ and meets the coordinate
boundary strip of width $4\rho(B)$, then $\Wind_4(f;D)=0$.
\end{lemma}

\begin{proof}
Put $\rho=\rho(B)$, $L=L_1(B)$, and $\mu=\mu(B)$.  Properness of $B$ makes
both coordinate widths positive; $L$, $\mu$, and $\rho$ are positive by
their finite certificates and \cref{lem:protected-radius-positive}.

Meeting the strip supplies a displayed $z\in D$ and one of four rational
side inequalities.  Consider the left-side case
\[
                  \Ree z\le x_-^B+4\rho;
\]
the other three cases are the same calculation with the relevant coordinate
reversed.  Since $D\subseteq B$, also $x_-^B\le\Ree z$.  Project $z$ to the
displayed parent-boundary point
\[
                         b=x_-^B+i\Imm z.
\]
Then $b$ satisfies the left-edge parametrization predicate of $B$ and
\[
             \Norm{b-z}=(\Ree z-x_-^B)^2\le16\rho^2.
\]
Let $q=a_D$ be the fixed lower-left anchor of $D$.  Membership of $z$ and
$q$ in $D$ gives
\[
                     \Norm{z-q}\le\diamq(D)\le\rho^2.
\]
Using $b-q=(b-z)+(z-q)$ and the finite two-term quadrance inequality,
\begin{align*}
 \Norm{b-q}
 &\le2\bigl(\Norm{b-z}+\Norm{z-q}\bigr)\\
 &\le2(16\rho^2+\rho^2)=34\rho^2\le64\rho^2.
\end{align*}
The point $b$ satisfies the boundary-membership predicate of $B$.  Also
$q\in D$ and $D\subseteq B$, so $q$ satisfies the membership predicate of
$B$.  We may therefore apply \eqref{eq:local-l1} to the displayed pair
$(b,q)$.  Substituting the preceding distance estimate gives
\[
 \Norm{f(b)-f(q)}\le L\Norm{b-q}\le L(64\rho^2)=64L\rho^2,
\]
and hence
\begin{equation}
\label{eq:protected-boundary-variation}
                 \Norm{f(b)-f(q)}\le64L\rho^2.
\end{equation}

Set
\[
                 A=L\diamq(D),\qquad V=64L\rho^2.
\]
The definition of $\rho$ gives
\[
 A\le L\rho^2,\qquad \rho\le\frac1{512},\qquad
 2L\rho\le\mu.
\]
Since $130/512<2$ and $L,\rho>0$, these rational inequalities yield
\begin{equation}
\label{eq:protected-boundary-budget}
       2(A+V)\le130L\rho^2<2L\rho\le\mu.
\end{equation}

The boundary-clearance certificate at $b$, followed by the two-term
quadrance inequality and
\eqref{eq:protected-boundary-variation}, gives
\[
 \mu\le\Norm{f(b)}
   \le2\bigl(\Norm{f(q)}+\Norm{f(b)-f(q)}\bigr)
   \le2\bigl(\Norm{f(q)}+V\bigr).
\]
If $\Norm{f(q)}\le A$, this contradicts
\eqref{eq:protected-boundary-budget}, because the preceding display would
then give the impossible chain
\[
                  \mu\le2(A+V)<\mu.
\]
Consequently
\begin{equation}
\label{eq:protected-anchor-dominates}
                         A<\Norm{f(q)}.
\end{equation}

Finally, let $x$ be any displayed point on one of the four edge
parametrizations of $D$.  Both $x$ and $q$ satisfy the membership predicate
of $D$, and hence, because $D\subseteq B$, that of $B$.  We may apply
\eqref{eq:local-l1} once more; the box-diagonal certificate then gives
\[
 \Norm{f(x)-f(q)}\le L\Norm{x-q}\le L\diamq(D)=A<\Norm{f(q)}.
\]
Moreover $A\ge0$, so \eqref{eq:protected-anchor-dominates} implies
$f(q)\ne0$.  In the exact polynomial decomposition
\[
                    f=f(q)+\bigl(f-f(q)\bigr),
\]
the nonzero constant polynomial $f(q)$ therefore strictly dominates the
second summand on the complete rational boundary of $D$.
\Cref{lem:dominance-homotopy} identifies $\Wind_4(f;D)$ with the winding of
that constant, which is zero by \cref{prop:product}.  Hence
$\Wind_4(f;D)=0$, as required.
\end{proof}

\leanref{Purelib.Polynomial.Geometry.ProtectedBoundary.nearbyBoundaryPoint;
Purelib.Polynomial.Geometry.ProtectedBoundary.protectedBoundaryRadius_dominates;
Purelib.Polynomial.Geometry.ProtectedBoundary.quadrupledBoxWinding_eq_zero_of_boundaryStripWitness;
Purelib.Polynomial.Geometry.ProtectedBoundary.strictly_inside_protectedBoundary;
Purelib.RootIsolation.Isolation.LocalExtraction.quadrupledBoxWinding_eq_zero_of_anchorResidualBound_lt}

\subsection{Newton center and extraction}

\begin{lemma}[small nonzero box gives a rational Newton center]
\label{lem:newton-center}
Let $C\subseteq B$ be a rational box with $\Wind_4(f;C)\ne0$, fixed lower-left
anchor $a=a_C\in\Qi$, and squared diagonal $\delta^2$.  Then
\[
     \Norm{f(a)}\le L_1(B)\delta^2.
\]
If in addition $L_1(B)\delta^2\le\tau(B)$, then
$\Norm{f'(a)}\ge\eta(B)$, and consequently
\[
       c=a-\frac{f(a)}{f'(a)}\in\Qi,
       \qquad
       \Norm{c-a}\le\frac{L_1(B)\delta^2}{\eta(B)}.
\]
\end{lemma}

\begin{proof}
Every displayed boundary point $z$ of $C$ lies with $a$ in $B$, and
$\Norm{z-a}\le\delta^2$.  Hence \cref{lem:local-horner-estimates} gives
\[
       \Norm{f(z)-f(a)}\le L_1(B)\delta^2.
\]
If the claimed residual bound failed, the nonzero constant $f(a)$ would
strictly dominate $f-f(a)$ on all four edges of $C$.  The finite Rouch\'e rule
would then give $\Wind_4(f;C)=0$, contrary to the supplied certificate.
Therefore $\Norm{f(a)}\le L_1(B)\delta^2$.

Under the additional hypothesis this is at most $\tau(B)$, so
\cref{lem:derivative-margin} gives
$\Norm{f'(a)}\ge\eta(B)>0$.  In particular $f'(a)\ne0$, and field arithmetic
constructs the displayed $c\in\Qi$.  Finally,
\[
 \Norm{c-a}
 =\Norm{\frac{f(a)}{f'(a)}}
 =\frac{\Norm{f(a)}}{\Norm{f'(a)}}
 \le\frac{L_1(B)\delta^2}{\eta(B)},
\]
using multiplicativity of the Gaussian-rational norm and the two preceding
bounds.
\end{proof}

\leanref{Purelib.RootIsolation.Isolation.LocalExtraction.anchor_residual_le_of_quadrupledBoxWinding_ne_zero;
Purelib.RootIsolation.Isolation.LocalExtraction.computeNewtonCenterOfRawWinding}

\begin{definition}[local extraction scales]
\label{def:local-extraction-scales}
For a proper cleared box $B$ and a request $\eps>0$, define
\[
 r_{\mathrm{ext}}(B,\eps)=
 \min\left\{\frac{\rho(B)}2,1,
       \frac{\eta(B)}{648L_2(B)},\frac\eps8\right\},
\]
and, abbreviating this radius to $r$, define the source-diagonal threshold
\[
 \theta_{\mathrm{src}}(B,\eps)=
 \min\left\{\rho(B)^2,
       \frac{\tau(B)}{L_1(B)},
       \frac{\eta(B)r^2}{L_1(B)}\right\}.
\]
\end{definition}

\begin{lemma}[extraction-scale budgets]
\label{lem:local-extraction-scale-budgets}
Let $B$ and $\eps>0$ satisfy the hypotheses of
\cref{def:local-extraction-scales}.  Put $r=r_{\mathrm{ext}}(B,\eps)$ and
$\theta=\theta_{\mathrm{src}}(B,\eps)$.  Then $r,\theta>0$ and
\[
\begin{gathered}
 2r\le\rho(B),\qquad r\le1,\qquad
 648L_2(B)r\le\eta(B),\qquad 8r\le\eps,\\
 \theta\le\rho(B)^2,\qquad
 L_1(B)\theta\le\tau(B),\qquad
 L_1(B)\theta\le\eta(B)r^2.
\end{gathered}
\]
\end{lemma}

\begin{proof}
The quantities $\rho(B),\tau(B),\eta(B),L_1(B),L_2(B)$ and $\eps$ are
positive, so every entry in both displayed finite minima is positive.  The
four bounds on $r$ are the four minimum projections, with positive
denominators cleared.  The three bounds on $\theta$ follow identically from
its three minimum projections.
\end{proof}

\leanref{Purelib.RootIsolation.Newton.Extraction.extractionRadius;
Purelib.RootIsolation.Newton.Extraction.extractionRadius_pos;
Purelib.RootIsolation.Newton.Extraction.extractionRadius_twice_le_protected;
Purelib.RootIsolation.Newton.Extraction.extractionRadius_le_one;
Purelib.RootIsolation.Newton.Extraction.extractionRadius_le_eighth;
Purelib.RootIsolation.Newton.Extraction.extractionRadius_taylor;
Purelib.RootIsolation.Newton.Extraction.sourceDiagonalThreshold;
Purelib.RootIsolation.Newton.Extraction.sourceDiagonalThreshold_pos;
Purelib.RootIsolation.Newton.Extraction.sourceDiagonalThreshold_le_protected;
Purelib.RootIsolation.Newton.Extraction.sourceDiagonalThreshold_le_derivative;
Purelib.RootIsolation.Newton.Extraction.sourceDiagonalThreshold_le_displacement}

\begin{definition}[Newton band]
\label{def:newton-band}
Let $c\in\Qi$ and $r\in\Q_{>0}$.  A \emph{Newton band} of centre $c$ and
radius $r$ is a rational box $D$ satisfying
\[
\begin{aligned}
 \Ree c-2r&\le x_-^D\le\Ree c-r,
 &\Ree c+r&\le x_+^D\le\Ree c+2r,\\
 \Imm c-2r&\le y_-^D\le\Imm c-r,
 &\Imm c+r&\le y_+^D\le\Imm c+2r.
\end{aligned}
\]
Thus $c$ lies strictly inside $D$, both side lengths are at most $4r$, and
$\diamq(D)\le32r^2$.  An \emph{eliminant-safe Newton band} additionally
carries the four full-line clearance certificates obtained by selecting its
sides with the bounded safe-cut searches.
\end{definition}

\leanref{Purelib.RootIsolation.Newton.Band.NewtonBandBox;
Purelib.RootIsolation.Newton.Band.NewtonBandBox.center_strictly_inside;
Purelib.RootIsolation.Newton.Band.NewtonBandBox.diagonalSq_le}

\begin{theorem}[local extraction]
\label{thm:extract-one}
Fix the displayed squarefree and axis-elimination certificates.  Let $B$ be
a proper rational box carrying certificates for ambient containment,
boundary clearance, and $\Wind_4(f;B)\ne0$.  There is a terminating algorithm
\[
 \begin{aligned}
 &\operatorname{extractOne}(f,B,\mathcal W,\eps),\\
 &\mathcal W:\Wind_4(f;B)\ne0,\qquad \eps\in\Q_{>0}.
 \end{aligned}
\]
which returns rational coordinates
\[
       x_-<x_+,
       \qquad y_-<y_+,
\]
strictly inside $B$ such that:
\begin{enumerate}[label=(\alph*)]
\item the four full axis-parallel cut lines through these coordinates are safe within $B$;
\item the central box $D=(x_-,x_+,y_-,y_+)$ has $\Wind_4(f;D)=4$;
\item splitting $B$ along the four lines gives a finite $3\times3$ grid with clear boundaries, the central member being $D$;
\item $\diamq(D)\le\eps^2$.
\end{enumerate}
\end{theorem}

\begin{proof}
Put $\rho=\rho(B)$, $r=r_{\mathrm{ext}}(B,\eps)$, and
$\theta=\theta_{\mathrm{src}}(B,\eps)$.  Their positivity and all required
rational budgets are given by \cref{lem:local-extraction-scale-budgets}.
At each subdivision round, apply the safe-cut construction of
\cref{lem:safe-central-third-cut} to the current $x$-coordinate interval with
$H=E_x$, and to the $y$-coordinate interval with $H=E_y$.  Pairing the two
$x$-interval children with the two $y$-interval children gives four rational
child boxes, each of whose side widths is less than $2/3$ of the corresponding
parent width.  Run this finite four-child recursion to its
precomputed geometric depth and retain the first nonzero-winding child $C$ with
$\diamq(C)\le\theta$.  By \cref{lem:protected-strip}, $C$ lies in the
protected core of $B$.

Let $a=a_C$ be its anchor and let
\[
             c=a-\frac{f(a)}{f'(a)}
\]
be the rational Newton center from \cref{lem:newton-center}.  The definition
of $\theta$ and \cref{lem:local-extraction-scale-budgets,lem:newton-center}
give $\Norm{c-a}\le r^2$, while the protected-core inequalities and
$2r\le\rho$ keep the square of half-side $2r$ about $c$ strictly inside $B$.

For $z$ in that square,
\[
 \Norm{z-c}\le8r^2,
 \qquad
 \Norm{z-a}\le2\bigl(\Norm{z-c}+\Norm{c-a}\bigr)\le18r^2.
\]
Put
\[
       \ell(z)=f(a)+f'(a)(z-a)=f'(a)(z-c).
\]
On any candidate rectangular boundary surrounding $c$ whose coordinate
distance from $c$ is at least $r$ and at most $2r$,
\[
       \Norm{\ell(z)}=\Norm{f'(a)}\Norm{z-c}\ge\eta(B) r^2.
\]
The exact Taylor identity and the two corresponding radius budgets in
\cref{lem:local-extraction-scale-budgets} give the finite dominance estimate
\[
 \Norm{f(z)-\ell(z)}
    \le324L_2(B)r^4
    \le\frac12\eta(B)r^2.
\]
Thus, for every such rectangular boundary enclosing $c$, the quantitative
Rouché rule and \cref{prop:product,prop:coordinate-normalization} give
$\Wind_4(f;E)=\Wind_4(\ell;E)=4$ for its enclosed box $E$.

Finally run the bounded eliminant search of \cref{lem:safe-cut} in each of
the four rational bands
\[
 [\Ree c-2r,\Ree c-r],\quad[\Ree c+r,\Ree c+2r],
 \quad[\Imm c-2r,\Imm c-r],\quad[\Imm c+r,\Imm c+2r].
\]
The selected full cut lines are safe, and their central rectangle $D$
surrounds $c$ while staying in the affine-dominance square.  Hence
$\Wind_4(f;D)=4$, and the four lines produce the certified $3\times3$ grid.
Both side lengths are at most $4r$, so
$\diamq(D)\le32r^2\le\eps^2$.
\end{proof}

\leanref{Purelib.RootIsolation.Newton.Extraction.extract}

\section{Finite derivative--Taylor track separation}
\label{sec:track-separation}

This section turns coefficient bounds and squarefree certificates into
separation certificates for RATs.  Separation is the input needed later for
finite maximality.

The local extractor can be called inside any nonzero-winding box.  To
bound the length of every finite family of separated outputs in a bounded
square, we first compute a scale on which the derivative term cannot be
cancelled by both a quadratic Taylor error and two sufficiently small
displayed residuals.  This is the same coefficient machinery already used
by local extraction.

Every observation considered in this section is a RAT, so its support record
in \Defref{def:finite-observation-separation} is a proper rational box.
Later refinements move the relevant endpoints inward, preserving the same
displayed inequality and gap.  This witnessed finite relation is what
\emph{distinct} means below.

The underlying idea is uniqueness of a good linear approximation.  If two
separated tracks both fitted inside one sufficiently small box, their
rational anchors would be distinct by a witnessed gap, yet both polynomial
values could be made arbitrarily small.  Taylor's identity would then say
that the non-small derivative times the anchor difference is simultaneously
large enough to detect that gap and small enough to be absorbed by the two
residuals and the quadratic error.  The explicit scale $\upsilon_f$ is where
those two conclusions become incompatible.

\begin{lemma}[finite uniqueness and residual budgets]
\label{lem:separation-budgets}
On the coefficient isolation square $K_R$, put
\[
 \upsilon_f=
 \frac{\eta(K_R)}{8L_2(f,K_R)}>0.
\]
For every displayed coordinate gap $h>0$, the rational quantities
\[
 \beta_h=\min\left\{\tau(K_R),
        \frac{\eta(K_R)h^2}{32}\right\}>0,
 \qquad
 r_h=\min\left\{1,\frac{\beta_h}{L_1(f,K_R)}\right\}>0
\]
are computable.  Calling a winding-one RAT with request $r_h$ returns an
anchor $u$ satisfying
\[
                         \Norm{f(u)}\le\beta_h.
\]
\end{lemma}

\begin{proof}
Positivity follows by finite rational comparison from the positive
squarefree derivative margin and the positive Horner bounds.  The returned
box has squared diagonal at most $r_h^2\le r_h$.  Applying the anchor
residual estimate and the definition of $r_h$ gives
$L_1(f,K_R)r_h^2\le\beta_h$.
\end{proof}

\leanref{Purelib.RootIsolation.Isolation.TrackSeparation.uniquenessDiagonalThreshold;
Purelib.RootIsolation.Isolation.TrackSeparation.uniquenessDiagonalThreshold_pos;
Purelib.RootIsolation.Isolation.TrackSeparation.pairResidualThreshold;
Purelib.RootIsolation.Isolation.TrackSeparation.pairResidualThreshold_pos;
Purelib.RootIsolation.Isolation.TrackSeparation.residualRefinementRadius;
Purelib.RootIsolation.Isolation.TrackSeparation.refined_anchor_residual_le}

\begin{theorem}[root-track separation]
\label{thm:root-separation}
The positive rational $\upsilon_f$ of
\cref{lem:separation-budgets} has the following finite property.  Two
separated winding-one RATs cannot both have current boxes contained in one
rational box $C\subseteq K_R$ with
$\diamq(C)<\upsilon_f$.
\end{theorem}

\begin{proof}
Suppose two separated tracks have current boxes in such a $C$.  Their
separation certificate gives a positive rational coordinate gap $h$.
Request the residual radius $r_h$ of
\cref{lem:separation-budgets} from both tracks, and call the returned anchors
$u,v$.  Put $t=v-u$.  Separation and common-box membership give
\[
             0<h^2\le\Norm t\le\diamq(C)<\upsilon_f.
\]
The two residual requests give
\[
                       \Norm{f(u)},\Norm{f(v)}\le\beta_h,
\]
while the derivative margin gives
\[
             \eta(K_R)\le\Norm{f'(u)}.
\]
Write the exact finite Taylor identity as
\[
             f(v)-f(u)=f'(u)t+E,
 \qquad
             \Norm E\le L_2(f,K_R)\Norm t^2.
\]
The inequality $\Norm t<\upsilon_f$ bounds $8\Norm E$ by
$\eta(K_R)\Norm t$.  The definition of $\beta_h$, the two residual bounds,
and $h^2\le\Norm t$ likewise bound
$8\Norm{f(v)-f(u)}$ by $\eta(K_R)\Norm t$.  Two applications of the finite
two-term quadrance inequality would then force
\[
 4\eta(K_R)\Norm t\le2\eta(K_R)\Norm t,
\]
contradicting the displayed positivity of both factors.  Thus the common
box cannot exist.
\end{proof}

\leanref{Purelib.RootIsolation.Isolation.TrackSeparation.separated_tracks_not_supported_in_small_common_box}

\begin{corollary}[finite packing bound]
\label{cor:packing}
There is a computable natural number $P_f$ such that no finite list of pairwise separated winding-one RATs supported in $K_R$ has length greater than $P_f$.
\end{corollary}

\begin{proof}
Choose a rational mesh $s>0$ with $8s^2<\upsilon_f$ and subdivide the
endpoint record for $[-R,R]^2$ into an explicit finite grid list of half-open
cell records of side at most $s$.  Ask every member of a proposed finite list
for a box of coordinate width below $s/4$, and assign it to the cell whose
membership predicate accepts its rational anchor.  If two anchors receive
the same cell, both refined boxes lie in the coordinatewise enlargement of
that cell by $s/4$; this common rational box has squared diagonal below
$\upsilon_f$, contradicting \cref{thm:root-separation}.  Finite pigeonhole
induction therefore bounds the list length by the number of cells; take that
number as $P_f$.
\end{proof}

\leanref{Purelib.RootIsolation.Isolation.Packing.packingBound;
Purelib.RootIsolation.Isolation.Packing.family_card_le_packingBound}

\section{Nonnegativity and exact extraction of \texorpdfstring{$d$}{d} tracks}
\label{sec:exact-isolation}

This section combines the global winding budget, local extraction, and
separation to construct exactly $d$ RATs, proves ownership and finite
maximality, and relates small residuals to local and global refinement.

At this point $\Wind_4$ is an additive integer, and every nonzero-winding box
contains a finitely extractable winding-one subbox.  The separation
theorem forces the remaining sign.

The argument is a finite counting trap.  If a clear box had negative
winding, extracting a winding-one box from any nonzero remainder and removing its
contribution would leave the remainder negative, so the operation could be
repeated any prescribed finite number of times.  Repeating it more than the
finite packing bound would produce too many separated winding-one boxes.
Negative winding is therefore impossible.  Once every local contribution is
nonnegative, the global total $4d$ can be spent only as exactly $d$
contributions of four; this is the accounting step that produces the initial
tracks.

\begin{theorem}[nonnegativity]
\label{thm:nonnegative}
Every clear rational box $B\subseteq K_R$ satisfies
\[
             \Wind_4(f;B)\ge0.
\]
The statement also holds for a finite disjoint list of clear boxes, with
total winding defined by finite summation.
\end{theorem}

\begin{proof}
Assume a finite list of rational-box records has negative total winding.
Choose a member
with nonzero winding and apply \cref{thm:extract-one} with request $1$.
Remove the central winding-one box and replace the old member by the
other eight grid boxes.  Additivity decreases the total by exactly four,
so it remains negative and the operation may be repeated.

Perform this finite loop exactly $P_f+1$ times.  It returns $P_f+1$ pairwise distinct winding-one boxes, each carrying the deterministic refinement routine of \cref{thm:extract-one}.  This contradicts the finite packing bound \cref{cor:packing}.  Therefore the original total could not be negative.
\end{proof}

\leanref{Purelib.RootIsolation.Isolation.GlobalNonnegativity.quadrupledBoxWinding_nonnegative}

\begin{theorem}[exact finite isolation]
\label{thm:isolation}
Starting from $K_R$, a recursion on $d$ returns exactly $d$ pairwise separated clear rational boxes $B_1,\dots,B_d$ with
\[
       \Wind_4(f;B_j)=4.
\]
Every remaining grid box has winding zero.
\end{theorem}

\begin{proof}
The initial total is $4d$ by \cref{thm:global-wind}.  If the current
total is positive, some member has nonzero winding; apply
\cref{thm:extract-one} with request $1$ and replace its parent by the eight
complements.  By additivity the remaining total decreases by four.  By
\cref{thm:nonnegative} it never becomes negative.  After exactly $d$
iterations the remaining total is zero, and nonnegativity forces every
remaining member to have winding zero.
\end{proof}

\leanref{Purelib.RootIsolation.Isolation.ExactIsolation.isolateLargeSquare}

Operationally, a RAT is a certified zoom command: supply its current finite
record and a requested scale, and it returns one smaller certified record.
The bounded search may inspect several candidate boxes internally, but one
request has one observable answer.  The word \emph{track} denotes the ability
to repeat this interaction.  The zero/one/many distinction belongs to the
enclosing isolation algorithm, which returns an exact terminal or a finite
list of tracks.

\begin{lemma}[total strict winding-one refinement]
\label{lem:track-refinement}
Let $B$ be a rational box carrying a certified winding-one state, with the
displayed squarefree and axis-elimination certificates.  For every
$\eps\in\Q_{>0}$ the local
Newton-band extractor returns a rational box $D$ carrying the next certified
state and satisfying
\[
 D\subset B,\qquad \Wind_4(f;D)=4,
 \qquad \diamq(D)\le\eps^2.
\]
The same finite construction, stored as the state transition, can be called
again on $D$.
\end{lemma}

\begin{proof}
Apply \cref{thm:extract-one} to the current state and the requested
$\eps$.  Its returned four line-clearance certificates, $\Wind_4$
calculation, strict-containment inequalities, and diagonal estimate are
exactly the fields of the next finite state.  Reapplying the same definition
therefore gives a total request-driven refinement track.
\end{proof}

\leanref{Purelib.RootIsolation.Isolation.ShrinkingTrack.refineState_accurate}

\paragraph{Output ownership.}
Fix the output of \cref{thm:isolation}.  Its partition has output boxes
$B_1,\ldots,B_d$ and zero-winding remainder cells.  Let $m_{\Pi}>0$ be the
minimum of the finitely many displayed cell-clearance bounds, and define
\[
 \rho_{\rm own}:=min\left\{1,\frac{m_{\Pi}}{8L_1(f,K_f)}\right\}>0.
\]
For a competing RAT $T$ supported in the same coefficient isolation square,
make the single call at this radius and put
\[
                 D_T:=\operatorname{box}(T(\rho_{\rm own})).
\]
An \emph{ownership certificate} for $T$ is an index
$j\in\{1,\ldots,d\}$ and a proof of strict containment
\[
                              D_T\subset B_j.
\]
The deterministic selector scans the finite partition for the anchor of
$D_T$ and proves that the whole box, not merely its anchor, lies strictly
inside the selected cell.  A zero-winding remainder cell cannot be selected,
so the scan returns an output index $j$.  We write
$\operatorname{own}(T)=j$ and call the $j$th returned RAT the \emph{owner} of
$T$.

\begin{proof}[Proof of the polynomial part of \cref{thm:main}]
If $d=1$, exact finite coefficient normalization and division compute the
displayed Gaussian-rational zero and verify it by evaluation.  Suppose
$d\ge2$.  Use \cref{thm:isolation} for the $d$ initial states and
\cref{lem:track-refinement} for every accuracy call.  Let $q_B$ be the fixed
rational anchor of a returned box.  The nonzero-winding argument in
\cref{lem:newton-center} gives
\[
 \Norm{f(q_B)}\le L_1(f,A)\diamq(B).
\]
For an arbitrary displayed $z\in B$, \cref{eq:lip1} and the finite sum
inequality give the uniform, deliberately looser estimate
\[
 \Norm{f(z)}
 \le2\bigl(\Norm{f(z)-f(q_B)}+\Norm{f(q_B)}\bigr)
 \le4L_1(f,A)\diamq(B).
\]
Equation \eqref{eq:projective-box-upper} proves the bound on Hermitian
projective quadrance, and substitution of \eqref{eq:box-projective-exact} gives
the opposite-corner Hermitian-projective-quadrance--residual formula in
\cref{thm:main}.  If $q_{B'}$ is the anchor of a
refining box, both anchors satisfy the old membership predicate, hence
$\Norm{q_{B'}-q_B}\le\diamq(B)$.

The isolation routine leaves a finite safe rectangular partition consisting
of the $d$ output boxes and boxes with $\Wind_4=0$.  Let an additional
winding-one RAT be given.  The positive lower bounds on the finitely many
partition edges, together with the residual estimate just proved, give a
finite refinement after which its box meets no partition edge; hence it lies
in one partition member.  It cannot lie in a zero-winding member: safe cuts
around its smaller winding-one box would make additivity and nonnegativity
give total winding at least four there.  It therefore lies
in an output box $B_j$.

If it were separated from the output RAT in $B_j$, refine the two to disjoint
boxes with a positive rational gap.  Safe-cut searches in that gap extend
them to two disjoint members of a finite grid inside $B_j$.  Nonnegativity
and additivity would give $\Wind_4(f;B_j)\ge8$, contradicting
$\Wind_4(f;B_j)=4$.  Thus every pairwise separated family injects, by a
finite cell assignment, into the $d$ output boxes.  Its length is at most
$d$.
\end{proof}

\leanref{Purelib.RootIsolation.Isolation.Maximality.family_card_le_natDegree;
Purelib.RootIsolation.Isolation.BoxFTA.provePositive}

The main theorem gives the easy direction: a sufficiently small
winding-one box makes the residual at its anchor small.  The next result asks
for a controlled converse.  Starting from a displayed sample with small
residual, can we build a certified nearby track?  The answer must be local:
we also need a protected box and its squarefree derivative margin, since the
bare number $\Norm{f(q)}$ changes when $f$ is rescaled.  With that context,
one rational Newton step proposes a center and the same Taylor--Rouch\'e
calculation constructs the box around it.

\begin{definition}[protected parent]
\label{def:protected-parent}
For a displayed polynomial $f$ and fixed ambient box $A$, a
\emph{protected parent} is a proper rational box $B$, a refinement
$B\subseteq A$, and a boundary-clearance certificate for $f$ on $B$.
No winding value is part of this record.  A winding-one state supplies a
protected parent by forgetting its winding certificate.
\end{definition}

\leanref{Purelib.RootIsolation.Winding.Residual.ProtectedAmbient;
Purelib.RootIsolation.Winding.Residual.ProtectedAmbient.ofWindingOne}

\begin{theorem}[local residual-to-Hermitian-projective-quadrance converse]
\label{thm:residual-projective-converse}
Let $f$ have degree at least two and carry a squarefree certificate, let $B$
be a rational box carrying a certified winding-one state, and let $q\in\Qi$
satisfy the strict interior predicate of $B$.  For every $\eps\in\Q_{>0}$ one can compute a
positive rational $\Delta(f,B,q,\eps)$ with the following property.  If
\[
                    \Norm{f(q)}<\Delta(f,B,q,\eps),
\]
then exact evaluation either returns the finite certificate $f(q)=0$, or a
finite Newton--Rouché construction returns a rational box $D$ and a RAT
beginning at $D$ such that
\[
 D\subset B,\qquad \Wind_4(f;D)=4,
 \qquad \diamq(D)\le32\eps^2,
 \qquad \chi^2(q,z)\le18\eps^2\quad(z\in D).
\]
The output exposes the selected positive rational radius $r\le\eps$, the
winding certificate $4$, the four strict-containment inequalities,
and the stronger displayed estimate
\[
                   \Norm{z-q}\le18r^2\quad(z\in D).
\]
If only ordinary membership $q\in B$ is initially known, a second computable
positive target first uses one quarter of the certified boundary residual
clearance to prove strict membership and then gives the same conclusion.
\end{theorem}

\begin{proof}
Let $C_2(f,B)>0$ denote the explicit finite Taylor coefficient used by the
selected safe Newton band.  Compute $r$ as the minimum of
\[
 1,\ \frac{\Ree q-x_-}{6},\ \frac{x_+-\Ree q}{6},\
 \frac{\Imm q-y_-}{6},\ \frac{y_+-\Imm q}{6},\ \eps,
 \frac{\tau(B)}{4L_1(f,B)},\ \frac{\eta(B)}{C_2(f,B)},
\]
and put
\begin{equation}
\label{eq:residual-projective-target}
 \Delta(f,B,q,\eps)=
 \min\{\,2L_1(f,B)r^2,\ \eta(B) r^2\,\}.
\end{equation}
Every entry is a displayed positive rational, so $r$ and $\Delta$ are
positive.  The Bézout identity and the first residual budget give
$\Norm{f'(q)}\ge\eta(B)>0$.  Thus the rational Newton point
\[
                        n=q-\frac{f(q)}{f'(q)}
\]
is defined, and exact norm multiplicativity gives
\[
             \Norm{n-q}=\frac{\Norm{f(q)}}{\Norm{f'(q)}}\le r^2.
\]
The four geometric entries keep the eliminant-safe band around $n$ strictly
inside $B$.  The Taylor entry makes the nonlinear remainder subordinate to
the affine factor $f'(q)(Z-n)$ on every selected side.  Finite Rouché transfer
therefore gives $\Wind_4(f;D)=4$, and the standard refiner supplies the total
RAT.

Every band point is within two coordinate radii of $n$, while the Newton
displacement puts each coordinate of $n$ within one radius of $q$.  Hence
both coordinate differences from $q$ are at most $3r$, so
\[
             \Norm{z-q}\le(3r)^2+(3r)^2=18r^2\le18\eps^2.
\]
Equation~\eqref{eq:projective-box-upper} gives the Hermitian-projective
quadrance conclusion.  For the
ordinary-membership form, the minimum additionally includes one quarter of
the finite boundary-clearance lower bound and replaces the four pointwise
gaps by the computed uniform interior-strip radius.
\end{proof}

\leanref{Purelib.Polynomial.Geometry.ResidualChordal.ofResidual;
Purelib.Polynomial.Geometry.ResidualChordal.observeOfResidual}

The preceding theorem accepts a requested radius $\eps$ and tests whether the
residual is small enough for it.  The next corollary runs that relation in the
other numerical direction: starting from the residual $R$, it computes a
rational radius whose square is within a chosen factor $\kappa>1$ of
$R/\mu$.  The factor $\kappa$ is deliberate rational slack that replaces an
exact square root; the finite dyadic bracket is not a new geometric object.

\begin{corollary}[requested-factor rational residual bound]
\label{cor:requested-residual-projective}
Under the hypotheses of \cref{thm:residual-projective-converse}, put
\[
 \mu(f,B)=\min\{\eta(B),\,2L_1(f,B)\}>0
\]
and let $s=s(f,B,q)>0$ be the reference radius obtained from the geometric,
derivative-threshold, and Taylor caps in that proof with the requested cap
$\eps=1$.  For any displayed rational $\kappa>1$, define
\[
 \Lambda_\kappa(f,B,q)=
 \min\left\{\mu(f,B),\frac{\mu(f,B)s^2}{\kappa}\right\}>0.
\]
If $0<R=\Norm{f(q)}<\Lambda_\kappa(f,B,q)$, a bounded dyadic scan followed
by a precomputed finite number of exact rational comparisons computes a
positive rational
$\rho$ satisfying
\[
       \frac{R}{\mu(f,B)}<\rho^2
       \le\frac{\kappa R}{\mu(f,B)}.
\]
The resulting strictly nested winding-one box $D$ and its RAT obey
\[
 \chi^2(q,z)\le\frac{18\kappa}{\mu(f,B)}R\quad(z\in D),
 \qquad
 \diamq(D)\le\frac{32\kappa}{\mu(f,B)}R.
\]
These estimates are local to the supplied winding-one state on $B$.
\end{corollary}

\begin{proof}
Write $\mu=\mu(f,B)$ and put $x=R/\mu$.  Since
$0<R<\Lambda_\kappa\le\mu$, the displayed
rational $x$ lies strictly between zero and one.  Starting at $1$, repeatedly
halve the rational radius.  An explicit natural depth computed from $x$ makes the
final dyadic square no larger than $x$, so this is a bounded recursion rather
than a square-root operation or an unbounded search.  It returns the last
radius $r$ whose square is above $x$; one halving step gives
\[
                         x<r^2\le4x.
\]

Put
\[
              g=\min\left\{1,\frac{\kappa-1}{3}\right\}>0.
\]
The rational inequalities $g\le1$ and $3g\le\kappa-1$ give
\[
              (1+g)^2=1+2g+g^2\le1+3g\le\kappa.
\]
The dyadic bracket supplies the initial rational interval
\[
 \ell_0=\frac r2,\qquad u_0=r,\qquad
 \ell_0^2\le x<u_0^2,\qquad u_0-\ell_0=\ell_0.
\]
Compute the natural number
\[
             N=\left\lceil\frac2g\right\rceil.
\]
The finite geometric-depth estimate of \cref{lem:reciprocal-search} gives
$(2/3)^N\le g$, and hence $(1/2)^N\le g$.  Perform exactly $N$ rational
bisections.  At each step compare $x$ with the square of the midpoint and
retain the half with endpoints $\ell_j,u_j$ satisfying
$\ell_j^2\le x<u_j^2$.  Thus
\[
 u_N-\ell_N
   =2^{-N}(u_0-\ell_0)
   =2^{-N}\ell_0
   \le g\ell_N,
\]
where the last inequality also uses $\ell_0\le\ell_N$.  Consequently
\[
 u_N\le(1+g)\ell_N,
 \qquad
 x<u_N^2\le(1+g)^2\ell_N^2\le\kappa x.
\]
Taking $\rho=u_N$ proves the requested bracket.  Every recursion depth and
every comparison has been computed from the displayed rationals before the
finite recursion begins.

The threshold makes $\kappa(R/\mu)<s^2$, so $\rho<s$.  Moreover
$R<\mu\rho^2$, while $\mu\le\eta(B)$ and
$\mu\le2L_1(f,B)$; hence both residual seed budgets in
\cref{thm:residual-projective-converse} hold.  The same finite
Newton--Rouch\'e construction therefore applies.  Composing its bounds
$\chi^2(q,z)\le18\rho^2$ and $\diamq(D)\le32\rho^2$ with the requested
bracket gives the two displayed estimates.
\end{proof}

\leanref{Purelib.Polynomial.Geometry.DyadicBracket.selectWithin;
Purelib.Polynomial.Geometry.ResidualChordal.AdjustableLinearResidualChordalResult.chordalLinear}

\begin{remark}[why the converse is local and normalized]
A small value of $\Norm{f(q)}$ alone does not identify a track or even a
relevant ambient box: multiplying $f$ by a nonzero Gaussian rational rescales
the residual without changing its winding tracks.  The converse in
\cref{thm:residual-projective-converse} therefore includes the displayed
squarefree Bézout margins and a known containing winding state.
\end{remark}

The protected parent of \Defref{def:protected-parent} need not carry winding
one: the Newton construction uses its quantitative boundary clearance and
ambient refinement together with the already supplied squarefree and
axis-elimination certificates.  This smaller interface permits a global
finite classifier.

The classifier removes the need for the caller to choose that parent.  If the
sample lies outside the coefficient isolation square, leading-term dominance
computes a positive lower bound on its residual, so the small-residual branch
is impossible.  If it lies inside, the outer square itself supplies the
protected parent; the resulting local track is then assigned to one of the
already constructed global tracks.  The three output branches below are just
the decidable outcomes of those two tests plus exact evaluation.

\begin{theorem}[global displayed-residual classifier]
\label{thm:global-residual-classifier}
Let $f$ have degree $d\ge2$ and carry the finite squarefree certificate used
by the box theorem.  Let $K_f=[-R,R]^2$ be its coefficient-computed
isolation square, let $M_f>0$ be the leading--tail dominance margin of
\Defref{def:leading-tail-margin}, and put
\[
 U=\Norm{\operatorname{lc}(f)}(2R^2)^d,
 \qquad
 \Omega(f)=\min\left\{M_f,\frac{M_f^2}{4U}\right\}>0.
\]
For every displayed $q\in\Qi$ and every $\eps\in\Q_{>0}$ one can compute a
positive rational threshold $\Delta_{\rm glob}(f,q,\eps)$ and return exactly
one of the following finite records:
\begin{enumerate}
\item the exact evaluation certificate $f(q)=0$;
\item a winding-one box $D\subset K_f$, its total shrinking RAT, the estimates
\[
 \diamq(D)\le32\eps^2,
 \qquad
 \chi^2(q,z)\le18\eps^2\quad(z\in D),
\]
and a deterministic owner of this RAT among the $d$ tracks already returned
by global isolation;
\item the displayed inequality
\[
             \Delta_{\rm glob}(f,q,\eps)\le\Norm{f(q)},
\]
certifying failure of the explicit small-residual hypothesis.
\end{enumerate}
Inside $K_f$, the threshold is the uniform protected-parent target from
\cref{thm:residual-projective-converse}; outside $K_f$, it is $\Omega(f)$.  In
particular, no sample outside $K_f$ can enter the assigned-box branch with
residual below the stated threshold.
\end{theorem}

\begin{proof}
Whenever a displayed $q$ satisfies
$R^2\le\Norm q\le2R^2$,
Lemma~\ref{lem:leading-tail-annular-margin} gives
\[
 M_f\le \Norm{L_f(q)}-\Norm{T_f(q)}.
\]
For displayed Gaussian rationals $a,b$, a
two-coordinate Cauchy identity gives the finite inequality
\[
 0<m\le\Norm a-\Norm b
 \quad\Longrightarrow\quad
 \frac{m^2}{4\Norm a}\le\Norm{a+b}.
\]
On the annulus, $\Norm{L_f(q)}\le U$, so
$M_f^2/(4U)\le\Norm{f(q)}$.  Beyond the annulus the leading estimate is at
least twice the tail estimate.  The elementary bound
\[
 \frac{\Norm a}{2}-\Norm b\le\Norm{a+b}
\]
then gives $M_f\le\Norm{f(q)}$.  Failure of the closed-square membership
predicate implies
$R^2<\Norm q$ by one of four decidable coordinate comparisons, so the two
cases prove $\Omega(f)\le\Norm{f(q)}$ for every displayed $q$ rejected by
the membership predicate of $K_f$.

For $q\in K_f$, use the boundary clearance of $K_f$ as a
\texttt{ProtectedAmbient}.  A residual below the uniform target first proves
strict interior membership, then the same Newton--Rouch\'e calculation builds
$D$ and its track.  The track refines $K_f$.  The finite maximality assignment
therefore refines it to the coefficient-computed partition scale and returns
the selected output owner; a zero-winding remainder cell is excluded by
additivity and nonnegativity.  Every branch is selected by exact rational
comparisons.
\end{proof}

\leanref{Purelib.Polynomial.Geometry.GlobalResidual.classify;
Purelib.Polynomial.Geometry.GlobalResidual.AssignedResult.owner}

\section{Irreducibility and general polynomials}
\label{sec:general-polynomials}

This section constructs multiplicity presentations, applies the squarefree
RAT construction to their factors, and transfers the resulting finite
observations to the original polynomial.

\begin{definition}[finite irreducibility certificate]
\label{def:irreducibility-certificate}
A \emph{finite irreducibility certificate} for a displayed polynomial
$f\in\Qi[Z]$ records that $f$ is nonzero and has positive degree, together
with the following proper-divisor rule.  For every displayed nonzero
$h,k\in\Qi[Z]$ satisfying $f=hk$ and $\deg h<\deg f$, it certifies
$\deg h=0$.
\end{definition}

\leanref{Purelib.Polynomial.Irreducible.IrreducibleCertificate}

\begin{proposition}[irreducibility supplies squarefreeness]
\label{prop:irreducible-squarefree}
From a finite irreducibility certificate for $f$, the Euclidean algorithm
constructs the squarefree certificate of
\cref{def:squarefree-certificate}.  If $\deg f>1$, the same certificate
excludes every displayed Gaussian-rational linear factor.
\end{proposition}

\begin{proof}
Compute $g=\gcd(f,f')$ and the exact factorization $f=gh$.  The derivative
calculation gives $\deg g\le\deg f'<\deg f$, so the proper-divisor rule makes
$g$ constant.  Rescaling the computed Bézout cofactors gives
$Uf+Vf'=1$.  If $\deg f>1$, a displayed linear factor would itself be a
nonconstant proper divisor, contrary to the same rule.
\end{proof}

\leanref{Purelib.Polynomial.Irreducible.squarefreeOfIrreducible}

For completeness, arbitrary polynomials can be reduced constructively.

The purpose of this preprocessing is narrower than full factorization.  A
repeated factor should describe one geometric site carrying a larger integer
weight, not several indistinguishable tracks.  The algorithm therefore splits
only until it has pairwise-coprime squarefree factors and records how often
each occurs.  The box theorem is then run once per squarefree factor, while
the multiplicity is transported as finite integer data.

\begin{lemma}[squarefree certificate or smaller exact split]
\label{lem:squarefree-or-split}
Let $p\in\Qi[Z]$ be displayed, normalized, nonzero, and of positive degree.
A finite Euclidean calculation returns exactly one of the following records:
\begin{enumerate}[label=(\roman*)]
\item a squarefree certificate for $p$ in the sense of
\Defref{def:squarefree-certificate};
\item displayed normalized nonzero polynomials $g,q$ of positive degree,
each of degree strictly below $\deg p$, together with the identities
\[
                 p=gq,
 \qquad
                 \deg g+\deg q=\deg p.
\]
\end{enumerate}
\end{lemma}

\begin{proof}
Compute the extended Euclidean GCD of $p$ and $p'$.  Positive degree makes
$p'$ nonzero and gives $\deg p'<\deg p$, while nonzeroness of $p$ makes the
computed GCD $g$ nonzero.  Its degree is decidably either zero or positive.
If $\deg g=0$, rescale the returned B\'ezout identity exactly as in the proof
of Proposition~\ref{prop:irreducible-squarefree}; this gives the squarefree certificate
of \Defref{def:squarefree-certificate}.

If $\deg g>0$, exact Euclidean division returns the normalized cofactor $q$
and the identity $p=gq$.  Since $g$ divides $p'$, one has
$\deg g\le\deg p'<\deg p$.  Exact degree additivity gives
$\deg g+\deg q=\deg p$; hence $q$ is also nonzero, normalized, of positive
degree, and of degree strictly below $\deg p$.
\end{proof}

\leanref{Purelib.Polynomial.Squarefree.squarefreeOutcome;
Purelib.Polynomial.Factorization.factor_natDegree_add;
Purelib.Polynomial.Factorization.decomposeNormalized}

\begin{lemma}[finite coprime multiplicity fold]
\label{lem:coprime-multiplicity-fold}
Let $P$ be a displayed normalized positive-degree polynomial, and suppose a
nonempty finite list of normalized positive-degree squarefree occurrences
has product $P$ and degree sum $\deg P$.  A terminating finite GCD
fold converts that list into a finite squarefree-multiplicity presentation of
$P$ in the sense of \Defref{def:squarefree-multiplicity-presentation}.
In particular, the returned positive-degree basis is nonempty, its factors
are pairwise coprime with displayed B\'ezout certificates, and its weighted
product and degree identities recover those of the occurrence list.
\end{lemma}

\begin{proof}
Insert the occurrences successively into an initially empty basis.  When a
new occurrence is coprime to the current head, retain the head and recurse
through the tail.  Otherwise compute their nonconstant GCD and replace the
overlap by its common factor and the two exact cofactors.  Squarefreeness
supplies the pairwise-coprimality certificates among those three parts and
preserves coprimality with the unvisited tail.  Every recursive insertion
consumes one entry of the displayed finite list.

The insertion invariant records pairwise coprimality, the weighted product
identity, and the weighted degree identity.  After the fold, remove only
degree-zero basis parts and multiply their displayed values into one nonzero
unit.  The input has positive degree, so the remaining positive-degree basis
is nonempty.  These invariants are precisely the data required by
\Defref{def:squarefree-multiplicity-presentation}.
\end{proof}

\leanref{Purelib.Polynomial.Multiplicity.insertPart;
Purelib.Polynomial.Multiplicity.buildBasis;
Purelib.Polynomial.Factorization.Decomposition.positiveMultiplicityProductIdentity;
Purelib.Polynomial.Factorization.Decomposition.positiveMultiplicityEntries_pairwise;
Purelib.Polynomial.Factorization.Decomposition.positiveMultiplicityEntries_nonempty;
Purelib.Polynomial.Factorization.Decomposition.positiveMultiplicityDegree}

\begin{proposition}[finite squarefree and multiplicity preprocessing]
\label{prop:factor}
There is a total coefficient algorithm which distinguishes the zero
polynomial, a nonzero constant, and a positive-degree polynomial.  In the
positive-degree case it returns a nonempty finite list of normalized,
positive-degree, squarefree occurrences whose product is the
normalized input and whose degrees sum to its degree.  It also returns the
resulting finite squarefree-multiplicity presentation in the sense of
\Defref{def:squarefree-multiplicity-presentation}.
\end{proposition}

\begin{proof}
The normalized coefficient list decides whether the input is zero, a nonzero
constant, or of positive degree.  In the positive-degree branch, repeatedly
apply Lemma~\ref{lem:squarefree-or-split}.  Retain a returned squarefree factor;
in the split branch recurse on the two strictly smaller natural degrees and
append their occurrence lists.  Polynomial multiplication and degree
additivity preserve the two stated invariants, and recursion
terminates by the strict degree decrease.  Apply
Lemma~\ref{lem:coprime-multiplicity-fold} to the resulting nonempty occurrence
list to obtain the presentation of
\Defref{def:squarefree-multiplicity-presentation}.
\end{proof}

\leanref{Purelib.Polynomial.Factorization.preprocess}

Run the box theorem once for each positive-degree basis factor and attach its displayed multiplicity.  Repeated factors therefore do not create duplicate box machines.  This preprocessing is intentionally not claimed to be an irreducible factorization algorithm; no such algorithm is needed for the squarefree box theorem proved here.

The next theorem has two complementary readings.  If the value of a finite
product is very small, a bounded scan can exhibit at least one small factor.
If a particular factor has already been selected and the product of the
others has a positive lower bound, then smallness can be transferred back and
forth with the selected factor, with its multiplicity appearing as an
ordinary natural power.  No radical is needed to undo that power.

\begin{theorem}[finite multiplicity residual transfer]
\label{thm:multiplicity-residual}
Write the computed positive multiplicity basis as
\[
             p(q)=u\prod_{j=1}^{s}h_j(q)^{m_j},
 \qquad u\ne0,\quad m_j\in\N_{>0},
\]
where this is an equality of displayed Gaussian-rational evaluations.  Put
$U_0=\Norm u>0$ and $N=\sum_jm_j$.  For every displayed
$\delta\in\Q_{>0}$,
\[
             \Norm{p(q)}<U_0\delta^N
\]
constructively selects an index $j$ with
$\Norm{h_j(q)}<\delta$.

For a preselected entry $h=h_j$ of multiplicity $m=m_j$, expose the finite
list split before and after that entry and let $c(q)$ be the displayed product
of all the other positive factors.  If a displayed certificate gives
\[
       0<\beta\le U_0\Norm{c(q)}
\]
and a local residual theorem gives
$0\le X\le\kappa\Norm{h(q)}$ with $0\le\kappa$, then
\begin{equation}
\label{eq:multiplicity-projective-power}
 X^m\le
 \left(\kappa\max\{1,\beta^{-1}\}\right)^m\Norm{p(q)}.
\end{equation}
In particular $X$ may be the affine Hermitian projective quadrance supplied
by the requested-factor bound of
\cref{cor:requested-residual-projective}, with its
displayed coefficient used as $\kappa$ here.  The result uses the natural power $m$;
it does not extract an $m$-th root or introduce a fractional exponent.
\end{theorem}

\begin{proof}
Exact norm multiplicativity turns the basis identity into
\[
 \Norm{p(q)}=U_0\prod_j\Norm{h_j(q)}^{m_j}.
\]
A bounded scan tests the finitely many inequalities
$\Norm{h_j(q)}<\delta$.  If every test failed, monotonicity of natural powers
and finite products would give the reverse bound
$U_0\delta^N\le\Norm{p(q)}$, a contradiction.  The scan therefore returns
the first proof-bearing small factor.

For the preselected entry, the exact list split gives
\[
 \Norm{p(q)}=U_0\Norm{c(q)}\Norm{h(q)}^m,
 \qquad
 \beta\Norm{h(q)}^m\le\Norm{p(q)}.
\]
Raise the local inequality to the positive natural power $m$.  Since
$M=\max\{1,\beta^{-1}\}\ge1$ and $m>0$, one has
$\beta^{-1}\le M\le M^m$.  Multiplying the preceding inequalities proves
\eqref{eq:multiplicity-projective-power}.  All products and scans are structural
recursions on displayed finite lists.
\end{proof}

\leanref{Purelib.Polynomial.MultiplicityResidual.selectSmallFactor;
Purelib.Polynomial.MultiplicityResidual.chordalPow_le_integerPowerConstant_mul_sourceResidual}

\section{A finite rational-function argument principle}
\label{sec:rfap}

This section combines a reduced quotient with multiplicity presentations to
produce jointly separated signed sites and compute the rational-function
action principle.

Let $F,G\in\Qi[Z]$ be displayed dense polynomials and suppose that $G$ is
nonzero.  The identically zero numerator is returned by a separate
constructor.  In the remainder of this section $F$ is also nonzero.

Conceptually, the construction performs three operations.  It first cancels
the common polynomial factor, since a displayed quotient has no zero/pole
distinction there.  It then applies the polynomial box theorem separately to
the squarefree factors of the reduced numerator and denominator.  Finally it
labels numerator sites positively, denominator sites negatively, and uses
Bézout certificates to separate every pair across both sides.  The outer
winding is consequently the positive degree total minus the negative degree
total.

\subsection{Executable cancellation}

Run extended Euclid on the normalized pair.  Its finite output supplies
displayed polynomials
\[
 H,\quad F_0,\quad G_0,\quad U,\quad V\in\Qi[Z]
\]
and the identities
\begin{equation}
\label{eq:rfap-reduction}
 F=HF_0,\qquad G=HG_0,\qquad UF_0+VG_0=1.
\end{equation}
The Bézout identity is the finite certificate that the displayed reduced
cofactors have no common nonconstant factor.
Here $H$ and $G_0$ are nonzero, and $F_0$ is nonzero when $F$ is.  The same
finite division records prove
\begin{equation}
\label{eq:rfap-degree-balance}
 \deg H+\deg F_0=\deg F,\qquad
 \deg H+\deg G_0=\deg G.
\end{equation}
At a displayed $z\in\Qi$ for which
$G(z)\ne0$, cancellation gives the exact value identity
\[
       F(z)G(z)^{-1}=F_0(z)G_0(z)^{-1}.
\]
No value is assigned at a point where the denominator vanishes.

On any rational boundary clear for both reduced cofactors, the generic
winding of \cref{def:rational-function-winding}, realized explicitly in
\cref{sec:winding-realization}, applies to $F_0/G_0$.  In particular,
\cref{thm:rfap-quotient} computes its winding by subtracting the polynomial
windings.  Cancellation therefore introduces no second, quotient-specific
notion of winding.

The coefficient bounds from \cref{lem:eval-bound} compute a positive radius
for each nonzero polynomial beyond which its leading term dominates.  Take
the maximum of the two radii.  The resulting finite symmetric square is
clear for both reduced cofactors, so \cref{thm:global-wind,thm:rfap-quotient}
give the following result, including constant cofactors.

\begin{corollary}[outer reduced-degree winding]
\label{cor:rfap-outer}
There is one coefficient-computed positive rational $R$ such that
\begin{equation}
\label{eq:rfap-outer}
 \Wind_4(F_0/G_0;K_R)=4(\deg F_0-\deg G_0).
\end{equation}
\end{corollary}

\begin{proof}
Choose $R$ as the maximum of the two coefficient-computed dominance radii.
Theorem~\ref{thm:global-wind} gives the two polynomial windings on $K_R$;
Theorem~\ref{thm:rfap-quotient} subtracts them, giving
\eqref{eq:rfap-outer}.
\end{proof}

\leanref{Purelib.RationalFunctions.FunctionArgumentPrinciple.NonzeroReducedInput.outer_rationalWinding_eq_degreeDifference}

\subsection{Signed multiplicity sites}

Apply \cref{prop:factor} separately to the two reduced cofactors.  It returns
nonzero units and pairwise-coprime squarefree positive-degree factors
\begin{equation}
\label{eq:rfap-weighted-factorization}
 F_0=u_+\prod_{j=1}^{r}p_j^{m_j},\qquad
 G_0=u_-\prod_{k=1}^{s}q_k^{n_k},
\end{equation}
with $m_j,n_k>0$ and
\begin{equation}
\label{eq:rfap-weighted-degrees}
 \sum_{j=1}^{r}m_j\deg p_j=\deg F_0,
 \qquad
 \sum_{k=1}^{s}n_k\deg q_k=\deg G_0.
\end{equation}
Factors on the same side carry the finite pairwise Bézout certificates from
the multiplicity construction.  Equation \eqref{eq:rfap-reduction} also
gives a Bézout certificate between every $p_j$ and every $q_k$: exact finite
factorization transports coprimality to displayed nonzero factors.

Each squarefree factor is handled once.  A degree-one factor returns its
explicit Gaussian-rational zero with an exact evaluation certificate.  It is
the only branch in which a zero or pole point is returned.  A factor of
degree at least two returns exactly its degree many shrinking winding-one
tracks from \cref{thm:main}.  Thus the output is a finite family of
\emph{sites}; a site is either an exact rational terminal or a box track.

This is the finite rational-function analogue of a geometric zero or pole
location.  The sign records whether the quotient turns in the numerator or
denominator direction, the multiplicity $m$ records the contribution
$\pm4m$, and the terminal/track alternative says how it is
observed: exactly when linear and by repeatable boxes otherwise.

The separation cases now require no additional definition.  An exact
terminal $q=a+ib$ contributes the degenerate support record $(a,a,b,b)$,
whereas a RAT contributes the proper box returned by a finite call.
Substituting two degenerate records, one of each kind, or two proper boxes
into the four inequalities of
\Defref{def:finite-observation-separation} gives respectively the
terminal/terminal, terminal/RAT, and RAT/RAT certificates.  Signs and
multiplicities do not alter this geometric witness.

Fix a numerator factor $p_j$ and one of its higher-degree tracks.  The
Bézout evaluation bounds compute one positive rational refinement request
small enough to protect the track simultaneously from
\[
 C_j=u_+\prod_{\ell\ne j}p_\ell^{m_\ell}
 \quad\hbox{and}\quad G_0.
\]
On the returned box, $C_j$ and $G_0$ have positive boundary clearances and
zero winding.  The finite power and product laws therefore give
\[
 \Wind_4(p_j^{m_j}C_j/G_0;B)=4m_j.
\]
The presentation gives $F_0=p_j^{m_j}C_j$.  The denominator calculation is
symmetric and gives
\[
 \Wind_4(F_0/q_k^{n_k}D_k;B)=-4n_k,
 \qquad
 D_k=u_-\prod_{\ell\ne k}q_\ell^{n_\ell}.
\]
Joint separation is also quantitative and finite.  Tracks belonging to one
squarefree factor are separated by its polynomial isolation output.  For two
different factors, refine the first until its residual forces a positive
Bézout lower bound for the second, and refine the second until its residual
is below that bound.  Overlap would contradict the two rational inequalities.
For an exact rational terminal and a track, the same lower bound excludes the
displayed point from a finite refinement of the track.  Two exact terminals
of coprime factors are unequal because substituting a common displayed point
in their Bézout identity would give $1=0$.

\begin{proof}[Proof of \Thmref{thm:rfap-general}]
Cancellation, quotient winding, multiplicity preprocessing, local signed
winding, and the two degree totals were proved above.  Enumerate every focused
position in each finite multiplicity list.  Store the exact terminal or
higher isolation result at that position, its internal separation and atomic
ownership data, and all same-side separation witnesses.  Then enumerate the
finite Cartesian product of numerator and denominator positions and apply
the cross-factor Bézout construction just described.  Constants contribute
empty site families.  Finally store \eqref{eq:rfap-outer} in the same result
record; \cref{def:rational-function-winding} gives the normalized identity in
\Thmref{thm:rfap-general}.  All iterations are over displayed finite lists or
natural degrees.
\end{proof}

\leanref{Purelib.RationalFunctions.Sites.prove;
Purelib.Papers.RFAP.Paper.rationalFunctionArgumentPrinciple}

The ownership sentence in \cref{thm:rfap-general} is deliberately atomic.
It would be false to say that every arbitrary nonzero-winding box has exactly
one site: a larger rational box may surround several zeros and poles.  The
proved maximality statement concerns winding-one tracks for one squarefree
factor.  The returned finite partition has zero-winding remainder cells for
that factor, and every competing atomic track is assigned to one output cell.

\begin{corollary}[squarefree signed specialization]
\label{cor:rfap-squarefree}
Suppose both nonzero reduced cofactors $F_0$ and $G_0$ are supplied with
finite squarefree certificates.  The RFAP algorithm has exactly $\deg F_0$
numerator sites and $\deg G_0$ denominator sites.  Constants contribute no
site, degree-one cofactors contribute exact Gaussian-rational terminals, and
higher cofactors contribute degree-indexed box tracks.  Every numerator
track has winding $+4$, every denominator track has winding $-4$, the two
families are jointly separated, and \eqref{eq:rfap-outer} holds.
\end{corollary}

\begin{proof}
Apply the squarefree box theorem directly to each reduced cofactor instead
of multiplicity preprocessing.  Protect each higher track from the opposite
cofactor using their reduced Bézout certificate.  The opposite winding is
zero, so \cref{thm:rfap-quotient} gives $4-0$ on a numerator track and
$0-4$ on a denominator track.  The degree-one and constant branches, joint separation,
atomic ownership, and outer winding use the constructions already described.
\end{proof}

\leanref{Purelib.RationalFunctions.FunctionArgumentPrinciple.proveSquarefree;
Purelib.Papers.RFAP.Paper.squarefreeRationalFunctionArgumentPrinciple}

\subsection{Two-way metric and residual calibration}

The word here is \emph{residual}, not the meromorphic residue.  For displayed
$P,Q$ and a displayed $q$ with $Q(q)\ne0$, put
\begin{equation}
\label{eq:quotient-residual}
       \mathcal R_{P/Q}(q)=\Norm{P(q)/Q(q)}
       =\frac{\Norm{P(q)}}{\Norm{Q(q)}}.
\end{equation}

There are again two directions.  Moving far enough down a numerator track
makes this quotient residual uniformly small because the numerator becomes
small while the protected denominator stays bounded away from zero.
Conversely, a sufficiently small quotient residual makes the numerator
residual small, because the denominator also has an explicit upper bound; the
polynomial converse can then build a nearby numerator track.  At a denominator
site the same statement is applied to the reciprocal, so its smallness is the
finite certificate for the original quotient's displayed blow-up.

\begin{definition}[local quotient context and protected state]
\label{def:protected-quotient-state}
For displayed coprime polynomials $P,Q$ with $\deg P>0$ and a fixed ambient
box $A$, a \emph{local quotient context} consists of a RAT for $P$, a current
winding-one state on that RAT, and a displayed Bézout certificate
$UP+VQ=1$.  Finite refinement using this identity produces its
\emph{protected state}.  The same coefficient calculations provide positive
rationals
\[
  0<\eta_Q\le\Norm{Q(z)}\le M_Q
\]
for every displayed $z$ in that state or any refinement of it.  The lower
bound is the Bézout protection margin and the upper bound is a Horner
majorant on $A$.  At a denominator site, interchange $P$ and $Q$; the source
track is then a track for $Q$ protected from $P$.
\end{definition}

\leanref{Purelib.RationalFunctions.ResidualChordal.LocalQuotientContext;
Purelib.RationalFunctions.ResidualChordal.LocalQuotientContext.protectedState;
Purelib.RationalFunctions.ResidualChordal.LocalQuotientContext.oppositeMargin;
Purelib.RationalFunctions.ResidualChordal.LocalQuotientContext.oppositeMajorant}

\begin{theorem}[two-way RFAP residual bounds]
\label{thm:rfap-residual-bounds}
Let $B$ be a rational box carrying any winding-one state for $P$ that refines
the protected state inside the fixed ambient box $A$.
\begin{enumerate}[label=(\roman*)]
\item Every displayed $q\in B$ satisfies the forward estimate
\begin{equation}
\label{eq:rfap-forward-residual}
 \mathcal R_{P/Q}(q)
 \le \frac{4L_1(P,A)}{\eta_Q}\diamq(B).
\end{equation}
Thus a numerator track makes the reduced rational function small.  On a
denominator track, with $P,Q$ interchanged, the same estimate makes the
reciprocal rational function small.
\item Assume $P$ is squarefree of degree at least two, $q$ satisfies the
strict interior predicate of $B$, and $\eps>0$.  Let
$\Delta(P,B,q,\eps)$ be the polynomial target from
\cref{eq:residual-projective-target}.  The positive quotient target
\begin{equation}
\label{eq:rfap-converse-target}
       \Delta_{P/Q}(B,q,\eps)=\frac{\Delta(P,B,q,\eps)}{M_Q}
\end{equation}
has the following constructive converse: if
$\mathcal R_{P/Q}(q)<\Delta_{P/Q}(B,q,\eps)$, exact evaluation either returns
$P(q)=0$, or constructs a winding-one box $D\subset B$, its total RAT, and
\[
 \diamq(D)\le32\eps^2,
 \qquad \chi^2(q,z)\le18\eps^2\qquad(z\in D).
\]
\item For each $\eps>0$, a bounded rational refinement of the protected
$P$-track constructs a state on a box $B_\eps$ on which
\begin{equation}
\label{eq:rfap-epsilon-residual}
              \mathcal R_{P/Q}(q)<\eps^2
              \qquad(q\in B_\eps).
\end{equation}
In the pole orientation, $P=G_0$ and $Q=F_0$.  Exact evaluation either
returns the displayed pole certificate $G_0(q)=0$, or else
\begin{equation}
\label{eq:rfap-pole-blowup}
              \Norm{F_0(q)/G_0(q)}\ge\frac{1}{\eps^2}.
\end{equation}
\end{enumerate}
All constants and refinement requests are computed from displayed
coefficients, boxes, and Bézout data.
\end{theorem}

\begin{proof}
The polynomial estimate in \cref{thm:main} gives
$\Norm{P(q)}\le4L_1(P,A)\diamq(B)$, while protection gives
$\eta_Q\le\Norm{Q(q)}$; division proves
\eqref{eq:rfap-forward-residual}.  Conversely,
\eqref{eq:quotient-residual}, the upper bound $\Norm{Q(q)}\le M_Q$, and
\eqref{eq:rfap-converse-target} give
\[
 \Norm{P(q)}
 =\mathcal R_{P/Q}(q)\Norm{Q(q)}
 <\frac{\Delta(P,B,q,\eps)}{M_Q}M_Q
 =\Delta(P,B,q,\eps).
\]
Apply \cref{thm:residual-projective-converse}.

For (iii), refine the $P$-track at the explicit uniform residual request with
polynomial target $\eta_Q\eps^2$.  Every displayed point in the returned box
has $\Norm{P(q)}<\eta_Q\eps^2$ and
$\Norm{Q(q)}\ge\eta_Q$, proving
\eqref{eq:rfap-epsilon-residual}.  If both displayed values are nonzero,
exact norm multiplicativity gives
\[
 \mathcal R_{P/Q}(q)\mathcal R_{Q/P}(q)=1.
\]
Inverting the strict upper bound yields the weak lower bound
\eqref{eq:rfap-pole-blowup}.  If $P(q)=0$, the algorithm reports that exact
finite evaluation instead of assigning a value at the pole.
\end{proof}

\leanref{Purelib.RationalFunctions.ResidualChordal.LocalQuotientContext.quotientResidualSq_le_diagonal;
Purelib.RationalFunctions.ResidualChordal.LocalQuotientContext.oppositeQuotientSq_ge_inv_epsilonSq}

\begin{example}[coprime squarefree quotient]
For $F(Z)=Z$ and $G(Z)=Z-1$, the identity
\[
             1=F(Z)-G(Z)
\]
is a finite coprimality certificate.  Both derivatives are the constant one.
The two sites are therefore the exact displayed points $0$ and $1$, carrying
the exact numerator-zero and denominator-zero evaluations respectively; the
outer winding is $4-4=0$.  Their cross-separation witness is the rational
coordinate gap $1$.
\end{example}

\begin{example}[repeated factor]
Take $F(Z)=(Z-1)^2$, $G(Z)=1$, and the box
$B=(-2,2,-2,2)$.  The finite product calculation gives
\[
       \Wind_4(F/G;B)
       =\Wind_4(Z-1;B)+\Wind_4(Z-1;B)-\Wind_4(1;B)=8.
\]
This is proved from two explicit boundary-clearance certificates, the finite
product formula, and \cref{thm:rfap-quotient}.  It demonstrates why
cancellation alone does not imply the $\pm4$ specialization and why a
repeated geometric site carries multiplicity two rather than being duplicated.
\end{example}

\section{Finite size of the RFAP-I calculation}
\label{sec:rfap-one-calculation-size}

This section bounds the finite arithmetic used to produce one requested
refinement state.  In the glossary's terms, it covers coefficient clearing,
the two eliminants and their B\'ezout data, the safe cuts, and every Sturm
winding used by repeated subdivision.  The purpose is not an optimized
complexity classification.  It is the more basic constructive assertion that
the sizes of the intermediate displayed records, as well as the number of
calculation clauses, are bounded before the calculation is run.

For this purpose the executable route carries an ordinary rational as a
coprime numerator and positive denominator.  Before addition it removes the
gcd of the two denominators; before multiplication it cross-cancels; after
each operation it performs a final constructive gcd reduction.  Gaussian
rationals are pairs of these reduced records.  These canonical records, and
the normalized polynomials, matrices, and chains built from them, are the
active algebraic objects used in the proofs above.  Structural equality of
canonical fractions is equivalent to the displayed cross-multiplication
identity, while normalization returns an explicit certificate of that
identity.  No erasure into a quotient-valued rational is part of the
paper-facing proof route.

This reduced-fraction layer is slightly more infrastructure than the
headline RFAP-I theorem alone requires.  We have embedded it in the Lean
development partly to enable the companion approximation work, where bounds
on the numerators and denominators carried by actual calculations become
mathematical output rather than merely an implementation concern.  Keeping
the reduced records and their normalization certificates lets that later
work reuse the RFAP-I algorithms through the same stable canonical algebraic
interface, without assuming that a small rational value has a small
unreduced representative.

Write $\ell(n)=1+\lfloor\log_2 n\rfloor$ for $n>0$, with
$\ell(0)=1$.  For a reduced rational record, its component magnitude is the
maximum of the absolute numerator and denominator; for a Gaussian record it
is the maximum over both coordinates.  Polynomial and matrix magnitudes are
the corresponding maxima over their finite coefficient lists.  All uses of
$\ell$ below are implemented by recursion on natural numbers, not by real
logarithms.

\begin{definition}[RFAP-I calculation budgets]
\label{def:rfap-one-calculation-budgets}
Fix a displayed polynomial $f$, a proper boundary-clear winding state on a
rational box, its displayed vertical and horizontal elimination certificate,
and a positive rational refinement request.  Let $r$ be the geometric depth
computed from that request.

The magnitude recurrence $M$ is the maximum of the following source
recurrences:
\begin{enumerate}
\item the coefficient recurrence for the vertical separated-Sylvester matrix
and its complete fraction-free Schur inverse;
\item the same recurrence after the quarter turn used for the horizontal
eliminant;
\item the product denominator and cleared-coefficient envelope of the input
coefficient list;
\item the normalized safe-cut recurrence applied $r$ times to the four
endpoints of the initial box;
\item the normalized coefficient and candidate-evaluation envelope for the
two displayed elimination polynomials; and
\item the coefficient and endpoint-evaluation envelopes of all normalized
Sturm chains used for the four child windings in those $r$ rounds.
\end{enumerate}
Put $B=\ell(M)$.  The step budget $T$ is the sum of the displayed
coefficient-clearing count, the two finite Schur-recursion counts, the two
bounded candidate scans and child-selection count in every round, and the
field-operation budgets of those child windings.  Both $T$ and $B$ are
primitive finite recurrences in the supplied records.
\end{definition}

\begin{theorem}[RFAP-I calculation-size certificate]
\label{thm:rfap-one-calculation-size}
For every input to \cref{def:rfap-one-calculation-budgets}, the computed
numbers $T$ and $B$ have the following properties.
\begin{enumerate}
\item Coefficient clearing, both separated-Sylvester--Schur calculations,
all safe-coordinate scans, and all subdivision windings fit their respective
summands of $T$.
\item Every coefficient of both normalized separated-Sylvester scalars and
every polynomial entry of their scaled inverses has component width at most
$B$.  In particular this covers the finite polynomial data from which the
displayed B\'ezout cofactors are extracted.
\item The common coefficient scale and both integral components of every
cleared coefficient have width at most $B$.
\item Every endpoint of the retained box after the requested refinement has
width at most $B$.
\item For every child box examined during refinement, every coefficient of
every normalized signed-remainder-chain member and every normalized endpoint
evaluation used by its four edge indices has width at most $B$.
\end{enumerate}
The normalized chains compute exactly the headline Cauchy indices, so the
bounded execution produces the same windings and retained refinement state as
the construction proved above.
\end{theorem}

\begin{proof}
One normalized field operation has an explicit component recurrence.  Horner
induction therefore bounds evaluation of a polynomial with bounded
coefficients and length.  Long-division induction applies the same recurrence
to every signed-remainder-chain member; its degree drops strictly at each
recursive call.  Fraction-free Schur induction bounds the current matrix and
then every entry of the recursively assembled scaled inverse.  These two
inductions also retain coefficient-list bounds, so their evaluation envelopes
are source-computed rather than maxima of completed traces.

Each safe-coordinate search tests exactly $d_0+1$ candidates, where $d_0$ is
the supplied eliminant degree bound.  The selected coordinate is one of those
candidates, and its reduced numerator and denominator satisfy the same
fixed operation envelope.  Induction on $r$ carries this bound through every
retained child.  Summing the finite clause bounds gives $T$; inclusion in the
nested maximum $M$ and monotonicity of $\ell$ give $B$.  Finally, the erasure
theorems commute with polynomial arithmetic, long division, evaluation, and
the Schur construction.  They identify the normalized execution with the
coefficientwise identities and windings already used by the paper-facing
theorems.
\end{proof}

\leanref{Purelib.RootIsolation.Sturm.Normalized.erase_canonical_signedRemainderChain;
Purelib.RootIsolation.Sturm.Normalized.chainMember_resource;
Purelib.Polynomial.NormalizedEliminant.erase_normalizedEliminant;
Purelib.Polynomial.NormalizedBezoutResource.normalizedSeparatedData_within;
Purelib.Polynomial.ClearingCalculationSize.scaledGaussianCoefficient_component_le;
Purelib.RootIsolation.Isolation.CalculationSize.refineTo_box_within;
Purelib.RootIsolation.CalculationSize.stageThreeStepBudget;
Purelib.RootIsolation.CalculationSize.stageThreeBitBudget;
Purelib.Papers.RFAP.Paper.stageThreeResourceCertificate}

The estimates deliberately prefer transparent uniform recurrences to tight
ones.  Fraction-free subresultant estimates, a common subdivision grid, or a
more detailed machine-cost model could improve them.  None is needed for the
present conclusion: every named intermediate record on the requested RFAP-I
path has a displayed finite width bound.

\section{Exact cubic example: \texorpdfstring{$Z^3-(1+i)$}{Z cubed minus (1+i)}}

This section traces one displayed polynomial through its squarefree
certificate, coefficient isolation square, box winding, and winding-one
boxes.

Consider
\[
            f(Z)=Z^3-(1+i).
\]
Its normalized high-to-low coefficient list and derivative list are
\[
 [a_3,a_2,a_1,a_0]=[1,0,0,-(1+i)],\qquad
 [f']=[3,0,0].
\]
It is Eisenstein in $\Zi[Z]$ at the Gaussian prime $1+i$: the two
intermediate coefficients and the constant term are divisible by $1+i$,
the constant term is not divisible by $(1+i)^2$, and the leading coefficient
is not divisible by $1+i$.  Thus it is irreducible in the ordinary algebraic
terminology.

The box construction itself needs only an exact squarefree certificate.  Take
the two polynomials
\[
 U(Z)=\frac{-1+i}{2}=-\frac1{1+i},\qquad
 V(Z)=\frac{1-i}{6}Z=\frac{Z}{3(1+i)}.
\]
The coefficient of $Z^3$ in $Uf+Vf'$ and its constant coefficient are
respectively
\[
 \frac{-1+i}{2}+3\frac{1-i}{6}=0,
 \qquad
 -\frac{-1+i}{2}(1+i)=1,
\]
while every other coefficient is zero.  Hence the displayed coefficientwise
identity is $Uf+Vf'=1$.

We next run the coefficient construction of
\cref{def:coefficient-isolation-square}.  Since
\[
 C_f=3\bigl(\Norm{-(1+i)}+2\Norm0+4\Norm0\bigr)=6,
 \qquad
 R_f=\max\{1,6/\Norm1+1\}=7,
\]
the coefficient isolation square is $K_f=K_7=(-7,7,-7,7)$.  Its exact
leading--tail margin is
\[
 M_f=(\Norm1\cdot7^2-C_f)(7^2)^{3-1}
     =(49-6)49^2=103243>0.
\]
The single global application of \cref{thm:global-wind} therefore gives
$\Wind_4(f;K_f)=12$.

The coefficient polynomials actually passed to the eliminant and edge
calculations are obtained by expanding
\[
 f(x+iy)=A_f(x,y)+iB_f(x,y),
 \qquad
 A_f=x^3-3xy^2-1,\quad B_f=3x^2y-y^3-1.
\]
For example, on a horizontal edge put
$x(t)=\alpha+(\beta-\alpha)t$ and $y(t)=c$; the two input lists are the
coefficients of
\[
 x(t)^3-3c^2x(t)-1,
 \qquad
 3c x(t)^2-c^3-1.
\]
On a vertical edge one instead puts $x(t)=c$ and
$y(t)=\alpha+(\beta-\alpha)t$, obtaining
\[
 c^3-3c y(t)^2-1,
 \qquad
 3c^2y(t)-y(t)^3-1.
\]
The separated Schur--Sylvester recursion on these coefficient data constructs
the two axis eliminants and computes the common finite degree bound
$D_f=168$.

The local coefficient bounds are equally explicit.  On $K_7$ the rational
radius parameter in \cref{def:local-horner-bounds} is $R_{K_7}=98$, and the
Horner triples are
\[
 (0,0,0),\ (2,0,0),\ (392,4,0),\ (76832,1568,8),\
 (15059076,460992,4704).
\]
Consequently $L_1(f,K_7)=460992$ and $L_2(f,K_7)=4704$.  Applying the same
finite evaluation scan to the displayed cofactors gives
\[
 M_U=1,\qquad M_V=\frac{196}{9},\qquad
 \tau(K_7)=\frac1{16},\qquad \eta(K_7)=\frac9{3136}.
\]
Thus the principal numerical outputs of the global and local coefficient
stages are explicit; the degree-$168$ eliminant lists themselves are produced
by the stated finite recursion rather than typeset.

A certificate-level execution of the safe subdivision records the following
three dyadic winding-one boxes.  The edge columns display the four
Cauchy--Sturm contributions in positive boundary order; their sum is four in
every row.  Each contribution uses the factor-two convention fixed in
\cref{def:index}.  This is where the direct edge calculations are informative;
repeating them on $K_f$ would merely recompute the global winding already
supplied by \cref{thm:global-wind}.

\begin{center}
\begin{tabularx}{\textwidth}{@{}cXlc@{}}
\toprule
track & rational box $[x_-,x_+]\times[y_-,y_+]$ & edge indices & $\Wind_4(f;B)$\\
\midrule
$1$ & $[138/128,140/128]\times[36/128,38/128]$ & $(0,2,0,2)$ & $4$\\
$2$ & $[-102/128,-101/128]\times[101/128,102/128]$ & $(2,0,2,0)$ & $4$\\
$3$ & $[-38/128,-36/128]\times[-140/128,-138/128]$ & $(0,2,0,2)$ & $4$\\
\bottomrule
\end{tabularx}
\end{center}

For one concrete view of the coefficient arithmetic behind the table,
the bottom edge of the first box has
$x(t)=(69+t)/64$ and $y(t)=9/32$.  After clearing the common positive
denominator $64^3$, its two Cauchy--Sturm inputs are
\[
 A_e(t)=t^3+207t^2+13311t-703,
 \qquad
 B_e(t)=54t^2+7452t-10882.
\]
The finite signed-remainder calculation gives contribution zero, the first
entry in that row.  The other eleven edge pairs are obtained by the same two
substitution formulas above.

Taking the rational centers gives
\[
 q_1=\frac{139+37i}{128},\qquad
 q_2=\frac{-203+203i}{256},\qquad
 q_3=\frac{-37-139i}{128}.
\]
Exact evaluation at these rational samples gives
\begin{align*}
 \Norm{f(q_1)}=\Norm{f(q_3)}
   &=\frac{39952889}{549755813888}
     \approx 7.27\mathbin{\times}10^{-5},\\
 \Norm{f(q_2)}
   &=\frac{537358761}{35184372088832}
     \approx 1.53\mathbin{\times}10^{-5}.
\end{align*}

The first polynomial track admits the following entirely finite refinement trace.

\begin{center}
\small
\begin{tabularx}{\textwidth}{@{}c>{\raggedright\arraybackslash}Xc@{}}
\toprule
state & winding-one box & $\diamq(B)$\\
\midrule
$0$ & $[69/64,35/32]\times[9/32,19/64]$
    & $1/2048$\\
$1$ & $[555/512,1111/1024]\times[297/1024,149/512]$
    & $1/524288$\\
$2$ & $[71055/65536,4441/4096]\times[19039/65536,595/2048]$
    & $1/2147483648$\\
\bottomrule
\end{tabularx}
\end{center}
The edge-index certificates are respectively $(0,2,0,2)$,
$(0,2,0,2)$, and $(2,0,0,2)$, so every row has winding four.  Their rational
centers are nested approximants with residual norms
\[
 \begin{aligned}
 \frac{39952889}{549755813888}
   &\approx7.27\mathbin{\times}10^{-5},\\
 \frac{8767783105145}{9223372036854775808}
   &\approx9.51\mathbin{\times}10^{-7},\\
 \frac{546115072853076822625}{633825300114114700748351602688}
   &\approx8.62\mathbin{\times}10^{-10}.
 \end{aligned}
\]
These are successive finite refinement states.

\section{Conclusion}

A squarefree linear polynomial has a directly computed exact
Gaussian-rational zero.  In degree $d\ge2$, the polynomial theorem returns
exactly $d$ separated interfaces that answer every finite refinement request.
Each finite observation consists of rational boxes, rational witnesses, and
Sturm/eliminant certificates.  The same
squared-diagonal budget bounds the anchor's Hermitian projective quadrance
and, up to the
explicit factor $L_1(f,A)$, its polynomial residual quadrance.  The global
count is four times the degree, local windings are four, and distinctness has a
coefficient-computable separation radius.  Exact terminal outputs carry a
Gaussian-rational evaluation certificate.

The RFAP addendum preserves the same discipline.  It cancels a displayed
quotient by finite Euclid, attaches signed multiplicities to exact rational
terminals or independently shrinking tracks, separates numerator and
denominator sites using finite B\'ezout margins, and computes the outer
winding as four times the reduced degree difference.  Repeated factors remain
one geometric site with a larger winding; the squarefree specialization has
only $+4$ and $-4$ sites.

Every requested approximation is produced in finite time by the total
finite procedures proved above.  Moreover,
\cref{thm:rfap-one-calculation-size} computes in advance a finite operation
budget and component-width budget for the eliminant, clearing, Sturm,
winding, safe-cut, and requested-refinement calculations that produce it.

\appendix
\section{Formalization and dependency audit}
\label{sec:formalization-audit}

Our use of Lean should not be read as unqualified approval of its language
design.  Its small proof-checking kernel is valuable; its authoring language
is often remarkably unpleasant.  Much of the surface feels accumulated
rather than designed: dependent syntax, macros, tactic dialects, coercions,
typeclass search, and elaboration heuristics overlap in ways that are ad hoc
and difficult to predict.  A mathematically transparent step can produce
source text whose meaning depends on an invisible elaboration path, and a
failed step can produce diagnostics describing the elaborator rather than
the mathematical obstruction.  The result is too often noisy,
indecipherable, and unnecessarily ugly.  Although users can introduce local
notation, they cannot turn the system as a whole into a language-neutral
mathematical vernacular: core commands, library names, and diagnostics remain
overwhelmingly English.  It is difficult to reconcile these choices with the
mathematical tradition that good notation should expose structure with
economy and beauty.  This is not merely an aesthetic complaint.  Obscure
elaboration makes proofs harder to audit and makes unintended dependencies
easier to miss.

Mathlib presented the more consequential difficulty.  Its breadth is an
extraordinary resource, but many convenient general-purpose routes carry
\texttt{Classical.choice} into proofs of statements whose data and content
are entirely finite.  An import by itself does not put choice into the axiom
closure of a theorem.  The problem arises when an apparently elementary
lemma, bundled structure, normalization routine, or synthesized instance
uses a choice-bearing route, which then appears in \texttt{\#print axioms}
for the paper-facing result.  For ordinary classical developments this may
be harmless.  Here it would silently surrender exactly the claim under
investigation.  A machine-checked theorem is not a choice-free theorem merely
because the kernel accepts its proof term; its complete dependency closure
must also be choice-free.

We therefore audited actual proof dependencies and replaced the offending
routes rather than accepting them as library overhead.  The project-owned
\texttt{Purelib.*} modules provide explicit field laws,
constructive fractions and Gaussian rationals, executable dense polynomial
arithmetic and Euclid, finite bivariate and matrix calculations, and direct
separated-Sylvester eliminant certificates.  Elsewhere we replaced generic order,
inverse, normalization, and finite-cardinality arguments by direct sign
proofs and structural recursion on displayed finite data.  Some of these
local arguments prove facts that are elementary on paper; their length in
Lean records how difficult the library made it to use those facts without an
irrelevant selection principle.

The remaining dependence on propositional extensionality was no more
indispensable.  Simplifier-generated proposition equalities, constructor
injectivity certificates, and inherited Core proofs could introduce
\texttt{propext} into otherwise direct finite arguments.  We replaced those
routes by explicit implications, equality elimination, structural induction,
and proof-producing certificates, including tracked repairs to the Core
providers.  Finite arithmetic does not become extensional merely because a
simplifier chose to rewrite a proposition instead of using an implication.
Nor should users have to patch foundational library proofs to avoid axioms
that the mathematical argument never needed.

The repository-wide audit now reports an empty transitive axiom closure for
every selected mathematical declaration: neither \texttt{propext},
\texttt{Classical.choice}, \texttt{Quot.sound}, nor any other axiom appears.
The selection includes private mathematical declarations, registered compiler
rewrites, and active script entrypoints.  Every selected declaration's full
dependency closure is checked, including its generated proof helpers.
Independent compiler and elaborator helpers outside those closures are a
separate audit obligation; clean mathematical certificates do not establish
that their generators are axiom-free.  An independent nanoda check accepts the
complete exported proof dependencies under the same empty-axiom policy.

The earlier fraction implementation still exposed the axiom
\texttt{Quot.sound}.  It did so only after an explicit cross-multiplication
certificate had proved that two displayed representatives were equivalent.
That is not a selection principle, but it still obscured the finite
justification behind ambient quotient equality.  We therefore replaced
quotient fractions by computably canonical reduced records.  Equivalence and
normalization certificates remain first-class evidence, and equality of
canonical forms is proved directly.  The same discipline required structural
representations of the other finite objects and direct equations for several
recursive algorithms.
The former implementation is retired and remains available in Git history;
no active paper declaration imports it.  The result answers the interface
question constructively: whenever canonicalization is available, the kernel
need not be asked to identify quotient classes at all.

There is a broader library-design point.  Classical mathematics belongs in a
general mathematical library, but unnecessary axioms should
not be baked into results that admit finite direct proofs.  When the same
conclusion can be established without those axioms, removing them is a genuine
strengthening of the theorem, not a stylistic variant.  Constructive
primitives should sit below optional classical conveniences, and axiom
dependencies should be visible enough that users can preserve such a
boundary without rebuilding basic algebra.  Replacing a gratuitous
assumption by a proof that does not need it is an improvement that a
foundational library should actively prefer.  We offer the local replacements
in this development as evidence that this stronger organization is
practical.

\newpage
\paragraph{Reproducibility.}
The \texttt{papers.RFAP.RFAP} module exposes the principal executable
declarations in the \texttt{Purelib.Papers.RFAP.\allowbreak Paper} namespace:
\begin{verbatim}
Purelib.Papers.RFAP.Paper.isolateByBoxes
Purelib.Papers.RFAP.Paper.isolateGeneralByBoxes
Purelib.Papers.RFAP.Paper.rationalFunctionArgumentPrinciple
Purelib.Papers.RFAP.Paper.squarefreeRationalFunctionArgumentPrinciple
Purelib.Papers.RFAP.Paper.stageThreeResourceCertificate
\end{verbatim}
From the root of the accompanying source archive, the complete reproducibility
check is
\begin{verbatim}
scripts/run-lake.sh --wfail build
scripts/run-lake.sh --wfail build PurelibExamples Test \
  CCLExamples DivineProportionsExamples \
  RationalRotationsExamples
scripts/check-axioms.sh --configured-scripts --full
scripts/check-nanoda.sh
\end{verbatim}
The wrapper selects the pinned toolchain with its tracked Core proof repairs.
The warning-fatal builds check the reusable and publication sources together
with the retained tests and optional examples.  The full axiom checker also
builds and audits configured script owners; it requires empty complete
dependency closures for its selected declarations.  The final command
independently checks the complete exported proof dependencies from every
active source module while admitting no axiom declarations.  The repository
README documents toolchain preparation and the exact audit-root policy.

\small
\begin{thebibliography}{9}
\raggedright
\bibitem{Ruitenburg}
W. B. G. Ruitenburg,
\emph{Constructing roots of polynomials over the complex numbers},
in \emph{Computational Aspects of Lie Group Representations and Related
Topics}, CWI Tract 84, 1991, 107--128;
\url{https://ir.cwi.nl/pub/13205/13205D.pdf}.

\bibitem{Richman}
F. Richman,
\emph{The fundamental theorem of algebra: a constructive development without
choice},
Pacific Journal of Mathematics 196 (2000), 213--230;
\url{https://doi.org/10.2140/pjm.2000.196.213}.

\bibitem{GeuversWiedijkZwanenburg}
H. Geuvers, F. Wiedijk, and J. Zwanenburg,
\emph{A constructive proof of the Fundamental Theorem of Algebra without
using the rationals},
in \emph{Types for Proofs and Programs: TYPES 2000}, Lecture Notes in
Computer Science 2277, Springer, 2002, 96--111;
\url{https://doi.org/10.1007/3-540-45842-5_7}.

\bibitem{LiPaulson}
W. Li and L. C. Paulson,
\emph{Evaluating winding numbers and counting complex roots through Cauchy indices in Isabelle\slash HOL},
Journal of Automated Reasoning 64 (2020), 331--360;
\url{https://doi.org/10.1007/s10817-019-09521-3}; arXiv:1804.03922.

\bibitem{Eisermann}
M. Eisermann,
\emph{The Fundamental Theorem of Algebra made effective: an elementary real-algebraic proof via Sturm chains},
American Mathematical Monthly 119 (2012), 715--752;
\url{https://doi.org/10.4169/amer.math.monthly.119.09.715};
arXiv:0808.0097.

\bibitem{PerrucciRoyQuantitative}
D. Perrucci and M.-F. Roy,
\emph{Quantitative fundamental theorem of algebra},
Quarterly Journal of Mathematics 70 (2019), 1009--1037;
\url{https://doi.org/10.1093/qmath/haz008}; arXiv:1803.04358.

\bibitem{PerrucciRoyWinding}
D. Perrucci and M.-F. Roy,
\emph{Algebraic winding numbers},
Journal of Algebra 659 (2024), 581--610;
\url{https://doi.org/10.1016/j.jalgebra.2024.07.010}; arXiv:2305.08638.

\bibitem{Wilf}
H. S. Wilf,
\emph{A global bisection algorithm for computing the zeros of polynomials in
the complex plane},
Journal of the ACM 25 (1978), no.~3, 415--420;
\url{https://doi.org/10.1145/322077.322084}.

\bibitem{Wildberger}
N. J. Wildberger,
\emph{One dimensional metrical geometry},
arXiv:math/0701338, 2007;
\url{https://arxiv.org/abs/math/0701338}.
\end{thebibliography}

\end{document}
